\newMonth{Majus}
\setrunningtitles{Majus}{May}

% ---- martyrology/mart05/mart0501.htm
\needspace{10\baselineskip}
\begin{paracol}{2}
\selectlanguage{latin}
\begin{center}{\color{gregoriocolor} Kaléndis Maji. 
 Luna\dots\ }\end{center}
\switchcolumn
\selectlanguage{english}
\begin{center}{\color{gregoriocolor} The 1\textsuperscript{st} Day of May. The\dots\ Day of the Moon.}\end{center}
\end{paracol}

\noindent\begin{tabularx}{\linewidth}{*{19}{>{\centering\arraybackslash}X}}
 \textcolor{gregoriocolor}{a} & \textcolor{gregoriocolor}{b} & \textcolor{gregoriocolor}{c} & \textcolor{gregoriocolor}{d} & \textcolor{gregoriocolor}{e} & \textcolor{gregoriocolor}{f} & \textcolor{gregoriocolor}{g} & \textcolor{gregoriocolor}{h} & \textcolor{gregoriocolor}{i} & \textcolor{gregoriocolor}{k} & \textcolor{gregoriocolor}{l} & \textcolor{gregoriocolor}{m} & \textcolor{gregoriocolor}{n} & \textcolor{gregoriocolor}{p} & \textcolor{gregoriocolor}{q} & \textcolor{gregoriocolor}{r} & \textcolor{gregoriocolor}{s} & \textcolor{gregoriocolor}{t} & \textcolor{gregoriocolor}{u} \\
 4 & 5 & 6 & 7 & 8 & 9 & 10 & 11 & 12 & 13 & 14 & 15 & 16 & 17 & 18 & 19 & 20 & 21 & 22 \\
\end{tabularx}
\vspace{0.5\baselineskip}
\noindent\begin{tabularx}{\linewidth}{*{12}{>{\centering\arraybackslash}X}}
 \textcolor{gregoriocolor}{A} & \textcolor{gregoriocolor}{B} & \textcolor{gregoriocolor}{C} & \textcolor{gregoriocolor}{D} & \textcolor{gregoriocolor}{E} & F & \textcolor{gregoriocolor}{F} & \textcolor{gregoriocolor}{G} & \textcolor{gregoriocolor}{H} & \textcolor{gregoriocolor}{M} & \textcolor{gregoriocolor}{N} & \textcolor{gregoriocolor}{P} \\
 23 & 24 & 25 & 26 & 27 & 28 & 27 & 28 & 29 & 1 & 2 & 3 \\
\end{tabularx}

\begin{paracol}{2}
\selectlanguage{latin}
\lettrine[lines=2]{N}{atális} beatórum Philíppi et Jacóbi Apostolórum. Ex 
 his Philíppus, cum omnem fere Scythiam ad Christi fidem convertísset, tandem 
 apud Hierápolim, Asiæ civitátem, cruci affíxus et lapídibus 
 óbrutus, glorióso fine quiévit; Jacóbus vero, qui et frater Dómini légitur et primus 
 Hierosolymórum Epíscopus, e pinna Templi præcipitátus, confráctis inde 
 crúribus, ac fullónis fuste in cérebro percússus, intériit, ibique, non 
 longe a Templo, sepúltus est.
\switchcolumn
\selectlanguage{english}
\lettrine[lines=2]{T}{he} birthday of the blessed apostles Philip and James. 
 Philip, after having converted nearly all of Scythia to the faith of Christ, 
 went to Hieropolis, a city in Asia, where he was fastened to a cross and 
 stoned, and thus ended his life gloriously. James, who is also called 
 the brother of our Lord, was the first bishop of Jerusalem. Being 
 hurled down from a pinnacle of the temple, his legs were broken, and being 
 struck on the head with a dyer's staff, he expired and was buried near the 
 temple.
\switchcolumn*
\selectlanguage{latin}
Romæ item natális sancti Pii Quinti, ex Ordine 
 Prædicatórum, Papæ et Confessóris; qui, ecclesiásticæ disciplínæ restituéndæ, 
 hærésibus exstirpándis et Christiáni nóminis hóstibus conteréndis strénue ac 
 felíciter incúmbens, vitæ ac legum sanctitáte cathólicam Ecclésiam 
 gubernávit. Ipsíus autem festívitas tértio Nonas mensis hujus 
 celebrátur.
\switchcolumn
\selectlanguage{english}
At Rome, Pope St. Pius V of the Order of Preachers, who 
 laboured zealously and successfully for the re-establishment of church 
 discipline, the stamping out of heresies, and the destruction of the enemies 
 of the Christian name. He governed the Catholic Church by holy laws, 
 and the example of a saintly life. His feast is observed on the fifth 
 day of May.
\switchcolumn*
\selectlanguage{latin}
In Ægypto sancti Jeremíæ Prophétæ, qui, a pópulo 
 lapídibus óbrutus, apud Taphnas occúbuit, ibíque sepúltus est; ad cujus 
 sepúlcrum fidéles (ut refert sanctus Epiphánius) supplicáre consuevérunt, 
 índeque sumpto púlvere, áspidum mórsibus medéntur.
\switchcolumn
\selectlanguage{english}
In Egypt, St. Jeremias, prophet, who was stoned to death 
 by the people at Taphnas, where he was buried. St. Epiphanius tells 
 that the faithful were accustomed to pray at his grave, and to take away 
 from it dust to heal those who were stung by serpents.
\switchcolumn*
\selectlanguage{latin}
In território Vivariénsi, in Gálliis, beáti Andéoli, 
 Subdiáconi, qui, una cum áliis, a beáto Polycárpo, Smyrnénsi Epíscopo, ex 
 Oriénte in Gálliam ad prædicándum verbum Dei missus est. Hic, sub 
 Sevéro Imperatóre, spinósis fústibus cæsus, demum, per ensem lígneum cápite 
 in quátuor partes in modum crucis conscísso, martyrium consummávit.
\switchcolumn
\selectlanguage{english}
In France, in the Province of Vivarias, blessed Andeol, 
 subdeacon, who was sent from the East into Gaul with others by St. Polycarp 
 to preach the word of God. Under Emperor Severus he was scourged with 
 thorny sticks, and having his head split with a wooden sword into four 
 parts, in the shape of a cross, he completed his martyrdom.
\switchcolumn*
\selectlanguage{latin}
Oscæ, in Hispánia, sanctórum Mártyrum Oréntii et 
 Patiéntiæ.
\switchcolumn
\selectlanguage{english}
At Huesca in Spain, the holy martyrs Orentius and 
 Patience.
\switchcolumn*
\selectlanguage{latin}
Apud Colúmnam vicum, in Aurelianénsi Gálliæ território, 
 pássio sancti Sigismúndi, Regis Burgundiónum, qui in púteum demérsus 
 occúbuit, ac póstea miráculis cláruit. Sacrum vero ipsíus corpus, e 
 púteo tandem extráctum, ad Ecclésiam Agaunénsis monastérii, intra Sedunénsis 
 diœcésis términos siti, delátum est ibíque honorífice collocátum.
\switchcolumn
\selectlanguage{english}
In the town of Columna, in the province of Orleans in 
 France, the martyrdom of St. Sigismund, king of Burgundy. He met death 
 by being drowned in a well, and was afterwards famous for his miracles. 
 His venerable body was later recovered and taken to the monastery of Agaune 
 in the diocese of Sitten where it was honorably entombed.
\switchcolumn*
\selectlanguage{latin}
Antisiodóri sancti Amatóris, Epíscopi et Confessóris.
\switchcolumn
\selectlanguage{english}
At Auxerre, St. Amator, bishop and confessor.
\switchcolumn*
\selectlanguage{latin}
Auscii, in Gállia, sancti Oriéntii Epíscopi.
\switchcolumn
\selectlanguage{english}
At Auch in France, Bishop St. Orientius.
\switchcolumn*
\selectlanguage{latin}
Elviæ, in Anglia, sancti Asaphi Epíscopi, cujus nómine 
 ipsa cívitas Episcopális póstmodum est insigníta.
\switchcolumn
\selectlanguage{english}
At Llanelwy in Wales, Bishop St. Asaph, in whose memory 
 the cathedral city was later named.
\switchcolumn*
\selectlanguage{latin}
Foro Lívii sancti Peregríni, ex Ordine Servórum beátæ 
 Maríæ Vírginis.
\switchcolumn
\selectlanguage{english}
At Forli, St. Peregrinus of the Order of Servites of the 
 Blessed Virgin Mary.
\switchcolumn*
\selectlanguage{latin}
Bérgomi sanctæ Gratæ 
 Víduæ.
\switchcolumn
\selectlanguage{english}
At Bergamo, St. Grata, widow.
\switchcolumn*
\selectlanguage{latin}
\end{paracol}

% \continueMonth{Majus}

% ---- martyrology/mart05/mart0502.htm
\needspace{10\baselineskip}
\begin{paracol}{2}
\selectlanguage{latin}
\begin{center}{\color{gregoriocolor} Sexto Nonas Maji. 
 Luna\dots\ }\end{center}
\switchcolumn
\selectlanguage{english}
\begin{center}{\color{gregoriocolor} The 
 2\textsuperscript{nd} Day of May. The\dots\ Day of the Moon.}\end{center}
\end{paracol}

\noindent\begin{tabularx}{\linewidth}{*{19}{>{\centering\arraybackslash}X}}
 \textcolor{gregoriocolor}{a} & \textcolor{gregoriocolor}{b} & \textcolor{gregoriocolor}{c} & \textcolor{gregoriocolor}{d} & \textcolor{gregoriocolor}{e} & \textcolor{gregoriocolor}{f} & \textcolor{gregoriocolor}{g} & \textcolor{gregoriocolor}{h} & \textcolor{gregoriocolor}{i} & \textcolor{gregoriocolor}{k} & \textcolor{gregoriocolor}{l} & \textcolor{gregoriocolor}{m} & \textcolor{gregoriocolor}{n} & \textcolor{gregoriocolor}{p} & \textcolor{gregoriocolor}{q} & \textcolor{gregoriocolor}{r} & \textcolor{gregoriocolor}{s} & \textcolor{gregoriocolor}{t} & \textcolor{gregoriocolor}{u} \\
 5 & 6 & 7 & 8 & 9 & 10 & 11 & 12 & 13 & 14 & 15 & 16 & 17 & 18 & 19 & 20 & 21 & 22 & 23 \\
\end{tabularx}
\vspace{0.5\baselineskip}
\noindent\begin{tabularx}{\linewidth}{*{12}{>{\centering\arraybackslash}X}}
 \textcolor{gregoriocolor}{A} & \textcolor{gregoriocolor}{B} & \textcolor{gregoriocolor}{C} & \textcolor{gregoriocolor}{D} & \textcolor{gregoriocolor}{E} & F & \textcolor{gregoriocolor}{F} & \textcolor{gregoriocolor}{G} & \textcolor{gregoriocolor}{H} & \textcolor{gregoriocolor}{M} & \textcolor{gregoriocolor}{N} & \textcolor{gregoriocolor}{P} \\
 24 & 25 & 26 & 27 & 28 & 29 & 28 & 29 & 1 & 2 & 3 & 4 \\
\end{tabularx}

\begin{paracol}{2}
\selectlanguage{latin}
\lettrine[lines=2]{A}{lexandríæ} natális 
 sancti Athanásii, ejúsdem urbis Epíscopi, Confessóris et Ecclésiæ Doctóris, 
 sanctitáte et doctrína claríssimi; in cujus persecutiónem univérsus fere 
 Orbis conjuráverat. Ipse tamen cathólicam fidem, a témpore Constantíni 
 usque ad Valéntem, advérsus Imperatóres ac Prǽsides et innúmeros Epíscopos 
 Ariános strénue propugnávit; a quibus plúrimas perpéssus insídias, prófugus 
 toto Orbe actus est, nec ullus ei tutus ad laténdum supérerat locus. 
 Tandem, ad suam Ecclésiam revérsus, illic, post multos agones multásque 
 patiéntiæ corónas, quadragésimo sexto sui sacerdótii anno migrávit ad 
 Dóminum, témpore Valentiniáni et Valéntis Imperatórum.
\switchcolumn
\selectlanguage{english}
\lettrine[lines=2]{A}{t} Alexandria, the birthday of St. 
 Athanasius, bishop of that city, confessor and doctor of the Church, most 
 celebrated for sanctity and learning. Although almost all of the world 
 had formed a conspiracy to persecute him, he courageously defended the 
 Catholic faith, from the reign of Constantine to that of Valens, against 
 emperors, governors, and a multitude of Arian bishops, whose underhanded 
 attacks forced him to wander as an exile over the whole earth without 
 finding a place of security. At length, however, he was restored to 
 his church, and after overcoming many trials, and winning many crowns by his 
 patience, he departed for heaven in the forty-sixth year of his priesthood, 
 in the time of the emperors Valentinian and Valens.
\switchcolumn*
\selectlanguage{latin}
Floréntiæ item natális 
 sancti Antoníni, ex Ordine Prædicatórum, Epíscopi et Confessóris, doctrína 
 et sanctitáte célebris. Ipsíus autem festívitas sexto Idus mensis 
 hujus recólitur.
\switchcolumn
\selectlanguage{english}
At Florence, Bishop St. Antoninus 
 of the Order of Preachers, renowned for sanctity and learning. His 
 feast is kept on the 10th of this month.
\switchcolumn*
\selectlanguage{latin}
Romæ sanctórum 
 Mártyrum Saturníni, Neópoli, Germáni et Cælestíni, qui, multa passi, in cárcerem demum conjécti, ibi in Dómino quievérunt.
\switchcolumn
\selectlanguage{english}
At Rome, the holy martyrs 
 Saturninus, Neopolus, Germanus, and Celestine, who after much suffering were 
 thrown into prison, where they found rest in the Lord.
\switchcolumn*
\selectlanguage{latin}
Eódem die sancti 
 Vindemiális, Epíscopi et Mártyris, qui, una cum sanctis Epíscopis Eugénio et 
 Longíno doctrína et miráculis advérsus Ariános decértans, a Rege Wandalórum 
 Hunneríco jubétur váriis torméntis afflígi ad tandem cápite obtruncári.
\switchcolumn
\selectlanguage{english}
The same day, St. Vindemial, bishop 
 and martyr, who with the holy bishops Eugene and Longinus, combated the 
 Arians by his teaching and miracles, and was beheaded by order of Hunneric, 
 king of the Vandals.
\switchcolumn*
\selectlanguage{latin}
Híspali, in Hispánia, sancti Felícis, Diáconi et Mártyris.
\switchcolumn
\selectlanguage{english}
At Seville in Spain, St. Felix, 
 deacon and martyr.
\switchcolumn*
\selectlanguage{latin}
Attalíæ, in Pamphylia, 
 sanctórum Mártyrum Exsupérii, et Zoes uxóris, atque Cyríaci et Theodúli 
 filiórum; qui, sub Hadriáno Imperatóre, cum servi essent cujúsdam viri 
 Pagáni, omnes, ipso hero jubénte, ob líberam Christiánæ fídei professiónem, 
 primum verberáti sunt ac veheménter torti, deínde, in accénsum clíbanum 
 injécti, ánimas suas Deo tradidérunt.
\switchcolumn
\selectlanguage{english}
At Attalia in Pamphylia, the holy 
 martyrs Exuperius and Zoe, his wife, with their sons, Cyriacus and Theodulus. 
 They were the slaves of a man named Paganus. During the reign of 
 Emperor Hadrian, because of their outspoken profession of the Christian 
 faith, their master ordered them to be scourged and severely tortured. 
 They were finally cast into an oven, and in this way gave up their souls to 
 God.
\switchcolumn*
\selectlanguage{latin}
\end{paracol}


% ---- martyrology/mart05/mart0503.htm
\needspace{10\baselineskip}
\begin{paracol}{2}
\selectlanguage{latin}
\begin{center}{\color{gregoriocolor} Quinto Nonas Maji. 
 Luna\dots\ }\end{center}
\switchcolumn
\selectlanguage{english}
\begin{center}{\color{gregoriocolor} The 
 3\textsuperscript{rd} Day of May. The\dots\ Day of the Moon.}\end{center}
\end{paracol}

\noindent\begin{tabularx}{\linewidth}{*{19}{>{\centering\arraybackslash}X}}
 \textcolor{gregoriocolor}{a} & \textcolor{gregoriocolor}{b} & \textcolor{gregoriocolor}{c} & \textcolor{gregoriocolor}{d} & \textcolor{gregoriocolor}{e} & \textcolor{gregoriocolor}{f} & \textcolor{gregoriocolor}{g} & \textcolor{gregoriocolor}{h} & \textcolor{gregoriocolor}{i} & \textcolor{gregoriocolor}{k} & \textcolor{gregoriocolor}{l} & \textcolor{gregoriocolor}{m} & \textcolor{gregoriocolor}{n} & \textcolor{gregoriocolor}{p} & \textcolor{gregoriocolor}{q} & \textcolor{gregoriocolor}{r} & \textcolor{gregoriocolor}{s} & \textcolor{gregoriocolor}{t} & \textcolor{gregoriocolor}{u} \\
 6 & 7 & 8 & 9 & 10 & 11 & 12 & 13 & 14 & 15 & 16 & 17 & 18 & 19 & 20 & 21 & 22 & 23 & 24 \\
\end{tabularx}
\vspace{0.5\baselineskip}
\noindent\begin{tabularx}{\linewidth}{*{12}{>{\centering\arraybackslash}X}}
 \textcolor{gregoriocolor}{A} & \textcolor{gregoriocolor}{B} & \textcolor{gregoriocolor}{C} & \textcolor{gregoriocolor}{D} & \textcolor{gregoriocolor}{E} & F & \textcolor{gregoriocolor}{F} & \textcolor{gregoriocolor}{G} & \textcolor{gregoriocolor}{H} & \textcolor{gregoriocolor}{M} & \textcolor{gregoriocolor}{N} & \textcolor{gregoriocolor}{P} \\
 25 & 26 & 27 & 28 & 29 & 30 & 29 & 1 & 2 & 3 & 4 & 5 \\
\end{tabularx}

\begin{paracol}{2}
\selectlanguage{latin}
\lettrine[lines=2]{H}{ierosólymis} Invéntio 
 Sacrosánctæ Crucis Domínicæ, sub Constantíno Imperatóre.
\switchcolumn
\selectlanguage{english}
\lettrine[lines=2]{A}{t} Jerusalem, in the time of 
 Emperor Constantine, the finding of the holy Cross of our Lord.
\switchcolumn*
\selectlanguage{latin}
Romæ, via Nomentána, 
 pássio sanctórum Mártyrum Alexándri Papæ Primi, Evéntii et Theodúli 
 Presbyterórum. Ex his Alexánder, sub Hadriáno Príncipe et Aureliáno 
 Júdice, post víncula, cárceres, equúleum, úngulas et ignes, punctis 
 crebérrimis per tota membra confóssus ac perémptus est; Evéntius vero et 
 Theodúlus, post longos cárceres, ígnibus examináti, ad 
 últimum decolláti 
 sunt.
\switchcolumn
\selectlanguage{english}
At Rome, on the Via Nomentana, the 
 holy martyrs Pope Alexander and the priests Eventius and Theodulus. 
 Alexander was bound, imprisoned, racked, lacerated with hooks, burned, and 
 had all his limbs pierced with pointed instruments, and finally met death, 
 under Emperor Hadrian and the judge Aurelian. Eventius and Theodulus 
 after a long imprisonment were exposed to flames and then beheaded.
\switchcolumn*
\selectlanguage{latin}
Nárniæ sancti Juvenáli, 
 Epíscopi et Confessóris.
\switchcolumn
\selectlanguage{english}
At Narni, St. Juvenal, bishop and 
 confessor.
\switchcolumn*
\selectlanguage{latin}
Apud montem Senárium, 
 in Etrúria, natális sanctórum Sostenǽi et Ugucciónis Confessórum, e septem 
 Fundatóribus Ordinis Servórum beátæ Maríæ Vírginis; qui, cælitus admóniti, 
 eádem die et hora, salutatiónem Angélicam recitántes, e vita migrárunt. 
 Ipsórum autem ac Sociórum festum prídie Idus Februárii celebrátur.
\switchcolumn
\selectlanguage{english}
On Mount Senario in Etruria, Saints 
 Sosteneo and Ugoccio, confessors, of the seven founders of the Order of 
 Servites of the Blessed Virgin Mary. Responding to a voice from 
 heaven, they departed this life on the same day and at the same hour, while 
 reciting the angelical salutation. Their feast is observed with the 
 rest of their companions on the 12th day of February.
\switchcolumn*
\selectlanguage{latin}
Constantinópoli 
 sanctórum Mártyrum Alexándri mílitis, et Antonínæ Vírginis. Hæc, in 
 persecutióne Maximiáni, sub Prǽside Festo, ad lupánar damnáta, et ab 
 Alexándro, qui pro ipsa ibi remánserat, mutátis véstibus clam edúcta, cum eo 
 póstmodum jussa est torquéri; et ambo simul in ignem, præcísis mánibus, pro 
 Christo sunt injécti, atque ita, egrégio perácto certámine, coronántur.
\switchcolumn
\selectlanguage{english}
At Constantinople, the holy martyrs 
 Alexander, soldier, and Antonina, virgin. In the persecution of 
 Maximian, under the governor Festus, Antonina, having been condemned to 
 remain in a place of debauchery, was delivered by Alexander, who secretly 
 exchanged garments with her, and took her place. They were tortured 
 together, both had their hands cut off, were cast into the fire, and 
 received their crowns at the end of their heroic combat for the faith.
\switchcolumn*
\selectlanguage{latin}
In Thebáide sanctórum 
 Mártyrum Timóthei et Mauræ cónjugis, quos Ariánus Præféctus, post multa 
 torménta, cruci jussit affígi; in qua, per novem dies cum vivi pependíssent 
 ac se ipsos in fide roborássent, martyrium consummárunt.
\switchcolumn
\selectlanguage{english}
In Thebais, the holy martyrs 
 Timothy and his wife Maura. The Arian prefect caused them to be 
 tortured, and then fastened to a cross, on which they remained alive for 
 nine days, encouraging each other to persevere in the faith, until they 
 completed their martyrdom.
\switchcolumn*
\selectlanguage{latin}
Aphrodísiæ, in Cária, 
 sanctórum Mártyrum Diodóri et Rhodopiáni, qui, in persecutióne Diocletiáni 
 Imperatóris, a cívibus suis lapidáti sunt.
\switchcolumn
\selectlanguage{english}
At Aphrodisia in Caria, the holy 
 martyrs Diodorus and Rodopian, who were stoned to death by their fellow 
 citizens, in the persecution of Diocletian.
\switchcolumn*
\selectlanguage{latin}
\end{paracol}


% ---- martyrology/mart05/mart0504.htm
\needspace{10\baselineskip}
\begin{paracol}{2}
\selectlanguage{latin}
\begin{center}{\color{gregoriocolor} Quarto Nonas Maji. 
 Luna\dots\ }\end{center}
\switchcolumn
\selectlanguage{english}
\begin{center}{\color{gregoriocolor} The 
 4\textsuperscript{th} Day of May. The\dots\ Day of the Moon.}\end{center}
\end{paracol}

\noindent\begin{tabularx}{\linewidth}{*{19}{>{\centering\arraybackslash}X}}
 \textcolor{gregoriocolor}{a} & \textcolor{gregoriocolor}{b} & \textcolor{gregoriocolor}{c} & \textcolor{gregoriocolor}{d} & \textcolor{gregoriocolor}{e} & \textcolor{gregoriocolor}{f} & \textcolor{gregoriocolor}{g} & \textcolor{gregoriocolor}{h} & \textcolor{gregoriocolor}{i} & \textcolor{gregoriocolor}{k} & \textcolor{gregoriocolor}{l} & \textcolor{gregoriocolor}{m} & \textcolor{gregoriocolor}{n} & \textcolor{gregoriocolor}{p} & \textcolor{gregoriocolor}{q} & \textcolor{gregoriocolor}{r} & \textcolor{gregoriocolor}{s} & \textcolor{gregoriocolor}{t} & \textcolor{gregoriocolor}{u} \\
 7 & 8 & 9 & 10 & 11 & 12 & 13 & 14 & 15 & 16 & 17 & 18 & 19 & 20 & 21 & 22 & 23 & 24 & 25 \\
\end{tabularx}
\vspace{0.5\baselineskip}
\noindent\begin{tabularx}{\linewidth}{*{12}{>{\centering\arraybackslash}X}}
 \textcolor{gregoriocolor}{A} & \textcolor{gregoriocolor}{B} & \textcolor{gregoriocolor}{C} & \textcolor{gregoriocolor}{D} & \textcolor{gregoriocolor}{E} & F & \textcolor{gregoriocolor}{F} & \textcolor{gregoriocolor}{G} & \textcolor{gregoriocolor}{H} & \textcolor{gregoriocolor}{M} & \textcolor{gregoriocolor}{N} & \textcolor{gregoriocolor}{P} \\
 26 & 27 & 28 & 29 & 30 & 1 & 1 & 2 & 3 & 4 & 5 & 6 \\
\end{tabularx}

\begin{paracol}{2}
\selectlanguage{latin}
\lettrine[lines=2]{A}{pud} Ostia Tiberína 
 sanctæ Mónicæ, beáti Augustíni matris, cujus ille præcláram vitam, in libro 
 nono Confessiónum, testátam relíquit.
\switchcolumn
\selectlanguage{english}
\lettrine[lines=2]{A}{t} Ostia, the birthday of St. 
 Monica, mother of blessed Augustine. He has left us in the ninth book 
 of his Confessions a beautiful sketch of her life.
\switchcolumn*
\selectlanguage{latin}
In metállo Phennénsi 
 Palæstínæ natális beáti Silváni, Gazæ Epíscopi, qui, in Diocletiáni 
 Imperatóris persecutióne, Galérii Maximiáni Cǽsaris mandáto, cum plúrimis 
 suis Cléricis, martyrio coronátus est.
\switchcolumn
\selectlanguage{english}
At the metal mines of Phennes in 
 Palestine, the birthday of blessed Silvanus, bishop of Gaza, who was crowned 
 with martyrdom with many of his clerics by the command of Caesar Galerius 
 Maximian, in the persecution of Diocletian.
\switchcolumn*
\selectlanguage{latin}
Hierosólymis sancti 
 Cyríaci Epíscopi, qui, cum loca sancta visitáret, ibídem, sub Juliáno 
 Apóstata, cæsus est.
\switchcolumn
\selectlanguage{english}
At Jerusalem, in the reign of 
 Julian the Apostate, St. Cyriacus, bishop, who was murdered while visiting 
 the holy places.
\switchcolumn*
\selectlanguage{latin}
Cameríni sancti 
 Porphyrii Presbyteri et Mártyris, qui, sub Décio Imperatóre et Antíocho 
 Prǽside, cum plúrimos (in quibus fuit Venántius) ad Christi fidem 
 convertísset, cápite amputátus est.
\switchcolumn
\selectlanguage{english}
At Camerinum, St. Porphyry, priest 
 and martyr. Because he converted many to the faith (among them 
 Venantius), he was beheaded during the reign of Emperor Decius and the 
 governor Antiochus.
\switchcolumn*
\selectlanguage{latin}
In metállo Phennénsi 
 Palæstínæ sanctórum trigínta novem Mártyrum, qui, ad metálla damnáti, demum, 
 post adustiónem candéntis ferri et ália torménta, simul cápite cæsi sunt.
\switchcolumn
\selectlanguage{english}
Also in the mines of Phennes, 
 thirty-nine holy martyrs, who were condemned to work there, to be branded 
 with hot irons, to undergo other torments, and finally all to be beheaded at 
 the same time.
\switchcolumn*
\selectlanguage{latin}
Lauréaci, in Nórico 
 Ripénsi, sancti Floriáni Mártyris, qui, sub Diocletiáno Imperatóre, Aquilíni 
 Prǽsidis jussu, in flumen Anísum, ligáto ad collum saxo, præcipitátus est.
\switchcolumn
\selectlanguage{english}
At Lorch in Austria, under Emperor 
 Diocletian and the governor Aquilinus, the martyr St. Florian, who was 
 thrown into the River Enns, with a stone tied about his neck.
\switchcolumn*
\selectlanguage{latin}
Colóniæ Agrippínæ 
 sancti Paulíni Mártyris.
\switchcolumn
\selectlanguage{english}
At Cologne, the martyr St. Paulinus.
\switchcolumn*
\selectlanguage{latin}
Tarsi, in Cilícia, 
 sanctæ Pelágiæ, Vírginis et Mártyris, quæ, sub Diocletiáno Imperatóre, in 
 bovem æneum candéntem inclúsa, martyrium complévit.
\switchcolumn
\selectlanguage{english}
At Tarsus, St. Pelagia, virgin, who 
 endured martyrdom under Diocletian by being shut up inside an ox made of 
 brass that had been heated to redness.
\switchcolumn*
\selectlanguage{latin}
Nicomedíæ natális 
 sanctæ Antóniæ Mártyris, quæ, nímium torta et váriis afflícta cruciátibus, 
 áltero bráchio tribus suspénsa diébus, et in cárcere biénnio deténta, ad 
 últimum a Priscilliáno Prǽside, in confessióne Dómini, flammis exústa est.
\switchcolumn
\selectlanguage{english}
At Nicomedia, the birthday of St. 
 Antonia, martyr, who was cruelly tortured, subjected to various torments, 
 suspended by one arm for three days, kept two years in prison, and finally 
 delivered to the flames for the confession of Christ by the governor 
 Priscillian.
\switchcolumn*
\selectlanguage{latin}
Medioláni sancti 
 Venérii Epíscopi, cujus virtútes sanctus Joánnes Chrysóstomus, in epístola 
 ad eum scripta, testátas relíquit.
\switchcolumn
\selectlanguage{english}
At Milan, St. Venerius, a bishop 
 whose virtues are attested to by St. John Chrysostom in the epistle which he 
 had written to him.
\switchcolumn*
\selectlanguage{latin}
In território 
 Petragoricénsi sancti Sacerdótis, Epíscopi Lemovicénsis.
\switchcolumn
\selectlanguage{english}
In the province of Perigord, St. 
 Sacerdos, bishop of Limoges.
\switchcolumn*
\selectlanguage{latin}
Hildeshémii, in 
 Saxónia, sancti Godehárdi, Epíscopi et Confessóris, qui ab Innocéntio Papa 
 Secúndo in Sanctórum censum relátus est.
\switchcolumn
\selectlanguage{english}
At Hildesheim in Saxony, St. Gothard, 
 bishop and confessor, who was ranked among the saints by Innocent II.
\switchcolumn*
\selectlanguage{latin}
Antisiodóri sancti 
 Curcódomi Diáconi.
\switchcolumn
\selectlanguage{english}
At Auxerre, St. Curcodomus, deacon.
\switchcolumn*
\selectlanguage{latin}
\end{paracol}


% ---- martyrology/mart05/mart0505.htm
\needspace{10\baselineskip}
\begin{paracol}{2}
\selectlanguage{latin}
\begin{center}{\color{gregoriocolor} Tértio Nonas Maji. 
 Luna\dots\ }\end{center}
\switchcolumn
\selectlanguage{english}
\begin{center}{\color{gregoriocolor} The 
 5\textsuperscript{th} Day of May. The\dots\ Day of the Moon.}\end{center}
\end{paracol}

\noindent\begin{tabularx}{\linewidth}{*{19}{>{\centering\arraybackslash}X}}
 \textcolor{gregoriocolor}{a} & \textcolor{gregoriocolor}{b} & \textcolor{gregoriocolor}{c} & \textcolor{gregoriocolor}{d} & \textcolor{gregoriocolor}{e} & \textcolor{gregoriocolor}{f} & \textcolor{gregoriocolor}{g} & \textcolor{gregoriocolor}{h} & \textcolor{gregoriocolor}{i} & \textcolor{gregoriocolor}{k} & \textcolor{gregoriocolor}{l} & \textcolor{gregoriocolor}{m} & \textcolor{gregoriocolor}{n} & \textcolor{gregoriocolor}{p} & \textcolor{gregoriocolor}{q} & \textcolor{gregoriocolor}{r} & \textcolor{gregoriocolor}{s} & \textcolor{gregoriocolor}{t} & \textcolor{gregoriocolor}{u} \\
 8 & 9 & 10 & 11 & 12 & 13 & 14 & 15 & 16 & 17 & 18 & 19 & 20 & 21 & 22 & 23 & 24 & 25 & 26 \\
\end{tabularx}
\vspace{0.5\baselineskip}
\noindent\begin{tabularx}{\linewidth}{*{12}{>{\centering\arraybackslash}X}}
 \textcolor{gregoriocolor}{A} & \textcolor{gregoriocolor}{B} & \textcolor{gregoriocolor}{C} & \textcolor{gregoriocolor}{D} & \textcolor{gregoriocolor}{E} & F & \textcolor{gregoriocolor}{F} & \textcolor{gregoriocolor}{G} & \textcolor{gregoriocolor}{H} & \textcolor{gregoriocolor}{M} & \textcolor{gregoriocolor}{N} & \textcolor{gregoriocolor}{P} \\
 27 & 28 & 29 & 30 & 1 & 2 & 2 & 3 & 4 & 5 & 6 & 7 \\
\end{tabularx}

\begin{paracol}{2}
\selectlanguage{latin}
\lettrine[lines=2]{S}{ancti} Pii Quinti, ex 
 Ordine Prædicatórum, Papæ et Confessóris, qui Kaléndis mensis hujus 
 obdormívit in Dómino.
\switchcolumn
\selectlanguage{english}
\lettrine[lines=2]{P}{ope} St. Pius V, confessor of the 
 Order of Preachers, who went to sleep in the Lord on the 1st of May.
\switchcolumn*
\selectlanguage{latin}
Romæ sancti Silváni 
 Mártyris.
\switchcolumn
\selectlanguage{english}
At Rome, the martyr St. Silvanus.
\switchcolumn*
\selectlanguage{latin}
Item Romæ sanctæ 
 Crescentiánæ Mártyris.
\switchcolumn
\selectlanguage{english}
Also at Rome, St. Crescentia, 
 martyr.
\switchcolumn*
\selectlanguage{latin}
Leocátæ, in Sicília, 
 sancti Angeli, ex Ordine Carmelitárum, Presbyteri et Mártyris, qui ab 
 hæréticis, ob defensiónem cathólicæ fídei, trucidátus est.
\switchcolumn
\selectlanguage{english}
At Leocata in Sicily, St. Angelus, 
 priest of the Order of Carmelites, who was murdered by the heretics because 
 of his defence of the Catholic faith.
\switchcolumn*
\selectlanguage{latin}
Alexandríæ sancti 
 Euthymii Diáconi, qui ob Christum quiévit in cárcere.
\switchcolumn
\selectlanguage{english}
At Alexandria, St. Euthymius, 
 deacon, who died in prison for the sake of Christ.
\switchcolumn*
\selectlanguage{latin}
Antisiodóri pássio 
 sancti Joviniáni Lectóris.
\switchcolumn
\selectlanguage{english}
At Auxerre, the martyrdom of St. 
 Jovinian, lector.
\switchcolumn*
\selectlanguage{latin}
Thessalonícæ natális 
 sanctórum Mártyrum Irenǽi, Peregríni et Irénes, qui, 
 ígnibus combústi, 
 palmas martyrii percepérunt.
\switchcolumn
\selectlanguage{english}
At Thessalonica, the birthday of 
 the holy martyrs Irenæus, Peregrinus, and Irene, who were burned alive.
\switchcolumn*
\selectlanguage{latin}
Hierosólymis sancti 
 Máximi Epíscopi, qui a Maximiáno Galério Cǽsare, post óculum effóssum 
 pedémque igníto ferro adústum, ad metálla damnátus est; atque, liber inde 
 abíre permíssus et Ecclésiæ Hierosolymitánæ præpósitus, ibi, confessiónis 
 glória præclárus, in pace tandem quiévit.
\switchcolumn
\selectlanguage{english}
At Jerusalem, St. Maximus, bishop, 
 whom Maximian Galerius Caesar condemned to the mines, after having plucked 
 out one of his eyes and branded him on the foot with a hot iron. He 
 was afterwards freed, and allowed to rule the church at Jerusalem, where he 
 died in peace, renowned for the glory of his confession.
\switchcolumn*
\selectlanguage{latin}
Edéssæ, in Syria, 
 sancti Eulógii, Epíscopi et Confessóris.
\switchcolumn
\selectlanguage{english}
At Edessa in Syria, St. Eulogius, 
 bishop and confessor.
\switchcolumn*
\selectlanguage{latin}
Areláte, in Gállia, 
 sancti Hilárii Epíscopi, doctrína et sanctitáte conspícui.
\switchcolumn
\selectlanguage{english}
At Arles in France, the bishop St. 
 Hilary, noted for his learning and sanctity.
\switchcolumn*
\selectlanguage{latin}
Viénnæ, in Gállia, 
 sancti Nicéti Epíscopi, venerábilis sanctitátis viri.
\switchcolumn
\selectlanguage{english}
At Vienne in France, the bishop St. 
 Nicetus, a man venerable for his piety.
\switchcolumn*
\selectlanguage{latin}
Bonóniæ sancti 
 Theodóri Epíscopi, méritis clari.
\switchcolumn
\selectlanguage{english}
At Bologna, St. Theodore, a bishop 
 who was eminent for merits.
\switchcolumn*
\selectlanguage{latin}
Medioláni sancti 
 Gerúntii Epíscopi.
\switchcolumn
\selectlanguage{english}
At Milan, the bishop St. Geruntius.
\switchcolumn*
\selectlanguage{latin}
Eódem die sancti 
 Sacerdótis, Epíscopi Saguntíni.
\switchcolumn
\selectlanguage{english}
On the same day, St. Sacerdos, 
 bishop of Saguntum.
\switchcolumn*
\selectlanguage{latin}
\end{paracol}


% ---- martyrology/mart05/mart0506.htm
\needspace{10\baselineskip}
\begin{paracol}{2}
\selectlanguage{latin}
\begin{center}{\color{gregoriocolor} Prídie Nonas Maji. 
 Luna\dots\ }\end{center}
\switchcolumn
\selectlanguage{english}
\begin{center}{\color{gregoriocolor} The 
 6\textsuperscript{th} Day of May. The\dots\ Day of the Moon.}\end{center}
\end{paracol}

\noindent\begin{tabularx}{\linewidth}{*{19}{>{\centering\arraybackslash}X}}
 \textcolor{gregoriocolor}{a} & \textcolor{gregoriocolor}{b} & \textcolor{gregoriocolor}{c} & \textcolor{gregoriocolor}{d} & \textcolor{gregoriocolor}{e} & \textcolor{gregoriocolor}{f} & \textcolor{gregoriocolor}{g} & \textcolor{gregoriocolor}{h} & \textcolor{gregoriocolor}{i} & \textcolor{gregoriocolor}{k} & \textcolor{gregoriocolor}{l} & \textcolor{gregoriocolor}{m} & \textcolor{gregoriocolor}{n} & \textcolor{gregoriocolor}{p} & \textcolor{gregoriocolor}{q} & \textcolor{gregoriocolor}{r} & \textcolor{gregoriocolor}{s} & \textcolor{gregoriocolor}{t} & \textcolor{gregoriocolor}{u} \\
 9 & 10 & 11 & 12 & 13 & 14 & 15 & 16 & 17 & 18 & 19 & 20 & 21 & 22 & 23 & 24 & 25 & 26 & 27 \\
\end{tabularx}
\vspace{0.5\baselineskip}
\noindent\begin{tabularx}{\linewidth}{*{12}{>{\centering\arraybackslash}X}}
 \textcolor{gregoriocolor}{A} & \textcolor{gregoriocolor}{B} & \textcolor{gregoriocolor}{C} & \textcolor{gregoriocolor}{D} & \textcolor{gregoriocolor}{E} & F & \textcolor{gregoriocolor}{F} & \textcolor{gregoriocolor}{G} & \textcolor{gregoriocolor}{H} & \textcolor{gregoriocolor}{M} & \textcolor{gregoriocolor}{N} & \textcolor{gregoriocolor}{P} \\
 28 & 29 & 30 & 1 & 2 & 3 & 3 & 4 & 5 & 6 & 7 & 8 \\
\end{tabularx}

\begin{paracol}{2}
\selectlanguage{latin}
\lettrine[lines=2]{R}{omæ} sancti Joánnis, 
 Apóstoli et Evangelístæ, ante Portam Latínam; qui, ab Epheso, jussu 
 Domitiáni, vinctus Romam est perdúctus, et, judicánte Senátu, ante eándem 
 portam in ólei fervéntis dólium missus, exívit inde púrior et vegétior quam 
 intrávit.
\switchcolumn
\selectlanguage{english}
\lettrine[lines=2]{A}{t} Rome, the Apostle and Evangelist 
 St. John before the Latin Gate. He was bound and brought to Rome from 
 Ephesus by the order of Domitian, and the Senate condemned him to be taken 
 to that gate and placed in a cauldron of boiling oil, from which he came 
 forth more healthy and vigorous than before.
\switchcolumn*
\selectlanguage{latin}
Damásci natális beáti 
 Joánnis Damascéni, Presbyteri, Confessóris et Ecclésiæ Doctóris, doctrína et 
 sanctitáte célebris. Hic, pro cultu sanctárum Imáginum, verbo et 
 scriptis advérsus Leónem Isáuricum strénue decertávit; cujus Imperatóris ob 
 calúmnias cum ipsi Joánni déxtera manus e Saracenórum Príncipe amputáta 
 esset, idem, beátæ Maríæ Vírgini, cujus Imágines 
 defénderat, se comméndans, 
 prótinus déxteram íntegram sanámque recépit. Ejus autem festívitas 
 sexto Kaléndas Aprílis celebrátur.
\switchcolumn
\selectlanguage{english}
At Damascus, the birthday of St. 
 John Damascene, priest and doctor of the Church, renowned for sanctity and 
 learning. By means of his writing and preaching, he courageously 
 resisted Leo the Isaurian, in defending the veneration paid to sacred 
 images. By order of this emperor his right hand was cut off, but 
 commending himself before an image of the Blessed Virgin Mary, which he had 
 defended, his hand was immediately restored to him, entire and sound. 
 His feast day is the 27th of March.
\switchcolumn*
\selectlanguage{latin}
Cyréne, in Líbya, sancti Lúcii Epíscopi, quem in Actibus Apostolórum sanctus Lucas commémorat.
\switchcolumn
\selectlanguage{english}
At Cyrene in Africa, Bishop St. 
 Lucius, who is mentioned by St. Luke in the Acts of the Apostles.
\switchcolumn*
\selectlanguage{latin}
Antiochíæ sancti 
 Evódii, qui (ut beátus Ignátius ad Antiochénses scribit), primus ibídem a 
 sancto Petro Apóstolo ordinátus Epíscopus, glorióso martyrio vitam finívit.
\switchcolumn
\selectlanguage{english}
At Antioch, St. Evodius, who, as 
 the blessed Ignatius wrote to the people of Antioch, was consecrated first 
 bishop of that city by the apostle St. Peter, and ended his life by a 
 glorious martyrdom.
\switchcolumn*
\selectlanguage{latin}
In Africa sanctórum 
 Mártyrum Heliodóri et Venústi, cum áliis septuagínta quinque
\switchcolumn
\selectlanguage{english}
In Africa, the holy martyrs 
 Heliodorus and Venustus and seventy-five others.
\switchcolumn*
\selectlanguage{latin}
In Cypro sancti 
 Theódoti, Epíscopi Cyríniæ, qui, sub Licínio Imperatóre, gravíssima passus 
 est, ac tandem, in Ecclésiæ pace, spíritum Deo réddidit.
\switchcolumn
\selectlanguage{english}
In Cyprus, St. Theodotus, bishop of 
 Cyrinia, who having undergone grievous afflictions under Emperor Licinius, 
 at length yielded his soul to God when peace was restored to the Church.
\switchcolumn*
\selectlanguage{latin}
Carrhis, in 
 Mesopotámia, sancti Protógenis, Epíscopi et Confessóris.
\switchcolumn
\selectlanguage{english}
At Carrhae in Mesopotamia, St. 
 Protogenes, bishop and confessor.
\switchcolumn*
\selectlanguage{latin}
In Anglia, sancti 
 Eadbérti, Epíscopi Lindisfarnénsis, doctrína et pietáte insígnis.
\switchcolumn
\selectlanguage{english}
In England, St. Eadbert, bishop of 
 Lindisfarne, famed for his teachings and his piety.
\switchcolumn*
\selectlanguage{latin}
Romæ sanctæ Benedíctæ 
 Vírginis.
\switchcolumn
\selectlanguage{english}
At Rome, the virgin St. Benedicta.
\switchcolumn*
\selectlanguage{latin}
Salérni Translátio 
 sancti Matthǽi, Apóstoli et Evangelístæ; cujus sacrum corpus, olim ex 
 Æthiópia ad divérsas regiónes et demum ad eam urbem delátum, ibídem, in 
 dedicáta ejus nómine Ecclésia, summo honóre cónditum fuit.
\switchcolumn
\selectlanguage{english}
At Salerno, the translation of St. 
 Matthew, apostle and evangelist. His revered body, previously 
 transferred from Ethiopia to various countries, was finally taken to 
 Salerno, and with great pomp was there placed in a church dedicated to his 
 name.
\switchcolumn*
\selectlanguage{latin}
\end{paracol}


% ---- martyrology/mart05/mart0507.htm
\needspace{10\baselineskip}
\begin{paracol}{2}
\selectlanguage{latin}
\begin{center}{\color{gregoriocolor} Nonis Maji. 
 Luna\dots\ }\end{center}
\switchcolumn
\selectlanguage{english}
\begin{center}{\color{gregoriocolor} The 
 7\textsuperscript{th} Day of May. The\dots\ Day of the Moon.}\end{center}
\end{paracol}

\noindent\begin{tabularx}{\linewidth}{*{19}{>{\centering\arraybackslash}X}}
 \textcolor{gregoriocolor}{a} & \textcolor{gregoriocolor}{b} & \textcolor{gregoriocolor}{c} & \textcolor{gregoriocolor}{d} & \textcolor{gregoriocolor}{e} & \textcolor{gregoriocolor}{f} & \textcolor{gregoriocolor}{g} & \textcolor{gregoriocolor}{h} & \textcolor{gregoriocolor}{i} & \textcolor{gregoriocolor}{k} & \textcolor{gregoriocolor}{l} & \textcolor{gregoriocolor}{m} & \textcolor{gregoriocolor}{n} & \textcolor{gregoriocolor}{p} & \textcolor{gregoriocolor}{q} & \textcolor{gregoriocolor}{r} & \textcolor{gregoriocolor}{s} & \textcolor{gregoriocolor}{t} & \textcolor{gregoriocolor}{u} \\
 10 & 11 & 12 & 13 & 14 & 15 & 16 & 17 & 18 & 19 & 20 & 21 & 22 & 23 & 24 & 25 & 26 & 27 & 28 \\
\end{tabularx}
\vspace{0.5\baselineskip}
\noindent\begin{tabularx}{\linewidth}{*{12}{>{\centering\arraybackslash}X}}
 \textcolor{gregoriocolor}{A} & \textcolor{gregoriocolor}{B} & \textcolor{gregoriocolor}{C} & \textcolor{gregoriocolor}{D} & \textcolor{gregoriocolor}{E} & F & \textcolor{gregoriocolor}{F} & \textcolor{gregoriocolor}{G} & \textcolor{gregoriocolor}{H} & \textcolor{gregoriocolor}{M} & \textcolor{gregoriocolor}{N} & \textcolor{gregoriocolor}{P} \\
 29 & 30 & 1 & 2 & 3 & 4 & 4 & 5 & 6 & 7 & 8 & 9 \\
\end{tabularx}

\begin{paracol}{2}
\selectlanguage{latin}
\lettrine[lines=2]{S}{ancti} Stanislái, 
 Epíscopi Cracoviénsis et Mártyris, qui sequénti die, corónam martyrii 
 consecútus est.
\switchcolumn
\selectlanguage{english}
\lettrine[lines=2]{S}{t.} Stanislas, bishop of Cracow 
 and martyr, who received the crown of martyrdom on the day following this.
\switchcolumn*
\selectlanguage{latin}
Tarracínæ, in Campánia, 
 natális beátæ Fláviæ Domitíllæ, Vírginis et Mártyris, quæ, cum esset fília 
 sanctæ Plautíllæ, soróris sancti Mártyris Flávii Cleméntis Cónsulis, et 
 sacro velámine fuísset a sancto Cleménte Pontífice consecráta, primum, in 
 persecutióne Domitiáni, ob testimónium Christi, in ínsulam Póntiam, cum 
 áliis plúrimis, exsílio deportáta, longum illic martyrium duxit. 
 Novíssime vero, Tarracínam dedúcta, ibi, cum plúrimos doctrína et miráculis 
 ad Christi fidem convertísset. Júdicis jussu incénso cubículo, in quo 
 simul cum suis virgínibus Euphrósyna et Theodóra morabátur, cursum gloriósi 
 martyrii consummávit. Ipsa vero Domitílla, una cum sanctis Martyribus 
 Néreo et Achílleo atque Pancrátio, festíva celebritáte quarto Idus mensis 
 hujus recólitur.
\switchcolumn
\selectlanguage{english}
At Terracina in Campania, the 
 birthday of blessed Flavia Domitilla, virgin and martyr, and niece of the 
 holy martyr, the Consul Flavius Clemens. She received the religious 
 veil at the hands of St. Clement, and in the persecution of Domitian was exiled with many others to the island of Pontia, where endured a long 
 martyrdom for Christ. Taken afterwards to Terracina, she converted 
 many to the faith of Christ by her teachings and miracles. The judge 
 ordered the room in which she was with the virgins Euphrosina and Theodora, 
 to be set on fire, and she thus completed her glorious martyrdom. She 
 is also mentioned with the holy martyrs Nereus, Achilleus and Pancras, on 
 the 12th day of this month.
\switchcolumn*
\selectlanguage{latin}
Eódem die sancti 
 Juvenális Mártyris.
\switchcolumn
\selectlanguage{english}
On the same day, St. Juvenal, 
 martyr.
\switchcolumn*
\selectlanguage{latin}
Nicomedíæ sanctórum 
 Mártyrum fratrum Flávii, Augústi et Augustíni.
\switchcolumn
\selectlanguage{english}
At Nicomedia, the holy martyrs 
 Flavius, Augustus and Augustine, all brothers.
\switchcolumn*
\selectlanguage{latin}
Ibídem sancti Quadráti 
 Mártyris, qui, in persecutióne Décii Imperatóris, sæpius ad torménta 
 repetítus, demum, cápite truncátus, martyrium complévit.
\switchcolumn
\selectlanguage{english}
In the same city, St. Quadratus, 
 martyr, who was frequently tortured in the persecution of Decius, and at 
 last beheaded.
\switchcolumn*
\selectlanguage{latin}
Eboráci, in Anglia, 
 sancti Joánnis Epíscopi, vita et miráculis clari.
\switchcolumn
\selectlanguage{english}
At York in England, St. John, 
 bishop, renowned for a saintly life and miracles.
\switchcolumn*
\selectlanguage{latin}
Papíæ sancti Petri 
 Epíscopi.
\switchcolumn
\selectlanguage{english}
At Pavia, Bishop St. Peter.
\switchcolumn*
\selectlanguage{latin}
Romæ Translátio 
 córporis sancti Stéphani Protomártyris, quod, Pelágio Primo Summo Pontífice, 
 e Constantinópoli ad Urbem allátum atque in sepúlcro sancti Lauréntii 
 Mártyris in agro Veráno pósitum, ibídem magna piórum religióne cólitur.
\switchcolumn
\selectlanguage{english}
At Rome, the translation of the 
 body of St. Stephen protomartyr, which was brought from Constantinople to 
 Rome by Pope Pelagius I, and laid in the sepulchre of the martyr St. Lawrence in the Agro 
 Verano, where it is honoured with great devotion by the pious faithful.
\switchcolumn*
\selectlanguage{latin}
\end{paracol}


% ---- martyrology/mart05/mart0508.htm
\needspace{10\baselineskip}
\begin{paracol}{2}
\selectlanguage{latin}
\begin{center}{\color{gregoriocolor} Octávo Idus Maji. 
 Luna\dots\ }\end{center}
\switchcolumn
\selectlanguage{english}
\begin{center}{\color{gregoriocolor} The 
 8\textsuperscript{th} Day of May. The\dots\ Day of the Moon.}\end{center}
\end{paracol}

\noindent\begin{tabularx}{\linewidth}{*{19}{>{\centering\arraybackslash}X}}
 \textcolor{gregoriocolor}{a} & \textcolor{gregoriocolor}{b} & \textcolor{gregoriocolor}{c} & \textcolor{gregoriocolor}{d} & \textcolor{gregoriocolor}{e} & \textcolor{gregoriocolor}{f} & \textcolor{gregoriocolor}{g} & \textcolor{gregoriocolor}{h} & \textcolor{gregoriocolor}{i} & \textcolor{gregoriocolor}{k} & \textcolor{gregoriocolor}{l} & \textcolor{gregoriocolor}{m} & \textcolor{gregoriocolor}{n} & \textcolor{gregoriocolor}{p} & \textcolor{gregoriocolor}{q} & \textcolor{gregoriocolor}{r} & \textcolor{gregoriocolor}{s} & \textcolor{gregoriocolor}{t} & \textcolor{gregoriocolor}{u} \\
 11 & 12 & 13 & 14 & 15 & 16 & 17 & 18 & 19 & 20 & 21 & 22 & 23 & 24 & 25 & 26 & 27 & 28 & 29 \\
\end{tabularx}
\vspace{0.5\baselineskip}
\noindent\begin{tabularx}{\linewidth}{*{12}{>{\centering\arraybackslash}X}}
 \textcolor{gregoriocolor}{A} & \textcolor{gregoriocolor}{B} & \textcolor{gregoriocolor}{C} & \textcolor{gregoriocolor}{D} & \textcolor{gregoriocolor}{E} & F & \textcolor{gregoriocolor}{F} & \textcolor{gregoriocolor}{G} & \textcolor{gregoriocolor}{H} & \textcolor{gregoriocolor}{M} & \textcolor{gregoriocolor}{N} & \textcolor{gregoriocolor}{P} \\
 30 & 1 & 2 & 3 & 4 & 5 & 5 & 6 & 7 & 8 & 9 & 10 \\
\end{tabularx}

\begin{paracol}{2}
\selectlanguage{latin}
\lettrine[lines=2]{I}{n} monte Gargáno 
 Apparítio sancti Michaélis Archángeli, quem Pius Papa Duodécimus Radiólogis 
 et Radiumtherapéuticis Patrónum et Protectórum constítuit.
\switchcolumn
\selectlanguage{english}
\lettrine[lines=2]{O}{n} Mount Gargano, the apparition of 
 St. Michael Archangel, whom Pope Pius XII named the patron and protector of 
 radiologists and radiotherapists.
\switchcolumn*
\selectlanguage{latin}
Cracóviæ, in Polónia, 
 natális sancti Stanislái, Epíscopi et Mártyris, qui a Bolesláo, ímpio Rege, 
 necátus est. Ipsíus autem festum prídie hujus diéi celebrátur.
\switchcolumn
\selectlanguage{english}
At Cracow in Poland, the birthday 
 of St. Stanislas, bishop and martyr, who was slain by the wicked King 
 Boleslas. His feast was celebrated on the previous day.
\switchcolumn*
\selectlanguage{latin}
Medioláni item natális 
 sancti Victóris Mártyris, qui, natióne Maurus et a primæva ætate Christiánus, 
 a Maximiáno, cum esset in castris imperiálibus miles, compúlsus ut idólis 
 sacrificáret, et in confessióne Dómini fortíssime persevérans, ideo, primum 
 gráviter fústibus cæsus, sed, Deo protegénte, dolóris expers; deínde 
 liquénti plumbo perfúsus, sed nihil pénitus læsus; novíssime gloriósi 
 martyrii cursum, cápite abscíssus, implévit.
\switchcolumn
\selectlanguage{english}
At Milan, the birthday of the holy 
 martyr Victor, a Moor. He became a Christian in his youth and served 
 in the imperial army. When Maximian wished to force him to offer 
 sacrifice to idols, he persevered with the greatest fortitude in the 
 confession of the Lord. He was first beaten with rods, but by God's 
 protection without feeling any pain. Following this, melted lead was 
 poured over him, which did him no injury whatever. The career of his 
 glorious martyrdom was finally ended by his being beheaded.
\switchcolumn*
\selectlanguage{latin}
Constantinópoli sancti 
 Agáthii Centuriónis, qui, in persecutióne Diocletiáni et Maximiáni, a Firmo 
 Tribúno delátus quod Christiánus esset, et a Júdice Perínthi Bibiáno 
 sævíssime tortus, Byzántii demum a Procónsule Flaccíno cápitis damnátus est. 
 Ipsíus corpus ad Scyllácium littus, in Calábria, divínitus póstea delátum 
 est, atque ibi honorífice asservátum.
\switchcolumn
\selectlanguage{english}
At Constantinople, St. Acathius, 
 who, being denounced as a Christian by the tribune Firmus, and cruelly 
 tortured at Perinthus by the judge Bibian, was finally condemned to death at 
 Byzantium by the procunsul Flaccinus. His body was afterwards 
 miraculously brought to the shore of Squillace in Calabria, where it is 
 preserved with honour.
\switchcolumn*
\selectlanguage{latin}
Romæ sancti Bonifátii 
 Papæ Quarti, qui Pántheon in honórem beátæ Maríæ ad Mártyres dedicávit.
\switchcolumn
\selectlanguage{english}
At Rome, Pope St. Boniface IV, who 
 dedicated the Pantheon to the honour of our Lady and the martyrs.
\switchcolumn*
\selectlanguage{latin}
Item Romæ sancti 
 Benedícti Secúndi, Papæ et Confessóris.
\switchcolumn
\selectlanguage{english}
Also at Rome, St. Benedict II, pope 
 and confessor.
\switchcolumn*
\selectlanguage{latin}
Viénnæ, in Gállia, 
 sancti Dionysii, Epíscopi et Confessóris.
\switchcolumn
\selectlanguage{english}
At Vienne in France, St. Denis, 
 bishop and confessor.
\switchcolumn*
\selectlanguage{latin}
Antisiodóri sancti 
 Helládii Epíscopi.
\switchcolumn
\selectlanguage{english}
At Auxerre, St. Helladius, bishop.
\switchcolumn*
\selectlanguage{latin}
In monastério Bellæ 
 Vallis, in território Bisuntíno, sancti Petri, qui ex Mónacho Cisterciénsi 
 factus est Tarentasiénsis in Sabáudia Epíscopus.
\switchcolumn
\selectlanguage{english}
In the monastery of Bella Vallis, 
 in the diocese of Besançon, St. Peter, Cistercian monk, who was made bishop 
 of Tarantaise in Savoy.
\switchcolumn*
\selectlanguage{latin}
Apud Ruræmóndam, in 
 Géldria, sancti Wirónis, Epíscopi Scoti.
\switchcolumn
\selectlanguage{english}
At Ruremonde in Holland, St. Wiro, 
 bishop of Scotland.
\switchcolumn*
\selectlanguage{latin}
\end{paracol}


% ---- martyrology/mart05/mart0509.htm
\needspace{10\baselineskip}
\begin{paracol}{2}
\selectlanguage{latin}
\begin{center}{\color{gregoriocolor} Séptimo Idus Maji. 
 Luna\dots\ }\end{center}
\switchcolumn
\selectlanguage{english}
\begin{center}{\color{gregoriocolor} The 
 9\textsuperscript{th} Day of May. The\dots\ Day of the Moon.}\end{center}
\end{paracol}

\noindent\begin{tabularx}{\linewidth}{*{19}{>{\centering\arraybackslash}X}}
 \textcolor{gregoriocolor}{a} & \textcolor{gregoriocolor}{b} & \textcolor{gregoriocolor}{c} & \textcolor{gregoriocolor}{d} & \textcolor{gregoriocolor}{e} & \textcolor{gregoriocolor}{f} & \textcolor{gregoriocolor}{g} & \textcolor{gregoriocolor}{h} & \textcolor{gregoriocolor}{i} & \textcolor{gregoriocolor}{k} & \textcolor{gregoriocolor}{l} & \textcolor{gregoriocolor}{m} & \textcolor{gregoriocolor}{n} & \textcolor{gregoriocolor}{p} & \textcolor{gregoriocolor}{q} & \textcolor{gregoriocolor}{r} & \textcolor{gregoriocolor}{s} & \textcolor{gregoriocolor}{t} & \textcolor{gregoriocolor}{u} \\
 12 & 13 & 14 & 15 & 16 & 17 & 18 & 19 & 20 & 21 & 22 & 23 & 24 & 25 & 26 & 27 & 28 & 29 & 30 \\
\end{tabularx}
\vspace{0.5\baselineskip}
\noindent\begin{tabularx}{\linewidth}{*{12}{>{\centering\arraybackslash}X}}
 \textcolor{gregoriocolor}{A} & \textcolor{gregoriocolor}{B} & \textcolor{gregoriocolor}{C} & \textcolor{gregoriocolor}{D} & \textcolor{gregoriocolor}{E} & F & \textcolor{gregoriocolor}{F} & \textcolor{gregoriocolor}{G} & \textcolor{gregoriocolor}{H} & \textcolor{gregoriocolor}{M} & \textcolor{gregoriocolor}{N} & \textcolor{gregoriocolor}{P} \\
 1 & 2 & 3 & 4 & 5 & 6 & 6 & 7 & 8 & 9 & 10 & 11 \\
\end{tabularx}

\begin{paracol}{2}
\selectlanguage{latin}
\lettrine[lines=2]{N}{aziánzi,} in 
 Cappadócia, natális beáti Gregórii Epíscopi, Confessóris et Ecclésiæ 
 Doctóris, ob singulárem divinárum rerum doctrínam cognoménto Theólogi; qui 
 collápsam Constantinópoli cathólicam fidem, ipsíus urbis Episcopátum gerens, 
 restítuit, hæresésque insurgéntes compréssit.
\switchcolumn
\selectlanguage{english}
\lettrine[lines=2]{A}{t} Nazianzum, the birthday of St. 
 Gregory, bishop, confessor, and doctor of the Church, surnamed the 
 Theologian because of his remarkable knowledge of divinity. At 
 Constantinople, he restored the Catholic faith which was fast waning, and 
 repressed the rising heresies.
\switchcolumn*
\selectlanguage{latin}
Romæ sancti Hermæ, 
 cujus Apóstolus Paulus in Epístola ad Romános méminit. Ipse autem 
 Hermas, digne semetípsum sacríficans acceptabilísque Deo hóstia factus, 
 virtútibus clarus cæléstia regna petívit.
\switchcolumn
\selectlanguage{english}
At Rome, St. Hermas, mentioned by 
 the apostle St. Paul in the Epistle to the Romans. Generously 
 sacrificing himself, he became an offering acceptable to God, and 
 outstanding for his virtues he took his departure for the heavenly kingdom.
\switchcolumn*
\selectlanguage{latin}
Cállii, via Flamínia, 
 pássio sancti Geróntii, Epíscopi Ficoclénsis.
\switchcolumn
\selectlanguage{english}
At Cagli, on the Flaminian Way, the 
 passion of St. Gerontius, bishop of Cervia.
\switchcolumn*
\selectlanguage{latin}
In Pérside sanctórum 
 Mártyrum trecentórum et decem.
\switchcolumn
\selectlanguage{english}
In Persia, three hundred and ten 
 holy martyrs.
\switchcolumn*
\selectlanguage{latin}
In Ægypto sancti 
 Pachómii Abbátis, qui plúrima eréxit in ea regióne monasteria, et régulam 
 Monachórum scripsit, quam ab Angelo dictánte didícerat.
\switchcolumn
\selectlanguage{english}
In Egypt, the abbot St. Pachomius, 
 who founded many monasteries in that country, and wrote a rule for monks 
 which he had learned from the dictation of an angel.
\switchcolumn*
\selectlanguage{latin}
In castro Vindecíno, 
 in Gállia, deposítio sancti Beáti Confessóris.
\switchcolumn
\selectlanguage{english}
In the town of Windisch in France, 
 the death of St. Beatus, confessor.
\switchcolumn*
\selectlanguage{latin}
Bonóniæ beáti Nicolái 
 Albergáti, Mónachi Carthusiáni, ejúsdem civitátis Epíscopi et sanctæ Románæ 
 Ecclésiæ Cardinális, sanctitáte et Apostólicis Legatiónibus clari; cujus 
 corpus Floréntiæ, apud Carthusiános, cónditum est.
\switchcolumn
\selectlanguage{english}
At Bologna, blessed Nicholas 
 Albergati, a Carthusian monk, bishop of that city, and cardinal of the Holy 
 Roman Church, celebrated for his sanctity and and for his work as an 
 apostolic legate. His body was buried at Florence in the monastery of 
 the Carthusians.
\switchcolumn*
\selectlanguage{latin}
Constantinópoli 
 Translátio sanctórum Andréæ Apóstoli, et Lucæ Evangelístæ de Achája; et 
 Timóthei, uníus ex discípulis beáti Pauli Apóstoli, ab Epheso. Corpus 
 autem sancti Andréæ, longo post témpore Amálphim devéctum, ibi pio fidélium 
 concúrsu honorátur; ex cujus sepúlcro liquor ad languóres curándos júgiter 
 manat.
\switchcolumn
\selectlanguage{english}
At Constantinople, the translation 
 of the apostle St. Andrew and the evangelist St. Luke, out of Achaia, and of 
 Timothy, disciple of the blessed apostle Paul, from Ephesus. The body 
 of St. Andrew, long after, was conveyed to Amalfi, where it is honoured by 
 the pious gatherings of the faithful. From his tomb there continually 
 flows a liquid which heals diseases.
\switchcolumn*
\selectlanguage{latin}
Romæ item Translátio 
 sancti Hierónymi Presbyteri, Confessóris et Ecclésiæ Doctóris, ex Béthlehem Judæ ad Basílicam sanctæ Maríæ ad Præsépe.
\switchcolumn
\selectlanguage{english}
At Rome, also, the translation of 
 St. Jerome, priest, confessor, and doctor of the Church. His body was 
 taken from Bethlehem of Judea to the basilica of St. Mary of the Manger.
\switchcolumn*
\selectlanguage{latin}
Bárii quoque, in 
 Apúlia, Translátio sancti Nicolái, Epíscopi et Confessóris, ex Myra, 
 civitáte Lyciæ.
\switchcolumn
\selectlanguage{english}
At Bari in Apulia, the translation 
 also of St. Nicholas, bishop and confessor, from Myra, a city of Lycia.
\switchcolumn*
\selectlanguage{latin}
\end{paracol}


% ---- martyrology/mart05/mart0510.htm
\needspace{10\baselineskip}
\begin{paracol}{2}
\selectlanguage{latin}
\begin{center}{\color{gregoriocolor} Sexto Idus Maji. 
 Luna\dots\ }\end{center}
\switchcolumn
\selectlanguage{english}
\begin{center}{\color{gregoriocolor} The 
 10\textsuperscript{th} Day of May. The\dots\ Day of the Moon.}\end{center}
\end{paracol}

\noindent\begin{tabularx}{\linewidth}{*{19}{>{\centering\arraybackslash}X}}
 \textcolor{gregoriocolor}{a} & \textcolor{gregoriocolor}{b} & \textcolor{gregoriocolor}{c} & \textcolor{gregoriocolor}{d} & \textcolor{gregoriocolor}{e} & \textcolor{gregoriocolor}{f} & \textcolor{gregoriocolor}{g} & \textcolor{gregoriocolor}{h} & \textcolor{gregoriocolor}{i} & \textcolor{gregoriocolor}{k} & \textcolor{gregoriocolor}{l} & \textcolor{gregoriocolor}{m} & \textcolor{gregoriocolor}{n} & \textcolor{gregoriocolor}{p} & \textcolor{gregoriocolor}{q} & \textcolor{gregoriocolor}{r} & \textcolor{gregoriocolor}{s} & \textcolor{gregoriocolor}{t} & \textcolor{gregoriocolor}{u} \\
 13 & 14 & 15 & 16 & 17 & 18 & 19 & 20 & 21 & 22 & 23 & 24 & 25 & 26 & 27 & 28 & 29 & 30 & 1 \\
\end{tabularx}
\vspace{0.5\baselineskip}
\noindent\begin{tabularx}{\linewidth}{*{12}{>{\centering\arraybackslash}X}}
 \textcolor{gregoriocolor}{A} & \textcolor{gregoriocolor}{B} & \textcolor{gregoriocolor}{C} & \textcolor{gregoriocolor}{D} & \textcolor{gregoriocolor}{E} & F & \textcolor{gregoriocolor}{F} & \textcolor{gregoriocolor}{G} & \textcolor{gregoriocolor}{H} & \textcolor{gregoriocolor}{M} & \textcolor{gregoriocolor}{N} & \textcolor{gregoriocolor}{P} \\
 2 & 3 & 4 & 5 & 6 & 7 & 7 & 8 & 9 & 10 & 11 & 12 \\
\end{tabularx}

\begin{paracol}{2}
\selectlanguage{latin}
\lettrine[lines=2]{S}{ancti} Antoníni, ex 
 Ordine Prædicatórum, Epíscopi Florentíni et Confessóris, cujus dies natális 
 sexto Nonas mensis hujus recensétur.
\switchcolumn
\selectlanguage{english}
\lettrine[lines=2]{S}{t.} Antoninus of the Order of 
 Preachers, confessor and archbishop of Florence, whose birthday is the 2nd 
 of May.
\switchcolumn*
\selectlanguage{latin}
Romæ, via Latína, natális sanctórum Mártyrum Gordiáni et Epímachi, quorum prior, pro 
 confessióne nóminis Christi, témpore Juliáni Apóstatæ, diu plumbátis cæsus 
 et ad últimum cápite truncátus, noctu a Christiánis sepúltus eádem via fuit 
 in crypta, in quam beáti Epímachi Mártyris relíquiæ paulo ante translátæ 
 fúerant ab Alexandría, ubi ipse, pro Christi fide, martyrium compléverat 
 prídie Idus Decémbris.
\switchcolumn
\selectlanguage{english}
At Rome, on the Via Latina, the 
 birthday of the holy martyrs Gordian and Epimachus. In the time of 
 Julian the Apostate, the former was a long time scourged and finally 
 beheaded for confessing the name of Christ. He was buried at night by 
 the Christians, in a crypt to which, shortly before, the remains of the 
 blessed martyr Epimachus had been transferred from Alexandria, where he had 
 been martyred for the faith of Christ on the 12th of December.
\switchcolumn*
\selectlanguage{latin}
In terra Hus sancti 
 Job Prophétæ, admirándæ patiéntiæ viri.
\switchcolumn
\selectlanguage{english}
In the land of Hus, the holy 
 prophet Job, a man of wonderful patience.
\switchcolumn*
\selectlanguage{latin}
Romæ beáti Calepódii, 
 Presbyteri et Mártyris; quem Alexánder Imperátor gládio fecit occídi, et 
 corpus ejus per civitátem trahi, atque in Tíberim jactári, quod invéntum 
 Callístus Papa sepelívit. Decollátus est étiam Palmátius Consul cum 
 uxóre et fíliis et áliis promíscui sexus quadragínta duóbus de domo sua, 
 Simplícius quoque Senátor cum uxóre et sexagínta octo de famila sua, item et 
 Felix cum uxóre sua Blanda; quorum cápita suspénsa sunt per divérsas portas 
 Urbis, ad exémplum Christianórum.
\switchcolumn
\selectlanguage{english}
At Rome, the blessed priest and 
 martyr Calepodius, who was killed with the sword by order of Emperor 
 Alexander. His body was dragged through the city and thrown into the 
 Tiber. It was afterwards found and buried by Pope Callistus. The 
 consul Palmatius was also beheaded with his wife, his sons, and forty-two of 
 both sexes belonging to his household; likewise the senator Simplicius with 
 his wife, and sixty-eight of his house; Felix also with his wife Blanda. 
 The heads of all these martyrs were exposed over different gates of the city 
 in order to terrify the Christians.
\switchcolumn*
\selectlanguage{latin}
Item Romæ, via Latína, 
 ad Centum Aulas, natális sanctórum Mártyrum Quarti et Quincti, quorum 
 córpora Cápuam transláta sunt.
\switchcolumn
\selectlanguage{english}
Also at Rome, on the Via Latina, 
 the birthday of the holy martyrs Quartus and Quinctus, whose bodies were 
 translated to Capua.
\switchcolumn*
\selectlanguage{latin}
Apud Leontínos, in 
 Sicília, sanctórum Mártyrum Alphii, Philadélphi et Cyríni.
\switchcolumn
\selectlanguage{english}
At Lentini in Sicily, the holy 
 martyrs Alphius, Philadelphis, and Cyrinus.
\switchcolumn*
\selectlanguage{latin}
Smyrnæ sancti 
 Dioscóridis Mártyris.
\switchcolumn
\selectlanguage{english}
At Smyrna, St. Dioscorides, martyr.
\switchcolumn*
\selectlanguage{latin}
Apud Taréntum sancti 
 Catáldi Epíscopi, miráculis clari.
\switchcolumn
\selectlanguage{english}
At Taranto, St. Cataldus, a bishop 
 renowned for miracles.
\switchcolumn*
\selectlanguage{latin}
Matríti sancti Isidóri 
 Agrícolæ, quem, miráculis clarum, Gregórius Papa Décimus quintus, una cum 
 sanctis Ignátio, Francísco Xavério, Terésia et Philíppo Nério, in Sanctórum 
 númerum rétulit.
\switchcolumn
\selectlanguage{english}
At Madrid, St. Isidore the Farmer. 
 Being well known for his miracles, Pope Gregory XV placed him in the number 
 of saints at the same time with St. Ignatius, St. Francis Xavier, St. 
 Teresa, and St. Philip Neri.
\switchcolumn*
\selectlanguage{latin}
Medioláni Invéntio 
 sanctórum Mártyrum Nazárii et Celsi, in qua beátus Ambrósius Epíscopus 
 corpus sancti Nazárii recénti adhuc sánguine conspérsum réperit, atque ad 
 Basílicam Apostolórum tránstulit, una cum córpore beáti Celsi púeri, quem 
 idem ipse Nazárius nutríerat, et Anolínus, in Nerónis persecutióne, simul 
 cum eo feríri gládio jússerat quinto Kaléndas Augústi; quo die festívitas 
 gloriósi eórum martyrii celebrátur.
\switchcolumn
\selectlanguage{english}
At Milan, the finding of the bodies 
 of the holy martyrs Nazarius and Celsus. The blessed bishop Ambrose 
 found the body of St. Nazarius covered with blood still fresh, and 
 transferred it to the Basilica of the Apostles, together with the body of 
 the blessed Celsus, a youth whom Nazarius had taken care of, and whom 
 Anolinus, in the persecution of Nero, had ordered to be slain with the sword 
 on the 28th of July, on which day their martyrdom is commemorated.
\switchcolumn*
\selectlanguage{latin}
\end{paracol}


% ---- martyrology/mart05/mart0511.htm
\needspace{10\baselineskip}
\begin{paracol}{2}
\selectlanguage{latin}
\begin{center}{\color{gregoriocolor} Quinto Idus Maji. 
 Luna\dots\ }\end{center}
\switchcolumn
\selectlanguage{english}
\begin{center}{\color{gregoriocolor} The 
 11\textsuperscript{th} Day of May. The\dots\ Day of the Moon.}\end{center}
\end{paracol}

\noindent\begin{tabularx}{\linewidth}{*{19}{>{\centering\arraybackslash}X}}
 \textcolor{gregoriocolor}{a} & \textcolor{gregoriocolor}{b} & \textcolor{gregoriocolor}{c} & \textcolor{gregoriocolor}{d} & \textcolor{gregoriocolor}{e} & \textcolor{gregoriocolor}{f} & \textcolor{gregoriocolor}{g} & \textcolor{gregoriocolor}{h} & \textcolor{gregoriocolor}{i} & \textcolor{gregoriocolor}{k} & \textcolor{gregoriocolor}{l} & \textcolor{gregoriocolor}{m} & \textcolor{gregoriocolor}{n} & \textcolor{gregoriocolor}{p} & \textcolor{gregoriocolor}{q} & \textcolor{gregoriocolor}{r} & \textcolor{gregoriocolor}{s} & \textcolor{gregoriocolor}{t} & \textcolor{gregoriocolor}{u} \\
 14 & 15 & 16 & 17 & 18 & 19 & 20 & 21 & 22 & 23 & 24 & 25 & 26 & 27 & 28 & 29 & 30 & 1 & 2 \\
\end{tabularx}
\vspace{0.5\baselineskip}
\noindent\begin{tabularx}{\linewidth}{*{12}{>{\centering\arraybackslash}X}}
 \textcolor{gregoriocolor}{A} & \textcolor{gregoriocolor}{B} & \textcolor{gregoriocolor}{C} & \textcolor{gregoriocolor}{D} & \textcolor{gregoriocolor}{E} & F & \textcolor{gregoriocolor}{F} & \textcolor{gregoriocolor}{G} & \textcolor{gregoriocolor}{H} & \textcolor{gregoriocolor}{M} & \textcolor{gregoriocolor}{N} & \textcolor{gregoriocolor}{P} \\
 3 & 4 & 5 & 6 & 7 & 8 & 8 & 9 & 10 & 11 & 12 & 13 \\
\end{tabularx}

% NOTE: check that this layout is ok on p. 126-127
\begin{paracol}{2}
\selectlanguage{latin}
\lettrine[lines=2]{R}{omæ,} via Salária, 
 natális beáti Anthimi Presbyteri, qui, post virtútum et prædicatiónis 
 insígnia, in persecutióne Diocletiáni, in Tíberim præcipitátus, et ab Angelo 
 exínde eréptus, oratório próprio restitútus est; deínde, cápite punítus, 
 victor migrávit ad cælos.
\switchcolumn
\selectlanguage{english}
\lettrine[lines=2]{A}{t} Rome, on the Salarian Way, the 
 birthday of blessed Anthimus, priest, who, after having distinguished 
 himself by his virtues and preaching, was cast into the Tiber during the 
 persecution of Diocletian. He was rescued by an angel and restored to 
 his oratory. Afterwards he was beheaded, and went victoriously to 
 heaven.
\switchcolumn*
\selectlanguage{latin}
Ibídem sancti Evéllii 
 Mártyris, qui, cum esset de família Nerónis, ad passiónem sancti Torpétis in 
 Christum crédidit, pro quo et decollátus est.
\switchcolumn
\selectlanguage{english}
In the same place, St. Evelius, 
 martyr, who belonged to the household of Nero. By witnessing the 
 martyrdom of St. Torpes, he also believed in Christ, and for him was 
 beheaded.
\switchcolumn*
\selectlanguage{latin}
Item Romæ sanctórum 
 Mártyrum Máximi, Bassi et Fábii; qui sub Diocletiáno, via Salária, cæsi sunt.
\switchcolumn
\selectlanguage{english}
Also at Rome, on the Salarian Way, 
 the holy martyrs Maximus, Bassus, and Fabius, who were put to death during 
 the reign of Diocletian.
\switchcolumn*
\selectlanguage{latin}
Auximi, in Picéno, 
 sanctórum Mártyrum Sisínii Diáconi, Dioclétii et Floréntii, discipulórum 
 sancti Anthimi Presbyteri; qui, sub Diocletiáno, lapídibus óbruti, martyrium 
 complevérunt.
\switchcolumn
\selectlanguage{english}
At Osimo in Piceno, the holy 
 martyrs Sisinius, a deacon, Diocletius and Florentius, disciples of the 
 priest St. Anthimus, whose martyrdom was completed under Diocletian by their 
 being stoned.
\switchcolumn*
\selectlanguage{latin}
Cameríni sanctórum 
 Mártyrum Anastásii et Sociórum; qui, in persecutióne Décii, sub Antíocho Prǽside, cæsi sunt.
\switchcolumn
\selectlanguage{english}
At Camerino, the holy martyrs 
 Anastasius and his companions who were killed in the persecution of Decius, 
 under the governor Antiochus.
\switchcolumn*
\selectlanguage{latin}
Varénnis, in Gállia, 
 sancti Gangúlfi Mártyris.
\switchcolumn
\selectlanguage{english}
At Varennes in France, St. 
 Gangulphus, martyr.
\switchcolumn*
\selectlanguage{latin}
Viénnæ, in Gállia, 
 sancti Mamérti Epíscopi, qui, ob imminéntem cladem, solémnes ante 
 Ascensiónem Dómini triduánas in ea urbe Litanías instítuit; quem ritum 
 póstea universális Ecclésia recípiens comprobávit.
\switchcolumn
\selectlanguage{english}
At Vienne in France, St. Mamertus, 
 bishop, who, to avert an impending calamity, instituted in that city the 
 three days' Litanies immediately before the Ascension of our Lord. 
 This rite was afterwards received and approved by the universal Church.
\switchcolumn*
\selectlanguage{latin}
Apud Silviníacum, in 
 Gállia, deposítio sancti Majóli, Abbátis Cluniacénsis, cujus vita sanctis 
 méritis fuit præclára.
\switchcolumn
\selectlanguage{english}
At Souvigny in France, the death of 
 St. Maieul, abbot of Cluny, whose life was distinguished for merits and 
 sanctity.
\switchcolumn*
\selectlanguage{latin}
Neápoli, in Campánia, 
 sancti Francísci de Hierónymo, in Tarentínæ diœcésis 
 óppido Cryptaleárum 
 orti, Sacerdótis e Societáte Jesu et Confessóris, exímiæ in salúte animárum 
 procuránda caritátis et patiéntiæ viri; quem Gregórius Papa Décimus sextus 
 in Sanctórum cánonem rétulit.
\switchcolumn
\selectlanguage{english}
At Naples in Campania, St. Francis 
 of Jerome, priest of the Society of Jesus, and confessor. He was born 
 in the town of Grottaglia, in the diocese of Taranto. Having been a 
 man of great patience and zeal for the salvation of souls, he was canonized 
 by Pope Gregory XVI.
\switchcolumn*
\selectlanguage{latin}
Apud Septempedános, in 
 Picéno, sancti Illumináti Confessóris.
\switchcolumn
\selectlanguage{english}
At San Severino in Piceno, St. 
 Illuminatus, confessor.
\switchcolumn*
\selectlanguage{latin}
Cálari, in Sardínia, sancti Ignátii a Lacóni, Confessóris, ex Ordine Minórum Capuccinórum, 
 humilitáte, caritáte et miráculis præclári; quem Pius Papa Duodécimus 
 Sanctórum honóribus decorávit.
\switchcolumn
\selectlanguage{english}
At Cagliari in Sardinia, St. 
 Ignatius of Laconi, confessor, of the Minor Order of Capuchins, 
 distinguished for his humility, charity and miracles. He was accorded 
 the honour of canonization by Pope Pius XII.
\switchcolumn*
\selectlanguage{latin}
\end{paracol}


% ---- martyrology/mart05/mart0512.htm
\needspace{10\baselineskip}
\begin{paracol}{2}
\selectlanguage{latin}
\begin{center}{\color{gregoriocolor} Quarto Idus Maji. 
 Luna\dots\ }\end{center}
\switchcolumn
\selectlanguage{english}
\begin{center}{\color{gregoriocolor} The 
 12\textsuperscript{th} Day of May. The\dots\ Day of the Moon.}\end{center}
\end{paracol}

\noindent\begin{tabularx}{\linewidth}{*{19}{>{\centering\arraybackslash}X}}
 \textcolor{gregoriocolor}{a} & \textcolor{gregoriocolor}{b} & \textcolor{gregoriocolor}{c} & \textcolor{gregoriocolor}{d} & \textcolor{gregoriocolor}{e} & \textcolor{gregoriocolor}{f} & \textcolor{gregoriocolor}{g} & \textcolor{gregoriocolor}{h} & \textcolor{gregoriocolor}{i} & \textcolor{gregoriocolor}{k} & \textcolor{gregoriocolor}{l} & \textcolor{gregoriocolor}{m} & \textcolor{gregoriocolor}{n} & \textcolor{gregoriocolor}{p} & \textcolor{gregoriocolor}{q} & \textcolor{gregoriocolor}{r} & \textcolor{gregoriocolor}{s} & \textcolor{gregoriocolor}{t} & \textcolor{gregoriocolor}{u} \\
 15 & 16 & 17 & 18 & 19 & 20 & 21 & 22 & 23 & 24 & 25 & 26 & 27 & 28 & 29 & 30 & 1 & 2 & 3 \\
\end{tabularx}
\vspace{0.5\baselineskip}
\noindent\begin{tabularx}{\linewidth}{*{12}{>{\centering\arraybackslash}X}}
 \textcolor{gregoriocolor}{A} & \textcolor{gregoriocolor}{B} & \textcolor{gregoriocolor}{C} & \textcolor{gregoriocolor}{D} & \textcolor{gregoriocolor}{E} & F & \textcolor{gregoriocolor}{F} & \textcolor{gregoriocolor}{G} & \textcolor{gregoriocolor}{H} & \textcolor{gregoriocolor}{M} & \textcolor{gregoriocolor}{N} & \textcolor{gregoriocolor}{P} \\
 4 & 5 & 6 & 7 & 8 & 9 & 9 & 10 & 11 & 12 & 13 & 14 \\
\end{tabularx}

\begin{paracol}{2}
\selectlanguage{latin}
\lettrine[lines=2]{R}{omæ,} via Ardeatína, sanctórum Mártyrum Nérei et Achíllei 
 fratrum, qui primo cum Flávia Domitílla, cujus erant eunúchi, in ínsula 
 Póntia longum pro Christo duxérunt exsílium; póstmodum gravíssimis 
 verbéribus attrectáti sunt; deínde, cum a Minútio Rufo, viro Consulári, 
 equúleo et flammis ad immolándum compelleréntur, diceréntque se, a beáto 
 Petro Apóstolo baptizátos, nulla ratióne posse idólis immoláre, cápite cæsi 
 sunt. Horum sacræ relíquiæ, simúlque Fláviæ Domitíllæ, ex Diaconía 
 sancti Hadriáni in antíquum eórum Títulum, ubi asservabántur olim recónditæ, 
 dénuo restaurátum, solémniter translátæ sunt prídie hujus diéi, jussu 
 Cleméntis Papæ Octávi; qui exínde hodiérna celebrándum die indíxit étiam 
 festum ipsíus beátæ Domitíllæ Vírginis, cujus pássio Nonis hujus mensis 
 recensétur.
\switchcolumn
\selectlanguage{english}
\lettrine[lines=2]{A}{t} Rome, on the Ardeatine Way, the holy martyrs Nereus 
 and Achilleus, brothers, who underwent a long exile for Christ in the island 
 of Pontia with Flavia Domitilla, whose chamberlains they were. 
 Afterwards they endured a most severe scourging. Finally, as the 
 judge, Minutius Rufus, endeavoured by using the rack and fire to force them 
 to offer sacrifices, they said that having been baptized by the blessed 
 apostle Peter, they could by no means sacrifice to idols. They were 
 beheaded, and their revered remains, with those of Flavia Domitilla, were, 
 by order of Pope Clement VIII, solemnly transferred the day before this, 
 from the sacristy of St. Adrian to the church in which they had been kept in 
 the first place, and which was now repaired. He also ordered today's 
 observance of the feast of St. Domitilla, the virgin, whose martyrdom was 
 mentioned on the 7th of May.
\switchcolumn*
\selectlanguage{latin}
Item Romæ, via Aurélia, sancti Pancrátii Mártyris, qui, 
 cum esset annórum quatuórdecim, sub Diocletiáno, cápitis obtruncatióne 
 martyrium complévit.
\switchcolumn
\selectlanguage{english}
In the same place, on the Aurelian Way, the holy martyr 
 Pancras who at fourteen years of age endured martyrdom by being beheaded 
 under Diocletian.
\switchcolumn*
\selectlanguage{latin}
Salamínæ, in Cypro, sancti Epiphánii Epíscopi, qui, 
 multíplici eruditióne et sacrárum sciéntia litterárum excéllens, vitæ quoque 
 sanctitáte, zelo cathólicæ fídei, munificéntia in páuperes et virtúte 
 miraculórum éxstitit admirándus.
\switchcolumn
\selectlanguage{english}
At Salamis in Cyprus, St. Epiphanius, a bishop of great 
 erudition, with a profound knowledge of the Holy Scriptures. He is to 
 be admired for the sanctity of his life, his zeal for the Catholic faith, 
 his charity to the poor, and the gift of miracles.
\switchcolumn*
\selectlanguage{latin}
Constantinópoli sancti Germáni Epíscopi, doctrína et 
 virtútibus insígnis, qui Leónem Isáuricum, edíctum advérsus sacras Imágines 
 promulgántem, magna fidúcia redárguit.
\switchcolumn
\selectlanguage{english}
At Constantinople, St. Germanus, a bishop distinguished 
 by his virtues and learning, who faithfully opposed Leo the Isaurian for 
 publishing an edict against sacred images.
\switchcolumn*
\selectlanguage{latin}
Tréviris sancti Modoáldi Epíscopi.
\switchcolumn
\selectlanguage{english}
At Treves, St. Modoaldus, bishop.
\switchcolumn*
\selectlanguage{latin}
Romæ sancti Dionysii, qui éxstitit pátruus sancti 
 Pancrátii Mártyris.
\switchcolumn
\selectlanguage{english}
At Rome, St. Denis, uncle of the martyr St. Pancras.
\switchcolumn*
\selectlanguage{latin}
Argyrii, in Sicília, sancti Philíppi Presbyteri, qui, a 
 Románo Pontífice ad eándem Sicíliam ínsulam missus, magnam illíus partem 
 convértit ad Christum. Ipsíus vero sánctitas in liberándis energúmenis 
 máxime declarátur.
\switchcolumn
\selectlanguage{english}
At Agirone in Sicily, St. Philip, a priest who was sent 
 to that island by the Roman Pontiff, and converted to Christ a great portion 
 of it. His sanctity is particularly manifested by the deliverance of 
 persons possessed.
\switchcolumn*
\selectlanguage{latin}
In civitáte 
 Calciaténsi, in Hispánia, sancti Domínici Confessóris.
\switchcolumn
\selectlanguage{english}
In the city of Calzada in Spain, 
 St. Dominic, confessor.
\switchcolumn*
\selectlanguage{latin}
\end{paracol}


% ---- martyrology/mart05/mart0513.htm
\needspace{10\baselineskip}
\begin{paracol}{2}
\selectlanguage{latin}
\begin{center}{\color{gregoriocolor} Tértio Idus Maji. 
 Luna\dots\ }\end{center}
\switchcolumn
\selectlanguage{english}
\begin{center}{\color{gregoriocolor} The 
 13\textsuperscript{th} Day of May. The\dots\ Day of the Moon.}\end{center}
\end{paracol}

\noindent\begin{tabularx}{\linewidth}{*{19}{>{\centering\arraybackslash}X}}
 \textcolor{gregoriocolor}{a} & \textcolor{gregoriocolor}{b} & \textcolor{gregoriocolor}{c} & \textcolor{gregoriocolor}{d} & \textcolor{gregoriocolor}{e} & \textcolor{gregoriocolor}{f} & \textcolor{gregoriocolor}{g} & \textcolor{gregoriocolor}{h} & \textcolor{gregoriocolor}{i} & \textcolor{gregoriocolor}{k} & \textcolor{gregoriocolor}{l} & \textcolor{gregoriocolor}{m} & \textcolor{gregoriocolor}{n} & \textcolor{gregoriocolor}{p} & \textcolor{gregoriocolor}{q} & \textcolor{gregoriocolor}{r} & \textcolor{gregoriocolor}{s} & \textcolor{gregoriocolor}{t} & \textcolor{gregoriocolor}{u} \\
 16 & 17 & 18 & 19 & 20 & 21 & 22 & 23 & 24 & 25 & 26 & 27 & 28 & 29 & 30 & 1 & 2 & 3 & 4 \\
\end{tabularx}
\vspace{0.5\baselineskip}
\noindent\begin{tabularx}{\linewidth}{*{12}{>{\centering\arraybackslash}X}}
 \textcolor{gregoriocolor}{A} & \textcolor{gregoriocolor}{B} & \textcolor{gregoriocolor}{C} & \textcolor{gregoriocolor}{D} & \textcolor{gregoriocolor}{E} & F & \textcolor{gregoriocolor}{F} & \textcolor{gregoriocolor}{G} & \textcolor{gregoriocolor}{H} & \textcolor{gregoriocolor}{M} & \textcolor{gregoriocolor}{N} & \textcolor{gregoriocolor}{P} \\
 5 & 6 & 7 & 8 & 9 & 10 & 10 & 11 & 12 & 13 & 14 & 15 \\
\end{tabularx}

\begin{paracol}{2}
\selectlanguage{latin}
\lettrine[lines=2]{S}{ancti} Robérti 
 Bellarmíno, e Societáte Jesu, Cardinális atque olim Epíscopi Capuáni, 
 Confessóris et Ecclésiæ Doctóris, cujus dies natális décimo quinto Kaléndas 
 Octóbris recensétur.
\switchcolumn
\selectlanguage{english}
\lettrine[lines=2]{S}{t.} Robert Bellarmine, of the 
 Society of Jesus, cardinal and one time bishop of Capua, confessor and 
 doctor of the Church, whose birthday is kept on the 17th of September.
\switchcolumn*
\selectlanguage{latin}
Romæ Dedicátio 
 Ecclésiæ sanctæ Maríæ ad Mártyres, quam beátus Bonifátius Papa Quartus, 
 expurgáto deórum ómnium véteri fano, quod Pántheon vocabátur, in honórem 
 beátæ semper Vírginis Maríæ et ómnium Mártyrum dedicávit, témpore Phocæ 
 Imperatóris. Ipsíus vero Dedicatiónis ánnuam solemnitátem póstmodum 
 Summus Póntifex Gregórius item Quartus ab univérsa Ecclésia, et in honórem 
 quidem ómnium Sanctórum, Kaléndis Novémbris agéndam esse constítuit.
\switchcolumn
\selectlanguage{english}
At Rome, in the time of Emperor 
 Phocas, the dedication of the church of St. Mary of the Martyrs, formerly a 
 temple of all the gods, called the Pantheon, which was purified and 
 dedicated by the blessed Pope Boniface IV to the honour of the Blessed Mary 
 ever Virgin, and of all the martyrs. The solemn anniversary of this 
 dedication was later ordered to be kept by Pope Gregory IV as the Feast of 
 All Saints on the 1st of November.
\switchcolumn*
\selectlanguage{latin}
Constantinópoli beáti 
 Múcii, Presbyteri et Mártyris; qui, sub Diocletiáno Imperatóre et Laudício 
 Procónsule, prius Amphípoli, in Macedónia, multis pœnis atque cruciátibus ob 
 Christi confessiónem afflíctus, póstea, Byzántium usque perdúctus, capitáli 
 senténtia occúbuit.
\switchcolumn
\selectlanguage{english}
At Constantinople, under Emperor 
 Diocletian and the proconsul Laudicius, the blessed Mucius, priest and 
 martyr, who endured many tribulations and torments for the confession of 
 Christ at Amphipolis, and then being taken to Byzantium, suffered death.
\switchcolumn*
\selectlanguage{latin}
Alexandríæ 
 commemorátio plurimórum sanctórum Mártyrum, qui, ob fidem cathólicam, ab 
 Ariánis in Ecclésia Theónæ occísi sunt.
\switchcolumn
\selectlanguage{english}
At Alexandria, the commemoration of 
 many holy martyrs, who were put to death for the Catholic faith by the 
 Arians in the church of St. Theonas.
\switchcolumn*
\selectlanguage{latin}
Heracléæ, in Thrácia, 
 sanctæ Glycériæ, Mártyris Románæ, quæ, sub Antoníno Imperatóre et Sabíno 
 Prǽside, cum fuísset plúrimis ac diris tentáta supplíciis et ab his divína 
 ope incólumis evasísset, tandem feris est objécta, earúmque una, morsum 
 córpori ipsíus infigénte, spíritum Deo réddidit.
\switchcolumn
\selectlanguage{english}
At Heraclea in Thrace, St. Glyceria, 
 a Roman martyr who suffered many severe torments under Emperor Antonius and 
 the governor Sabinus. By the help of God having escaped them all 
 unharmed, she was finally thrown to the wild beasts, and when the first one had 
 bitten her body, she rendered her soul to God.
\switchcolumn*
\selectlanguage{latin}
Apud Trajéctum sancti 
 Servátii, Tungrénsis Ecclésiæ Epíscopi; ad cujus méritum ómnibus 
 demonstrándum, dum témpore híemis ómnia in circúitu nix repléret, sepúlcrum 
 ejus numquam opéruit, donec, ex indústria cívium, Basílica super illud 
 ædificáta est.
\switchcolumn
\selectlanguage{english}
At Utrecht, St. Servatius, bishop 
 of Tongres, whose grave, as a public sign of his merit, was free from snow 
 during winter (although everything around was covered with it), until the 
 inhabitants built a church over it.
\switchcolumn*
\selectlanguage{latin}
In Palæstína sancti 
 Joánnis Silentiárii, qui, Coloniénsi in Arménia Episcopátu dimísso, in 
 sancti Sabbæ laura monásticam vitam duxit, et sancto fine quiévit.
\switchcolumn
\selectlanguage{english}
In Palestine, St. John the Silent, 
 who resigned the see of Colonia in Armenia and retired to the monastery of 
 St. Sabbas until his saintly death.
\switchcolumn*
\selectlanguage{latin}
Pódii, in diœcési 
 Pictaviénsi, sancti Andréæ Hubérti Fournet, Confessóris, olim párochi, 
 Institúti Filiárum a Cruce una cum sancta Elísabeth Bichier des Ages 
 Fundatóris, quem Pius Papa Undécimus Sanctórum fastis adscrípsit.
\switchcolumn
\selectlanguage{english}
At La Puye in the diocese of 
 Poitiers, St. André-Hubert Fournet, confessor and one time parish priest, 
 and founder with St. Elizabeth-Lucie Bichier des Ages of the Institute of 
 the Daughters of the Holy Cross. He was placed on the roll of the 
 saints by Pope Pius XI.
\switchcolumn*
\selectlanguage{latin}
\end{paracol}


% ---- martyrology/mart05/mart0514.htm
\needspace{10\baselineskip}
\begin{paracol}{2}
\selectlanguage{latin}
\begin{center}{\color{gregoriocolor} Prídie Idus Maji. 
 Luna\dots\ }\end{center}
\switchcolumn
\selectlanguage{english}
\begin{center}{\color{gregoriocolor} The 
 14\textsuperscript{th} Day of May. The\dots\ Day of the Moon.}\end{center}
\end{paracol}

\noindent\begin{tabularx}{\linewidth}{*{19}{>{\centering\arraybackslash}X}}
 \textcolor{gregoriocolor}{a} & \textcolor{gregoriocolor}{b} & \textcolor{gregoriocolor}{c} & \textcolor{gregoriocolor}{d} & \textcolor{gregoriocolor}{e} & \textcolor{gregoriocolor}{f} & \textcolor{gregoriocolor}{g} & \textcolor{gregoriocolor}{h} & \textcolor{gregoriocolor}{i} & \textcolor{gregoriocolor}{k} & \textcolor{gregoriocolor}{l} & \textcolor{gregoriocolor}{m} & \textcolor{gregoriocolor}{n} & \textcolor{gregoriocolor}{p} & \textcolor{gregoriocolor}{q} & \textcolor{gregoriocolor}{r} & \textcolor{gregoriocolor}{s} & \textcolor{gregoriocolor}{t} & \textcolor{gregoriocolor}{u} \\
 17 & 18 & 19 & 20 & 21 & 22 & 23 & 24 & 25 & 26 & 27 & 28 & 29 & 30 & 1 & 2 & 3 & 4 & 5 \\
\end{tabularx}
\vspace{0.5\baselineskip}
\noindent\begin{tabularx}{\linewidth}{*{12}{>{\centering\arraybackslash}X}}
 \textcolor{gregoriocolor}{A} & \textcolor{gregoriocolor}{B} & \textcolor{gregoriocolor}{C} & \textcolor{gregoriocolor}{D} & \textcolor{gregoriocolor}{E} & F & \textcolor{gregoriocolor}{F} & \textcolor{gregoriocolor}{G} & \textcolor{gregoriocolor}{H} & \textcolor{gregoriocolor}{M} & \textcolor{gregoriocolor}{N} & \textcolor{gregoriocolor}{P} \\
 6 & 7 & 8 & 9 & 19 & 11 & 11 & 12 & 13 & 14 & 15 & 16 \\
\end{tabularx}

\begin{paracol}{2}
\selectlanguage{latin}
\lettrine[lines=2]{T}{arsi,} in Cilícia, 
 natális sancti Bonifátii Mártyris, qui, sub Diocletiáno et Maximiáno passus, 
 deínde Romam advéctus, via Latína sepúltus est.
\switchcolumn
\selectlanguage{english}
\lettrine[lines=2]{A}{t} Tarsus in Cilicia, the birthday 
 of the holy martyr Boniface, who suffered under Diocletian and Maximian. 
 His body was subsequently taken to Rome and buried on the Via Latina.
\switchcolumn*
\selectlanguage{latin}
In Gállia sancti 
 Póntii Mártyris, cujus prædicatióne et indústria postquam duo Philíppi Cǽsares ad Christi fidem convérsi sunt, ipse, sub Valeriáno et Galliéno 
 Princípibus, martyrii palmam adéptus est.
\switchcolumn
\selectlanguage{english}
In France, St. Pontius, martyr. 
 Having by his preaching and his zeal converted to the faith of Christ the 
 two Caesars Philippi, he obtained the palm of martyrdom under the emperors 
 Valerian and Gallienus.
\switchcolumn*
\selectlanguage{latin}
In Syria sanctórum 
 Mártyrum Victóris et Corónæ, sub Antoníno Imperatóre; ex quibus Victor a 
 Sebastiáno Júdice váriis et horréndis afféctus est cruciátibus. Cum 
 autem ipsum Coróna, uxor cujúsdam mílitis, cœpísset beátum prædicáre ob 
 martyrii constántiam, vidit duas corónas de cælo lapsas, unam Victóri et 
 álteram sibi missam; cumque hoc audiéntibus cunctis testarétur, ipsa quidem 
 inter árbores scissa, Victor vero decollátus est.
\switchcolumn
\selectlanguage{english}
In Syria, the holy martyrs Victor 
 and Corona, under Emperor Antoninus. Victor was subjected to diverse 
 and horrible torments by the judge Sebastian. Just then, as Corona, 
 the the wife of a certain soldier, proclaimed him blessed for his 
 constancy in his sufferings, she saw two crowns falling from heaven, one for 
 Victor, the other for herself. She related this to all present, and 
 was torn to pieces between two trees, while Victor was beheaded.
\switchcolumn*
\selectlanguage{latin}
In Sardínia sanctárum 
 Mártyrum Justæ, Justínæ et Henedínæ.
\switchcolumn
\selectlanguage{english}
In Sardinia, the holy martyrs Justa, 
 Justina, and Henedina.
\switchcolumn*
\selectlanguage{latin}
Ferénti, in Túscia, 
 sancti Bonifátii Epíscopi, qui (ut refert beátus Gregórius Papa) sanctitáte 
 et miráculis a puerítia cláruit.
\switchcolumn
\selectlanguage{english}
At Ferentino in Tuscany, Bishop St. 
 Boniface, who was renowned for sanctity and miracles from his childhood as 
 is told by the blessed Pope Gregory.
\switchcolumn*
\selectlanguage{latin}
In pago Bethárram, 
 diœcésis Bajonénsis, sancti Michaélis Garicoïts, Confessóris, Congregatiónis 
 Presbyterórum Missionariórum a Sacro Corde Jesu Fundatóris, apostólico zelo 
 insígnis, quem Pius Papa Duodécimus Sanctórum fastis adscrípsit.
\switchcolumn
\selectlanguage{english}
In the town of Betharram in the 
 diocese of Bayonne, St. Michael Garricoits, confessor, and founder of the 
 Congregation of the Priests of the Sacred Heart, renowned for his apostolic 
 fervour. Pope Pius XII added him to the roll of saints.
\switchcolumn*
\selectlanguage{latin}
Níciæ, in Subalpínis, 
 sanctæ Maríæ Domínicæ Mazzaréllo, Confundatrícis Institúti Filiárum Maríæ 
 Auxiliatrícis, quæ humilitáte, prudéntia et caritáte præclára, in album 
 sanctárum Vírginum a Pio Papa Duodécimo fuit reláta.
\switchcolumn
\selectlanguage{english}
At Nizza Monferrato in Italy, St. 
 Mary Dominica Mazzarello, co-founder of the Daughters of Mary Help of 
 Christians, and renowned for her humility, prudence and charity. She 
 was added to the book of Virgins by Pope Pius XII.
\switchcolumn*
\selectlanguage{latin}
\end{paracol}


% ---- martyrology/mart05/mart0515.htm
\needspace{10\baselineskip}
\begin{paracol}{2}
\selectlanguage{latin}
\begin{center}{\color{gregoriocolor} Idibus Maji. 
 Luna\dots\ }\end{center}
\switchcolumn
\selectlanguage{english}
\begin{center}{\color{gregoriocolor} The 
 15\textsuperscript{th} Day of May. The\dots\ Day of the Moon.}\end{center}
\end{paracol}

\noindent\begin{tabularx}{\linewidth}{*{19}{>{\centering\arraybackslash}X}}
 \textcolor{gregoriocolor}{a} & \textcolor{gregoriocolor}{b} & \textcolor{gregoriocolor}{c} & \textcolor{gregoriocolor}{d} & \textcolor{gregoriocolor}{e} & \textcolor{gregoriocolor}{f} & \textcolor{gregoriocolor}{g} & \textcolor{gregoriocolor}{h} & \textcolor{gregoriocolor}{i} & \textcolor{gregoriocolor}{k} & \textcolor{gregoriocolor}{l} & \textcolor{gregoriocolor}{m} & \textcolor{gregoriocolor}{n} & \textcolor{gregoriocolor}{p} & \textcolor{gregoriocolor}{q} & \textcolor{gregoriocolor}{r} & \textcolor{gregoriocolor}{s} & \textcolor{gregoriocolor}{t} & \textcolor{gregoriocolor}{u} \\
 18 & 19 & 20 & 21 & 22 & 23 & 24 & 25 & 26 & 27 & 28 & 29 & 30 & 1 & 2 & 3 & 4 & 5 & 6 \\
\end{tabularx}
\vspace{0.5\baselineskip}
\noindent\begin{tabularx}{\linewidth}{*{12}{>{\centering\arraybackslash}X}}
 \textcolor{gregoriocolor}{A} & \textcolor{gregoriocolor}{B} & \textcolor{gregoriocolor}{C} & \textcolor{gregoriocolor}{D} & \textcolor{gregoriocolor}{E} & F & \textcolor{gregoriocolor}{F} & \textcolor{gregoriocolor}{G} & \textcolor{gregoriocolor}{H} & \textcolor{gregoriocolor}{M} & \textcolor{gregoriocolor}{N} & \textcolor{gregoriocolor}{P} \\
 7 & 8 & 9 & 10 & 11 & 12 & 12 & 13 & 14 & 15 & 16 & 17 \\
\end{tabularx}

\begin{paracol}{2}
\selectlanguage{latin}
\lettrine[lines=2]{S}{ancti} Joánnis 
 Baptístæ de La Salle, Presbyteri et Confessóris, qui Sodalitátis Fratrum 
 Scholárum Christianárum fuit Institútor, ac séptimo Idus Aprílis obdormívit 
 in Dómino.
\switchcolumn
\selectlanguage{english}
\lettrine[lines=2]{S}{t.} John Baptist de la Salle, 
 priest and confessor, who founded the Society of Brothers of the Christian 
 Schools. He went to rest in the Lord on the 7th of April.
\switchcolumn*
\selectlanguage{latin}
In Hispánia sanctórum 
 Torquáti, Ctesiphóntis, Secúndi, Indalétii, Cæcílii, Hesychii et Euphrásii; 
 qui Romæ a sanctis Apóstolis Epíscopi ordináti, et ad prædicándum verbum Dei 
 in Hispánias dirécti sunt. Cum autem hi váriis úrbibus evangelizássent, 
 et innúmeras multitúdines ad Christi fidem perduxíssent, divérsis in ea 
 província locis quievérunt; scílicet Torquátus Acci, Ctésiphon Vérgii, 
 Secúndus Abulæ, Indalétius Urci, Cæcílius Illíberi, Hesychius Cartéjæ et 
 Euphrásius Illitúrgi.
\switchcolumn
\selectlanguage{english}
In Spain, the Saints Torquatus, 
 Ctesiphon, Secundus, Indaletius, Cecilius, Hesychius, and Euphrasius, who 
 were consecrated bishops at Rome by the holy apostles, and sent to Spain to 
 preach the word of God. When they had evangelized various cities, and 
 brought innumerable multitudes under the yoke of Christ, they rested in 
 peace in different places in that country: Torquatus at Cadiz, Ctesiphon at 
 Vierco, Secundus at Avila, Indaletius at Portilla, Cecílius at Elvira, 
 Hesychius at Gibraltar, and Euphrasius at Anduxar.
\switchcolumn*
\selectlanguage{latin}
Fausínæ, in Sardínia, sancti Simplícii, Epíscopi et Mártyris, qui, Diocletiáni témpore, sub 
 Bárbaro Prǽside, perfóssus láncea martyrium consummávit.
\switchcolumn
\selectlanguage{english}
At Fausina in Sardinia, in the time 
 of Diocletian and the governor Barbarus, Bishop St. Simplicius, who was 
 pierced with a lance and thus gained martyrdom.
\switchcolumn*
\selectlanguage{latin}
Eboræ, in Lusitánia, 
 sancti Máncii Mártyris.
\switchcolumn
\selectlanguage{english}
At Evora in Portugal, St. Mancius, 
 martyr.
\switchcolumn*
\selectlanguage{latin}
In ínsula Chio natális 
 beáti Isidóri Mártyris, in cujus Basílica exstat púteus, in quem fertur 
 fuísse injéctus, de cujus aqua infírmi potáti sæpius sanántur.
\switchcolumn
\selectlanguage{english}
In the island of Chio, the birthday 
 of blessed Isidore, martyr, in whose church is a well into which he is said 
 to have been thrown. By drinking of the water of this well, the sick 
 are frequently cured.
\switchcolumn*
\selectlanguage{latin}
Lampásci, in 
 Hellespónto, pássio sanctórum Petri, Andréæ, Pauli et Dionysiæ.
\switchcolumn
\selectlanguage{english}
At Lampascum in the Hellespont, the 
 martyrdom of the Saints Peter, Andrew, Paul, and Dionysia.
\switchcolumn*
\selectlanguage{latin}
Arvérnis, in Gállia, 
 sanctórum Mártyrum Cássii, Victoríni, Máximi et Sociórum.
\switchcolumn
\selectlanguage{english}
In the Auvergne in France, the holy martyrs 
 Cassius, Victorinus, Maximus, and their companions.
\switchcolumn*
\selectlanguage{latin}
Ghelæ, in Brabántia, 
 sanctæ Dympnæ, Vírginis et Mártyris, fíliæ Regis Hibernórum; quæ, cum in 
 fide Christi et virginitáte servánda permanéret immóbilis, a patre jussa est 
 decollári.
\switchcolumn
\selectlanguage{english}
At Gheel in Brabant, St. Dymphna, 
 virgin and martyr, daughter of the king of Ireland. By order of her 
 father, she was beheaded for the faith of Christ and the preservation of her 
 virginity.
\switchcolumn*
\selectlanguage{latin}
\end{paracol}


% ---- martyrology/mart05/mart0516.htm
\needspace{10\baselineskip}
\begin{paracol}{2}
\selectlanguage{latin}
\begin{center}{\color{gregoriocolor} Décimo séptimo Kaléndas Júnii. 
 Luna\dots\ }\end{center}
\switchcolumn
\selectlanguage{english}
\begin{center}{\color{gregoriocolor} The 16\textsuperscript{th} Day of May. 
 The\dots\ Day of the Moon.}\end{center}
\end{paracol}

\noindent\begin{tabularx}{\linewidth}{*{19}{>{\centering\arraybackslash}X}}
 \textcolor{gregoriocolor}{a} & \textcolor{gregoriocolor}{b} & \textcolor{gregoriocolor}{c} & \textcolor{gregoriocolor}{d} & \textcolor{gregoriocolor}{e} & \textcolor{gregoriocolor}{f} & \textcolor{gregoriocolor}{g} & \textcolor{gregoriocolor}{h} & \textcolor{gregoriocolor}{i} & \textcolor{gregoriocolor}{k} & \textcolor{gregoriocolor}{l} & \textcolor{gregoriocolor}{m} & \textcolor{gregoriocolor}{n} & \textcolor{gregoriocolor}{p} & \textcolor{gregoriocolor}{q} & \textcolor{gregoriocolor}{r} & \textcolor{gregoriocolor}{s} & \textcolor{gregoriocolor}{t} & \textcolor{gregoriocolor}{u} \\
 19 & 20 & 21 & 22 & 23 & 24 & 25 & 26 & 27 & 28 & 29 & 30 & 1 & 2 & 3 & 4 & 5 & 6 & 7 \\
\end{tabularx}
\vspace{0.5\baselineskip}
\noindent\begin{tabularx}{\linewidth}{*{12}{>{\centering\arraybackslash}X}}
 \textcolor{gregoriocolor}{A} & \textcolor{gregoriocolor}{B} & \textcolor{gregoriocolor}{C} & \textcolor{gregoriocolor}{D} & \textcolor{gregoriocolor}{E} & F & \textcolor{gregoriocolor}{F} & \textcolor{gregoriocolor}{G} & \textcolor{gregoriocolor}{H} & \textcolor{gregoriocolor}{M} & \textcolor{gregoriocolor}{N} & \textcolor{gregoriocolor}{P} \\
 8 & 9 & 10 & 11 & 12 & 13 & 13 & 14 & 15 & 16 & 17 & 18 \\
\end{tabularx}

\begin{paracol}{2}
\selectlanguage{latin}
\lettrine[lines=2]{E}{ugúbii} sancti Ubáldi, 
 Epíscopi et Confessóris, miráculis clari.
\switchcolumn
\selectlanguage{english}
\lettrine[lines=2]{A}{t} Gubbio, St. Ubaldus, bishop and 
 confessor renowned for his miracles.
\switchcolumn*
\selectlanguage{latin}
Antisiodóri pássio 
 sancti Peregríni, qui fuit primus ejúsdem civitátis Epíscopus. Hic, a 
 beáto Xysto Papa Secúndo in Gállias, una cum áliis Cléricis, missus, ibi, 
 Evangélicæ prædicatiónis múnere expléto, capitáli senténtia damnátus corónam 
 méruit sempitérnam.
\switchcolumn
\selectlanguage{english}
At Auxerre, the passion of St. 
 Peregrinus, first bishop of that city. He was sent into France with 
 other clerics by the blessed Pope Sixtus II, and having accomplished his 
 work of preaching the Gospel, he was condemned to capital punishment, and 
 merited for himself an everlasting crown.
\switchcolumn*
\selectlanguage{latin}
In Pérside sanctórum 
 Mártyrum Audæ Epíscopi, septem Presbyterórum, novem Diaconórum, et septem 
 Vírginum; qui sub Isdegérde Rege, variis tormentórum genéribus cruciáti, 
 gloriósum martyrium complevérunt.
\switchcolumn
\selectlanguage{english}
In Persia, the holy martyrs Audas, 
 a bishop, seven priests, nine deacons and seven virgins, who endured various 
 kins of torments under King Isdegerdes, and thus gloriously completed their 
 martyrdom.
\switchcolumn*
\selectlanguage{latin}
Pragæ, in Bohémia, sancti Joánnis Nepomucéni, Metropolitánæ Ecclésiæ Canónici; qui, frustra 
 tentátus ut sigílli sacramentális fidem próderet, martyrii palmam, in flumen 
 Moldávam dejéctus, eméruit.
\switchcolumn
\selectlanguage{english}
At Prague in Bohemia, St. John 
 Nepomucene, a canon of the cathedral church, who, being tempted in vain to 
 betray the secret of confession, was cast into the River Moldau, and thus 
 won the palm of martyrdom.
\switchcolumn*
\selectlanguage{latin}
In Isáuria natális 
 sanctórum Mártyrum Aquilíni et Victoriáni.
\switchcolumn
\selectlanguage{english}
In Isauria, the birthday of the 
 holy martyrs Aquilinus and Victorian.
\switchcolumn*
\selectlanguage{latin}
Uzali, in Africa, 
 sanctórum Mártyrum Felícis et Gennádii.
\switchcolumn
\selectlanguage{english}
At Uzalis in Africa, the holy 
 martyrs Felix and Gennadius.
\switchcolumn*
\selectlanguage{latin}
In Palæstína pássio 
 sanctórum Monachórum, a Saracénis in laura sancti Sabbæ interfectórum.
\switchcolumn
\selectlanguage{english}
In Palestine, the martyrdom of the 
 holy monks massacred by the Saracens in the monastery of St. Sabbas.
\switchcolumn*
\selectlanguage{latin}
Janóviæ, apud Pinsk, 
 in Polésia, sancti Andréæ Bobóla, Sacerdótis e Societáte Jesu, qui, a 
 schismáticis innumerabília tormentórum génera perpéssus, illústri martyrio 
 coronátus est.
\switchcolumn
\selectlanguage{english}
At Janow, near Pinsk in Lithuania, 
 St. Andrew Bobola, priest of the Society of Jesus, who having suffered many 
 kinds of torments at the hands of the schismatics, was crowned with an 
 illustrious martyrdom.
\switchcolumn*
\selectlanguage{latin}
Ambiáni, in Gállia, 
 sancti Honoráti Epíscopi.
\switchcolumn
\selectlanguage{english}
At Amiens in France, St. Honoratus, 
 bishop.
\switchcolumn*
\selectlanguage{latin}
Apud Cenómanos, in 
 Gállia, sancti Dómnoli Epíscopi.
\switchcolumn
\selectlanguage{english}
At Le Mans in France, St. Domnolus, 
 bishop.
\switchcolumn*
\selectlanguage{latin}
Mirándulæ, in Æmília, 
 sancti Possídii, Calaménsis in Numídia Epíscopi, qui fuit sancti Augustíni 
 discípulus, ejúsque præcláram vitam scripsit.
\switchcolumn
\selectlanguage{english}
At Mirandola in Aemilia, St. 
 Possidius, bishop of Calamae, and disciple of St. Augustine, of whose 
 glorious life he wrote a history.
\switchcolumn*
\selectlanguage{latin}
In monastério 
 Enachduinénsi, in Hibérnia, tránsitus sancti Brendáni Presbyteri et 
 Clonferténsis Abbátis.
\switchcolumn
\selectlanguage{english}
In the monastery of Enachduin in 
 Ireland, the death of St. Brendan, abbot of Clonfert.
\switchcolumn*
\selectlanguage{latin}
Trecis, in Gállia, 
 sancti Fídoli Confessóris.
\switchcolumn
\selectlanguage{english}
At Treves in France, St. Fidolus, 
 confessor.
\switchcolumn*
\selectlanguage{latin}
Apud Forum Júlii, in 
 Gállia, sanctæ Máximæ Vírginis, quæ, multis clara virtútibus, in pace 
 quiévit.
\switchcolumn
\selectlanguage{english}
At Fréjus in France, St. Maxima, 
 virgin, who died in peace with a reputation for many virtues.
\switchcolumn*
\selectlanguage{latin}
\end{paracol}


% ---- martyrology/mart05/mart0517.htm
\needspace{10\baselineskip}
\begin{paracol}{2}
\selectlanguage{latin}
\begin{center}{\color{gregoriocolor} Sextodécimo Kaléndas Júnii. 
 Luna\dots\ }\end{center}
\switchcolumn
\selectlanguage{english}
\begin{center}{\color{gregoriocolor} The 
 17\textsuperscript{th} Day of May. The\dots\ Day of the Moon.}\end{center}
\end{paracol}

\noindent\begin{tabularx}{\linewidth}{*{19}{>{\centering\arraybackslash}X}}
 \textcolor{gregoriocolor}{a} & \textcolor{gregoriocolor}{b} & \textcolor{gregoriocolor}{c} & \textcolor{gregoriocolor}{d} & \textcolor{gregoriocolor}{e} & \textcolor{gregoriocolor}{f} & \textcolor{gregoriocolor}{g} & \textcolor{gregoriocolor}{h} & \textcolor{gregoriocolor}{i} & \textcolor{gregoriocolor}{k} & \textcolor{gregoriocolor}{l} & \textcolor{gregoriocolor}{m} & \textcolor{gregoriocolor}{n} & \textcolor{gregoriocolor}{p} & \textcolor{gregoriocolor}{q} & \textcolor{gregoriocolor}{r} & \textcolor{gregoriocolor}{s} & \textcolor{gregoriocolor}{t} & \textcolor{gregoriocolor}{u} \\
 20 & 21 & 22 & 23 & 24 & 25 & 26 & 27 & 28 & 29 & 30 & 1 & 2 & 3 & 4 & 5 & 6 & 7 & 8 \\
\end{tabularx}
\vspace{0.5\baselineskip}
\noindent\begin{tabularx}{\linewidth}{*{12}{>{\centering\arraybackslash}X}}
 \textcolor{gregoriocolor}{A} & \textcolor{gregoriocolor}{B} & \textcolor{gregoriocolor}{C} & \textcolor{gregoriocolor}{D} & \textcolor{gregoriocolor}{E} & F & \textcolor{gregoriocolor}{F} & \textcolor{gregoriocolor}{G} & \textcolor{gregoriocolor}{H} & \textcolor{gregoriocolor}{M} & \textcolor{gregoriocolor}{N} & \textcolor{gregoriocolor}{P} \\
 8 & 10 & 11 & 12 & 13 & 14 & 14 & 15 & 16 & 17 & 18 & 19 \\
\end{tabularx}

\begin{paracol}{2}
\selectlanguage{latin}
\lettrine[lines=2]{A}{pud} Villam Regálem, 
 in Hispánia, sancti Paschális, ex Ordine Minórum, Confessóris, miræ 
 innocéntiæ et pæniténtiæ viri; quem Leo Papa Décimus tértius cæléstem 
 eucharisticórum Cœtuum et Societátum a sanctíssima Eucharístia Patrónum 
 declarávit.
\switchcolumn
\selectlanguage{english}
\lettrine[lines=2]{A}{t} Villareal in Spain, St. Paschal 
 of the Order of Friars Minor, confessor. He was a man remarkable for 
 innocence of life and the spirit of penance, whom Pope Leo XIII declared to 
 be the heavenly patron of Eucharistic Congresses and of societies formed to 
 honour the Most Blessed Sacrament.
\switchcolumn*
\selectlanguage{latin}
Noviodúni, in Gálliis, 
 sanctórum Mártyrum Herádii, Pauli et Aquilíni, cum duóbus áliis.
\switchcolumn
\selectlanguage{english}
At Noyon in France, the holy 
 martyrs Heradius, Paul, and Aquilinus, with two others.
\switchcolumn*
\selectlanguage{latin}
Chalcédone sanctórum 
 Mártyrum Solochónis et Sociórum mílitum, sub Maximiáno Imperatóre.
\switchcolumn
\selectlanguage{english}
At Chalcedon, the holy martyrs 
 Solochan and his companions.
\switchcolumn*
\selectlanguage{latin}
Alexandríæ sanctórum 
 Mártyrum Adriónis, Victóris et Basíllæ.
\switchcolumn
\selectlanguage{english}
At Alexandria, the holy martyrs 
 Adrion, Victor, and Basilla.
\switchcolumn*
\selectlanguage{latin}
Eódem die sanctæ 
 Restitútæ, Vírginis et Mártyris; quæ, in Africa, Valeriáno imperánte, a 
 Próculo Júdice várie torta, et in navículam pice et stupa refértam, ut in mari comburerétur, impósita, tandem, cum in incensóres, immísso igne, flamma 
 converterétur, in oratióne spíritum Deo réddidit. Ipsíus corpus cum 
 eádem navícula, Dei nutu, ad Ænáriam ínsulam, prope Neápolim, in Campánia, 
 devéctum est, et a Christiánis magna veneratióne suscéptum; ac póstmodum in 
 ejus honórem Constantínus Magnus Basílicam in ipsa urbe Neápoli erigéndam 
 curávit.
\switchcolumn
\selectlanguage{english}
Also St. Restituta, virgin and 
 martyr, who was subjected to various kinds of tortures in Africa by the 
 judge Proculus, in the reign of Valerian, and then put in a boat filled with 
 pitch and oakum, to be burned to death on the sea. But the flame 
 turned on those who kindled it, and the saint yielded her soul to God in 
 prayer. Her body was, by Divine Providence, carried in the boat to the 
 island of Ischia, near Naples, where it was received by the Christians with 
 great veneration. A church was afterwards erected in her honour at 
 Naples by Constantine the Great.
\switchcolumn*
\selectlanguage{latin}
\end{paracol}


% ---- martyrology/mart05/mart0518.htm
\needspace{10\baselineskip}
\begin{paracol}{2}
\selectlanguage{latin}
\begin{center}{\color{gregoriocolor} Quintodécimo Kaléndas Júnii. 
 Luna\dots\ }\end{center}
\switchcolumn
\selectlanguage{english}
\begin{center}{\color{gregoriocolor} The 
 18\textsuperscript{th} Day of May. The\dots\ Day of the Moon.}\end{center}
\end{paracol}

\noindent\begin{tabularx}{\linewidth}{*{19}{>{\centering\arraybackslash}X}}
 \textcolor{gregoriocolor}{a} & \textcolor{gregoriocolor}{b} & \textcolor{gregoriocolor}{c} & \textcolor{gregoriocolor}{d} & \textcolor{gregoriocolor}{e} & \textcolor{gregoriocolor}{f} & \textcolor{gregoriocolor}{g} & \textcolor{gregoriocolor}{h} & \textcolor{gregoriocolor}{i} & \textcolor{gregoriocolor}{k} & \textcolor{gregoriocolor}{l} & \textcolor{gregoriocolor}{m} & \textcolor{gregoriocolor}{n} & \textcolor{gregoriocolor}{p} & \textcolor{gregoriocolor}{q} & \textcolor{gregoriocolor}{r} & \textcolor{gregoriocolor}{s} & \textcolor{gregoriocolor}{t} & \textcolor{gregoriocolor}{u} \\
 21 & 22 & 23 & 24 & 25 & 26 & 27 & 28 & 29 & 30 & 1 & 2 & 3 & 4 & 5 & 6 & 7 & 8 & 9 \\
\end{tabularx}
\vspace{0.5\baselineskip}
\noindent\begin{tabularx}{\linewidth}{*{12}{>{\centering\arraybackslash}X}}
 \textcolor{gregoriocolor}{A} & \textcolor{gregoriocolor}{B} & \textcolor{gregoriocolor}{C} & \textcolor{gregoriocolor}{D} & \textcolor{gregoriocolor}{E} & F & \textcolor{gregoriocolor}{F} & \textcolor{gregoriocolor}{G} & \textcolor{gregoriocolor}{H} & \textcolor{gregoriocolor}{M} & \textcolor{gregoriocolor}{N} & \textcolor{gregoriocolor}{P} \\
 10 & 11 & 12 & 13 & 14 & 15 & 15 & 16 & 17 & 18 & 19 & 20 \\
\end{tabularx}

\begin{paracol}{2}
\selectlanguage{latin}
\lettrine[lines=2]{C}{ameríni} sancti 
 Venántii Mártyris, qui, annos quíndecim natus, sub Décio Imperatóre et 
 Antíocho Præside, una cum áliis decem, gloriósi certáminis cursum, 
 cervícibus abscíssis implévit.
\switchcolumn
\selectlanguage{english}
\lettrine[lines=2]{A}{t} Camerino, the holy martyr Venantius, who, at fifteen years of age, along with ten others, ended a 
 glorious ordeal by being beheaded under Emperor Decius and the governor 
 Antiochus.
\switchcolumn*
\selectlanguage{latin}
Ravénnæ, natális 
 sancti Joánnis Primi, Papæ et Mártyris; qui, ab Ariáno Itáliæ Rege 
 Theodoríco illuc dolo evocátus, ibídem, ab eo propter orthodóxam fidem diu 
 macerátus in cárcere, vitam finívit. Ipsíus autem festum recólitur 
 sexto Kaléndas Júnii, quo die sacrum ejus corpus, Romam relátum, in Basílica sancti Petri, Apostolórum Príncipis, sepúltum est.
\switchcolumn
\selectlanguage{english}
The birthday of St. John I, pope 
 and martyr, who was called to Ravenna by the Arian king of Italy, Theodoric, 
 and died there after being in prison a long time for the true faith. 
 His feast, however, is celebrated on the 27th of May, the day on which his 
 revered body was taken to Rome and buried in the basilica of St. Peter, 
 prince of the apostles.
\switchcolumn*
\selectlanguage{latin}
Spoléti sancti 
 Felícis Epíscopi, qui, sub Maximiáno Imperatóre, martyrii palmam adéptus 
 est.
\switchcolumn
\selectlanguage{english}
At Spoleto, St. Felix, a bishop who 
 obtained the palm of martyrdom under Emperor Maximian.
\switchcolumn*
\selectlanguage{latin}
Heracléæ, in Ægypto, 
 sancti Potamónis Epíscopi, qui primum, sub Maximiáno Galério, Confessor fuit; 
 deínde, sub Constántio Imperatóre et Ariáno Prǽside Philágrio, martyrio 
 coronátus est. Ipsum vero beátum virum sancti Ecclésiæ Patres 
 Athanásius et Epiphánius suis láudibus celebrárunt.
\switchcolumn
\selectlanguage{english}
At Heraclea in Egypt, Bishop St. 
 Potamon, first a confessor under Maximian Galerius, and afterwards, a martyr 
 under Emperor Constantius, and the Arian governor Philagrius. 
 Athanasius and Epiphanius, Fathers of the Church, have sung the praises of 
 this holy man.
\switchcolumn*
\selectlanguage{latin}
In Ægypto sancti 
 Dióscori Lectóris, in quem Præses multa et vária torménta ita exércuit, ut 
 ungues ejus effóderet et lampádibus látera inflammáret, sed, cæléstis 
 lúminis fulgóre pertérriti, cecidérunt minístri; novíssime autem ipse 
 Dióscorus, láminis ardéntibus adústus, martyrium consummávit.
\switchcolumn
\selectlanguage{english}
In Egypt, St. Dioscorus, a lector, 
 who was subjected by the governor to many and diverse torments, such as the 
 tearing off of his nails and the burning of his sides with torches; but a 
 light from heaven having prostrated the executioners, the saint's martyrdom 
 was finally ended by having red-hot metal plates applied to his body.
\switchcolumn*
\selectlanguage{latin}
Ancyræ, in Galátia, sancti Theódoti Mártyris, et sanctárum septem Vírginum et Mártyrum, scílicet 
 Thecúsæ, quæ ipsíus Theódoti erat 
 ámita, Alexándræ, Cláudiæ, Phaínæ, 
 Euphrásiæ, Matrónæ et Julíttæ. Hæ a Prǽside primum prostitútæ, sed Dei 
 virtúte servátæ, deínde, lapídibus ad colla ligátis, in palúdem mersæ sunt; quarum relíquias colléctas cum Theódotus honorífice sepelísset, 
 ídeo, 
 comprehénsus a Prǽside, sævíssime dilaniátus est, ac demum, ense percússus, 
 martyrii corónam accépit.
\switchcolumn
\selectlanguage{english}
At Ancyra in Galatia, the martyr 
 St. Theodotus, and the holy virgins Thecusa, his aunt, Alexandra, Claudia, 
 Faina, Euphrasia, Matrona, and Julitta. They were at first taken to a 
 place of debauchery, but the power of God prevented them from evil, and they 
 later had stones fastened to their necks and were plunged into a lake. 
 For gathering the remains and burying them honorably, Theodotus was arrested 
 by the governor, and after having been horribly lacerated, was put to the 
 sword, and thus received the crown of martyrdom.
\switchcolumn*
\selectlanguage{latin}
Upsali, in Suécia, 
 sancti Eríci, Regis et Mártyris.
\switchcolumn
\selectlanguage{english}
At Upsal in Sweden, St. Erik, king 
 and martyr.
\switchcolumn*
\selectlanguage{latin}
Romæ sancti Felícis 
 Confessóris, ex Ordine Minórum Capuccinórum, evangélica simplicitáte et 
 caritáte conspícui; quem Clemens Undécimus, Póntifex Máximus, Sanctórum 
 fastis adscrípsit.
\switchcolumn
\selectlanguage{english}
At Rome, St. Felix, confessor of 
 the Order of Friars Minor Capuchin, celebrated for his evangelical 
 simplicity and charity. He was inscribed on the roll of the saints by 
 the Sovereign Pontiff Clement XI.
\switchcolumn*
\selectlanguage{latin}
\end{paracol}


% ---- martyrology/mart05/mart0519.htm
\needspace{10\baselineskip}
\begin{paracol}{2}
\selectlanguage{latin}
\begin{center}{\color{gregoriocolor} Quartodécimo Kaléndas Júnii. 
 Luna\dots\ }\end{center}
\switchcolumn
\selectlanguage{english}
\begin{center}{\color{gregoriocolor} The 
 19\textsuperscript{th} Day of May. The\dots\ Day of the Moon.}\end{center}
\end{paracol}

\noindent\begin{tabularx}{\linewidth}{*{19}{>{\centering\arraybackslash}X}}
 \textcolor{gregoriocolor}{a} & \textcolor{gregoriocolor}{b} & \textcolor{gregoriocolor}{c} & \textcolor{gregoriocolor}{d} & \textcolor{gregoriocolor}{e} & \textcolor{gregoriocolor}{f} & \textcolor{gregoriocolor}{g} & \textcolor{gregoriocolor}{h} & \textcolor{gregoriocolor}{i} & \textcolor{gregoriocolor}{k} & \textcolor{gregoriocolor}{l} & \textcolor{gregoriocolor}{m} & \textcolor{gregoriocolor}{n} & \textcolor{gregoriocolor}{p} & \textcolor{gregoriocolor}{q} & \textcolor{gregoriocolor}{r} & \textcolor{gregoriocolor}{s} & \textcolor{gregoriocolor}{t} & \textcolor{gregoriocolor}{u} \\
 22 & 23 & 24 & 25 & 26 & 27 & 28 & 29 & 30 & 1 & 2 & 3 & 4 & 5 & 6 & 7 & 8 & 9 & 10 \\
\end{tabularx}
\vspace{0.5\baselineskip}
\noindent\begin{tabularx}{\linewidth}{*{12}{>{\centering\arraybackslash}X}}
 \textcolor{gregoriocolor}{A} & \textcolor{gregoriocolor}{B} & \textcolor{gregoriocolor}{C} & \textcolor{gregoriocolor}{D} & \textcolor{gregoriocolor}{E} & F & \textcolor{gregoriocolor}{F} & \textcolor{gregoriocolor}{G} & \textcolor{gregoriocolor}{H} & \textcolor{gregoriocolor}{M} & \textcolor{gregoriocolor}{N} & \textcolor{gregoriocolor}{P} \\
 11 & 12 & 13 & 14 & 15 & 16 & 16 & 17 & 18 & 19 & 20 & 21 \\
\end{tabularx}

\begin{paracol}{2}
\selectlanguage{latin}
\lettrine[lines=2]{N}{atalis} sancti Petri 
 de Moróno, Confessóris, qui ex Anachoréta Summus Póntifex creátus, dictus 
 est Cælestínus Quintus. Sed Pontificátu se póstmodum abdicávit, et in 
 solitúdine religiósam vitam agens, virtútibus et miráculis clarus, migrávit 
 ad Dóminum.
\switchcolumn
\selectlanguage{english}
\lettrine[lines=2]{T}{he} birthday of St. Peter of Moroni 
 who, while leading the life of an anchoret, was created Sovereign Pontiff 
 and called Celestine V. He later abdicated the pontificate, and led a 
 religious life in solitude, where, renowned for virtues and miracles, he 
 went to the Lord.
\switchcolumn*
\selectlanguage{latin}
Romæ sanctæ Pudentiánæ 
 Vírginis, quæ, post innúmeros agónes, post multórum Mártyrum venerabíliter 
 exhíbitas sepultúras, post omnes facultátes suas pro Christo paupéribus 
 erogátas, tandem de terris ad cælos migrávit.
\switchcolumn
\selectlanguage{english}
At Rome, the saintly virgin 
 Pudentiana, who, after numberless tribulations, after burying with respect 
 many martyrs, and distributing all her goods to the poor for Christ's sake, 
 departed from this world to go to heaven.
\switchcolumn*
\selectlanguage{latin}
Ibídem sancti Pudéntis 
 Senatóris, qui ejúsdem sanctæ Pudentiánæ Vírginis ac sanctæ item Praxédis 
 Vírginis fuit pater; atque, ab Apostólicis Christo in baptísmo vestítus, 
 innocéntem túnicam usque ad vitæ corónam immaculáte custodívit.
\switchcolumn
\selectlanguage{english}
In the same city, St. Pudens, 
 senator, father of the virgins Pudentiana and Praxedes. He was clothed 
 with Christ in baptism by the apostles, and preserved the robe of innocence 
 unspotted until he received the crown of life.
\switchcolumn*
\selectlanguage{latin}
Item Romæ, via Appia, 
 natális sanctórum Calóceri et Parthénii eunuchórum; quorum prior, cum esset 
 præpósitus cubículi uxóris Décii Imperatóris, postérior vero altérius 
 múneris primicérius, et ambo nollent idólis sacrificáre, idcírco, ejúsdem 
 Imperatóris jussu váriis ac diris sunt excruciáti supplíciis, ac tandem, 
 cervícibus ardénti fuste contúsis, spíritum Deo tradidérunt.
\switchcolumn
\selectlanguage{english}
Also at Rome, on the Appian Way, 
 the birthday of the Saints Calocerus and Parthenius, eunuchs. The 
 former was chamberlain of the wife of Emperor Decius, and the latter chief 
 officer in another department. Because they refused to offer sacrifice 
 to idols they were tortured in many cruel ways, and finally when their necks 
 were broken with cudgels, they gave up their souls to God.
\switchcolumn*
\selectlanguage{latin}
Nicomedíæ sancti 
 Philóteri Mártyris, qui Paciáni Procónsulis fílius 
 éxstitit, et, sub 
 Diocletiáno Imperatóre, multa passus, martyrii corónam accépit.
\switchcolumn
\selectlanguage{english}
At Nicomedia, the martyr St. 
 Philoterus, son of the proconsul Pacian, who after suffering much under 
 Emperor Diocletian, received the crown of martyrdom.
\switchcolumn*
\selectlanguage{latin}
Ibídem sanctárum sex 
 Vírginum et Mártyrum, quarum præcípua erat Cyríaca, quæ, cum líbere 
 Maximiánum impietátis objurgásset, diríssime cæsa et dilaniáta est, atque ad 
 últimum, igne cremáta, martyrium consummávit.
\switchcolumn
\selectlanguage{english}
In the same city, six holy virgins 
 and martyrs. The principal one, named Cyriaca, having boldly reproved 
 Maximian for his impiety, was severely scourged and lacerated, and then 
 consumed by fire.
\switchcolumn*
\selectlanguage{latin}
Cantuáriæ, in Anglia, 
 sancti Dunstáni Epíscopi.
\switchcolumn
\selectlanguage{english}
At Canterbury in England, St. 
 Dunstan, bishop.
\switchcolumn*
\selectlanguage{latin}
Lohanéti, in Británnia minóre, sancti Ivónis Presbyteri et Confessóris; qui, pro Christi amóre, 
 causas pupillórum, viduárum ac páuperum defendébat.
\switchcolumn
\selectlanguage{english}
In Brittany, St. Ivo, priest and 
 confessor, who for the love of Christ, defended the interests of orphans, 
 widows and the poor.
\switchcolumn*
\selectlanguage{latin}
Ficécli, in Etrúria, 
 sancti Theóphili a Curte, Confessóris, Sacerdótis Ordinis Fratrum Minórum, 
 sacrórum recéssum propagatóris, quem Pius Papa Undécimus inter Sanctos 
 rétulit.
\switchcolumn
\selectlanguage{english}
At Fucecchio in Etruria, St. 
 Theophilus of Curte, confessor and priest of the Order of Friars Minor, who 
 was canonized by Pope Pius XI.
\switchcolumn*
\selectlanguage{latin}
\end{paracol}


% ---- martyrology/mart05/mart0520.htm
\needspace{10\baselineskip}
\begin{paracol}{2}
\selectlanguage{latin}
\begin{center}{\color{gregoriocolor} Tertiodécimo Kaléndas Júnii. 
 Luna\dots\ }\end{center}
\switchcolumn
\selectlanguage{english}
\begin{center}{\color{gregoriocolor} The 
 20\textsuperscript{th} Day of May. The\dots\ Day of the Moon.}\end{center}
\end{paracol}

\noindent\begin{tabularx}{\linewidth}{*{19}{>{\centering\arraybackslash}X}}
 \textcolor{gregoriocolor}{a} & \textcolor{gregoriocolor}{b} & \textcolor{gregoriocolor}{c} & \textcolor{gregoriocolor}{d} & \textcolor{gregoriocolor}{e} & \textcolor{gregoriocolor}{f} & \textcolor{gregoriocolor}{g} & \textcolor{gregoriocolor}{h} & \textcolor{gregoriocolor}{i} & \textcolor{gregoriocolor}{k} & \textcolor{gregoriocolor}{l} & \textcolor{gregoriocolor}{m} & \textcolor{gregoriocolor}{n} & \textcolor{gregoriocolor}{p} & \textcolor{gregoriocolor}{q} & \textcolor{gregoriocolor}{r} & \textcolor{gregoriocolor}{s} & \textcolor{gregoriocolor}{t} & \textcolor{gregoriocolor}{u} \\
 23 & 24 & 25 & 26 & 27 & 28 & 29 & 30 & 1 & 2 & 3 & 4 & 5 & 6 & 7 & 8 & 9 & 10 & 11 \\
\end{tabularx}
\vspace{0.5\baselineskip}
\noindent\begin{tabularx}{\linewidth}{*{12}{>{\centering\arraybackslash}X}}
 \textcolor{gregoriocolor}{A} & \textcolor{gregoriocolor}{B} & \textcolor{gregoriocolor}{C} & \textcolor{gregoriocolor}{D} & \textcolor{gregoriocolor}{E} & F & \textcolor{gregoriocolor}{F} & \textcolor{gregoriocolor}{G} & \textcolor{gregoriocolor}{H} & \textcolor{gregoriocolor}{M} & \textcolor{gregoriocolor}{N} & \textcolor{gregoriocolor}{P} \\
 12 & 13 & 14 & 15 & 16 & 17 & 17 & 18 & 19 & 20 & 21 & 22 \\
\end{tabularx}

\begin{paracol}{2}
\selectlanguage{latin}
\lettrine[lines=2]{A}{quilæ,} in Vestínis, 
 sancti Bernardíni Senénsis, Sacerdótis ex Ordine Minórum et Confessóris, qui 
 verbo et exémplo Itáliam illustrávit.
\switchcolumn
\selectlanguage{english}
\lettrine[lines=2]{A}{t} Aquila in Abruzzi, St. Bernardin 
 of Siena, priest of the Order of Friars Minor, who added to the glory of 
 Italy by his preaching and his example.
\switchcolumn*
\selectlanguage{latin}
Romæ sanctæ Plautíllæ, 
 féminæ Consuláris, quæ beatórum Mártyrum Flávii Cleméntis Cónsulis soror et 
 Fláviæ Domitíllæ Vírginis mater 
 éxstitit; atque, a sancto Petro Apóstolo 
 baptizáta, ómnium virtútum laude refúlgens, quiévit in pace.
\switchcolumn
\selectlanguage{english}
At Rome, St. Plautilla, wife of a 
 consul, sister of the consul Flavius Clemens, and mother of the holy virgin 
 Flavia Domitilla, both martyrs. She was baptized by the apostle St. 
 Peter, and after giving an example of all the virtues, she rested in peace.
\switchcolumn*
\selectlanguage{latin}
Item Romæ, via Salária, 
 natális sanctæ Basíllæ Vírginis, quæ, cum esset ex régio génere, et 
 illustríssimum sponsum habéret illúmque dimisísset, accusáta fuit ab eo quod 
 esset Christiána, et mox decrétum est a Galliéno Augústo ut aut sponsum 
 recíperet, aut gládio interíret; cumque ipsa Virgo, convénta de hoc, 
 respondísset se Regem regum habére sponsum, gládio transverberáta est.
\switchcolumn
\selectlanguage{english}
Also at Rome, on the Salarian Way, 
 the birthday of St. Basilla, virgin, who was of a royal family and betrothed 
 to a nobleman. When she refused to marry him, he accused her of being 
 a Christian. Emperor Gallienus gave orders that she should accept the 
 person to whom she had been engaged, or die by the sword. Being 
 informed of this, and answering that she had for her spouse the King of 
 kings, she was pierced with a sword.
\switchcolumn*
\selectlanguage{latin}
Nemáusi, in Gálliis, 
 sancti Baudélii Mártyris, qui comprehénsus est a Pagánis, et cum sacrificáre 
 nollet idólis et in Christi fide inter vérbera et torménta immóbilis 
 persísteret, martyrii palmam pretiósa morte suscépit.
\switchcolumn
\selectlanguage{english}
At Nimes in France, St. Baudelius, 
 martyr. Being arrested, but refusing to sacrifice to idols, and 
 remaining immovable in the faith of Christ, despite blows and tortures, he 
 gained the palm of martyrdom by his praiseworthy death.
\switchcolumn*
\selectlanguage{latin}
Edéssæ, apud Ægas, in 
 Cilícia, sanctórum Mártyrum Thalelǽi, Astérii, Alexándri et Sociórum, qui 
 sub Numeriáno Imperatóre passi sunt.
\switchcolumn
\selectlanguage{english}
At Edessa near Aegea in Cilicia, 
 the holy martyrs Thalalaeus, Asterius, Alexander, and their companions, who 
 suffered under Emperor Numerian.
\switchcolumn*
\selectlanguage{latin}
In Thebáide sancti 
 Aquilæ Mártyris, qui pectínibus pro Christo dilaniátus fuit.
\switchcolumn
\selectlanguage{english}
In Thebais, St. Aquila, martyr to 
 the faith, whose body was torn with iron combs.
\switchcolumn*
\selectlanguage{latin}
Apud Bitúrcias, in 
 Aquitánia, sancti Austregísili, Epíscopi et Confessóris.
\switchcolumn
\selectlanguage{english}
At Bourges in France, St. 
 Austregisil, bishop and confessor.
\switchcolumn*
\selectlanguage{latin}
Bríxiæ sancti 
 Anastásii Epíscopi.
\switchcolumn
\selectlanguage{english}
At Brescia, St. Anastasius, bishop.
\switchcolumn*
\selectlanguage{latin}
Papíæ sancti Theodóri 
 Epíscopi.
\switchcolumn
\selectlanguage{english}
At Pavia, St. Theodore, bishop.
\switchcolumn*
\selectlanguage{latin}
\end{paracol}


% ---- martyrology/mart05/mart0521.htm
\needspace{10\baselineskip}
\begin{paracol}{2}
\selectlanguage{latin}
\begin{center}{\color{gregoriocolor} Duodécimo Kaléndas Júnii. 
 Luna\dots\ }\end{center}
\switchcolumn
\selectlanguage{english}
\begin{center}{\color{gregoriocolor} The 
 21\textsuperscript{st} Day of May. The\dots\ Day of the Moon.}\end{center}
\end{paracol}

\noindent\begin{tabularx}{\linewidth}{*{19}{>{\centering\arraybackslash}X}}
 \textcolor{gregoriocolor}{a} & \textcolor{gregoriocolor}{b} & \textcolor{gregoriocolor}{c} & \textcolor{gregoriocolor}{d} & \textcolor{gregoriocolor}{e} & \textcolor{gregoriocolor}{f} & \textcolor{gregoriocolor}{g} & \textcolor{gregoriocolor}{h} & \textcolor{gregoriocolor}{i} & \textcolor{gregoriocolor}{k} & \textcolor{gregoriocolor}{l} & \textcolor{gregoriocolor}{m} & \textcolor{gregoriocolor}{n} & \textcolor{gregoriocolor}{p} & \textcolor{gregoriocolor}{q} & \textcolor{gregoriocolor}{r} & \textcolor{gregoriocolor}{s} & \textcolor{gregoriocolor}{t} & \textcolor{gregoriocolor}{u} \\
 24 & 25 & 26 & 27 & 28 & 29 & 30 & 1 & 2 & 3 & 4 & 5 & 6 & 7 & 8 & 9 & 10 & 11 & 12 \\
\end{tabularx}
\vspace{0.5\baselineskip}
\noindent\begin{tabularx}{\linewidth}{*{12}{>{\centering\arraybackslash}X}}
 \textcolor{gregoriocolor}{A} & \textcolor{gregoriocolor}{B} & \textcolor{gregoriocolor}{C} & \textcolor{gregoriocolor}{D} & \textcolor{gregoriocolor}{E} & F & \textcolor{gregoriocolor}{F} & \textcolor{gregoriocolor}{G} & \textcolor{gregoriocolor}{H} & \textcolor{gregoriocolor}{M} & \textcolor{gregoriocolor}{N} & \textcolor{gregoriocolor}{P} \\
 13 & 14 & 15 & 16 & 17 & 18 & 18 & 19 & 20 & 21 & 22 & 23 \\
\end{tabularx}

\begin{paracol}{2}
\selectlanguage{latin}
\lettrine[lines=2]{S}{ancti} Valéntis 
 Epíscopi, qui, una cum tribus púeris, necátus est.
\switchcolumn
\selectlanguage{english}
\lettrine[lines=2]{S}{t.} Valens, bishop, who was put to 
 death along with three children.
\switchcolumn*
\selectlanguage{latin}
Alexandríæ 
 commemorátio sanctórum Mártyrum Secúndi Presbyteri, et aliórum; quos sacris 
 diébus Pentecóstes, sub Constántio Imperatóre, Ariánus Epíscopus Geórgius 
 sævíssime occídi præcépit.
\switchcolumn
\selectlanguage{english}
At Alexandria, the commemoration of 
 the holy martyrs Secundus, a priest, and others, whom the Arian bishop 
 George ordered to be barbarously slain during the holy days of Pentecost, 
 under Emperor Constantius.
\switchcolumn*
\selectlanguage{latin}
In Mauritánia 
 Cæsariénsi natális sanctórum Mártyrum Diaconórum Timóthei, Pólii et Eutychii, 
 qui, in eádem regióne dísseminántes verbum Dei, páriter coronári meruérunt.
\switchcolumn
\selectlanguage{english}
In Morocco, the birthday of the 
 holy martyrs Timothy, Polius, and Eutychius, deacons, who merited to be 
 crowned together for spreading the word of God in that region.
\switchcolumn*
\selectlanguage{latin}
Cæsaréæ, in Cappadócia, 
 item natális sanctórum Mártyrum Polyéucti, Victórii et Donáti.
\switchcolumn
\selectlanguage{english}
At Caesarea in Cappadocia, the 
 birthday of the holy martyrs Polyeuctus, Victorinus, and Donatus.
\switchcolumn*
\selectlanguage{latin}
Córdubæ, in Hispánia, 
 sancti Secundíni Mártyris.
\switchcolumn
\selectlanguage{english}
At Cordova, the martyr St. 
 Secundinus.
\switchcolumn*
\selectlanguage{latin}
Eódem die sanctórum 
 Mártyrum Synésii et Theopómpi.
\switchcolumn
\selectlanguage{english}
The same day, the holy martyrs 
 Synesius and Theopompus.
\switchcolumn*
\selectlanguage{latin}
Cæsaréæ Philíppi 
 natális sanctórum Mártyrum Nicóstrati et Antíochi Tribunórum, cum áliis 
 milítibus.
\switchcolumn
\selectlanguage{english}
At Caesarea Philippi, the holy 
 martyrs Nicostratus and Antiochus, tribunes, with other soldiers.
\switchcolumn*
\selectlanguage{latin}
Alexandríæ 
 commemorátio sanctórum Episcopórum et Presbyterórum; qui, ab Ariánis exsílio 
 relegáti, sanctis Confessóribus sociári meruérunt.
\switchcolumn
\selectlanguage{english}
At Alexandria, the commemoration of 
 the saintly bishops and priests, who were banished by the Arians, and 
 merited to be numbered among the holy confessors.
\switchcolumn*
\selectlanguage{latin}
Níciæ, apud Varum 
 flúvium, sancti Hospítii Confessóris, abstinéntiæ virtúte ac prophetíæ 
 spíritu insígnis.
\switchcolumn
\selectlanguage{english}
At Nice in France, St. Hospitius, 
 confessor, distinguished by the virtue of abstinence and the spirit of 
 prophecy.
\switchcolumn*
\selectlanguage{latin}
\end{paracol}


% ---- martyrology/mart05/mart0522.htm
\needspace{10\baselineskip}
\begin{paracol}{2}
\selectlanguage{latin}
\begin{center}{\color{gregoriocolor} Undécimo Kaléndas Júnii. 
 Luna\dots\ }\end{center}
\switchcolumn
\selectlanguage{english}
\begin{center}{\color{gregoriocolor} The 
 22\textsuperscript{nd} Day of May. The\dots\ Day of the Moon.}\end{center}
\end{paracol}

\noindent\begin{tabularx}{\linewidth}{*{19}{>{\centering\arraybackslash}X}}
 \textcolor{gregoriocolor}{a} & \textcolor{gregoriocolor}{b} & \textcolor{gregoriocolor}{c} & \textcolor{gregoriocolor}{d} & \textcolor{gregoriocolor}{e} & \textcolor{gregoriocolor}{f} & \textcolor{gregoriocolor}{g} & \textcolor{gregoriocolor}{h} & \textcolor{gregoriocolor}{i} & \textcolor{gregoriocolor}{k} & \textcolor{gregoriocolor}{l} & \textcolor{gregoriocolor}{m} & \textcolor{gregoriocolor}{n} & \textcolor{gregoriocolor}{p} & \textcolor{gregoriocolor}{q} & \textcolor{gregoriocolor}{r} & \textcolor{gregoriocolor}{s} & \textcolor{gregoriocolor}{t} & \textcolor{gregoriocolor}{u} \\
 25 & 26 & 27 & 28 & 29 & 30 & 30 & 1 & 2 & 3 & 4 & 5 & 6 & 7 & 8 & 10 & 11 & 12 & 13 \\
\end{tabularx}
\vspace{0.5\baselineskip}
\noindent\begin{tabularx}{\linewidth}{*{12}{>{\centering\arraybackslash}X}}
 \textcolor{gregoriocolor}{A} & \textcolor{gregoriocolor}{B} & \textcolor{gregoriocolor}{C} & \textcolor{gregoriocolor}{D} & \textcolor{gregoriocolor}{E} & F & \textcolor{gregoriocolor}{F} & \textcolor{gregoriocolor}{G} & \textcolor{gregoriocolor}{H} & \textcolor{gregoriocolor}{M} & \textcolor{gregoriocolor}{N} & \textcolor{gregoriocolor}{P} \\
 14 & 15 & 16 & 17 & 18 & 19 & 19 & 20 & 21 & 22 & 23 & 24 \\
\end{tabularx}

\begin{paracol}{2}
\selectlanguage{latin}
\lettrine[lines=2]{R}{omæ} sanctórum 
 Mártyrum Faustíni, Timóthei et Venústi.
\switchcolumn
\selectlanguage{english}
\lettrine[lines=2]{A}{t} Rome, the holy martyrs Faustinus, 
 Timothy, and Venustus.
\switchcolumn*
\selectlanguage{latin}
In Africa sanctórum 
 Mártyrum Casti et Æmílii, qui per passiónis ignem martyrium consummárunt. 
 Hos (ut beátus Cypriánus scribit), in prima congressióne devíctos, Dóminus 
 victóres in secúndo prælio réddidit, ut fortióres ígnibus fíerent qui 
 ígnibus ante cessíssent.
\switchcolumn
\selectlanguage{english}
In Africa, the holy martyrs Castus 
 and Aemilius, who met their martyrdom by fire, St. Cyprian says that there 
 were overcome by the first trial, but that in the second God made them 
 victorious, so that those who had first weakened in the face of the fire 
 were made mightier than the flames.
\switchcolumn*
\selectlanguage{latin}
Cománæ, in Ponto, 
 sancti Basilísci Mártyris, qui, sub Maximiáno Imperatóre et Agríppa Prǽside, 
 férreas calceátus crépidas, ignítis clavis confíxas, múltaque ália passus, 
 demum, cápite obtruncátus et in flumen projéctus, martyrii glóriam 
 consecútus est.
\switchcolumn
\selectlanguage{english}
At Comana in Pontus, under Emperor 
 Maximian and the governor Agrippa, the holy martyr Basiliscus, who was 
 forced to wear iron shoes pierced with heated nails, and who endured many 
 other trials. He was finally beheaded and thrown into the river, which 
 gained for him the crown of martyrdom.
\switchcolumn*
\selectlanguage{latin}
In Córsica sanctæ 
 Júliæ Vírginis, quæ crucis supplício coronáta est.
\switchcolumn
\selectlanguage{english}
In Corsica, St. Julia, virgin, who 
 won her crown by being crucified.
\switchcolumn*
\selectlanguage{latin}
In Hispánia sanctæ 
 Quitériæ, Vírginis et Mártyris.
\switchcolumn
\selectlanguage{english}
In Spain, St. Quiteria, virgin and 
 martyr.
\switchcolumn*
\selectlanguage{latin}
Ravénnæ sancti 
 Marciáni, Epíscopi et Confessóris.
\switchcolumn
\selectlanguage{english}
At Ravenna, St. Marcian, bishop and 
 confessor.
\switchcolumn*
\selectlanguage{latin}
Pistórii, in Túscia, 
 beáti Atthónis Epíscopi, ex Ordine Vallis Umbrósæ.
\switchcolumn
\selectlanguage{english}
At Pistoia in Tuscany, the bishop, 
 blessed Attho, of the Order of Vallombrosa.
\switchcolumn*
\selectlanguage{latin}
In pago Antisiodorénsi 
 beáti Románi Abbátis, qui sancto Benedícto ministrávit in specu; inde, in 
 Gállias proféctus, ibi, ædificáto monastério relictísque multis sanctitátis 
 alúmnis, quiévit in Dómino.
\switchcolumn
\selectlanguage{english}
In the diocese of Auxerre, Abbot 
 St. Romanus, who ministered to St. Benedict in his cave. Going later 
 to France, he built a monastery there, and leaving many disciples and 
 imitators of his sanctity, went to rest in the Lord.
\switchcolumn*
\selectlanguage{latin}
Apud Aquínum sancti 
 Fulci Confessóris.
\switchcolumn
\selectlanguage{english}
At Aquino, St. Fulk, confessor.
\switchcolumn*
\selectlanguage{latin}
Antisiodóri sanctæ 
 Hélenæ Vírginis.
\switchcolumn
\selectlanguage{english}
At Auxerre, St. Helen, virgin.
\switchcolumn*
\selectlanguage{latin}
Cássiæ, in Umbria, 
 sanctæ Ritæ Víduæ, Moniális ex Ordine Eremitárum sancti Augustíni; quæ, post 
 sæculi núptias, ætérnum sponsum Christum únice diléxit.
\switchcolumn
\selectlanguage{english}
At Cascia in Umbria, St. Rita, a 
 widow and nun of the Order of the Hermits of St. Augustine, who, after being 
 disengaged from her earthly marriage, loved only her eternal spouse Christ.
\switchcolumn*
\selectlanguage{latin}
\end{paracol}


% ---- martyrology/mart05/mart0523.htm
\needspace{10\baselineskip}
\begin{paracol}{2}
\selectlanguage{latin}
\begin{center}{\color{gregoriocolor} Décimo Kaléndas Júnii. 
 Luna\dots\ }\end{center}
\switchcolumn
\selectlanguage{english}
\begin{center}{\color{gregoriocolor} The 
 23\textsuperscript{rd} Day of May. The\dots\ Day of the Moon.}\end{center}
\end{paracol}

\noindent\begin{tabularx}{\linewidth}{*{19}{>{\centering\arraybackslash}X}}
 \textcolor{gregoriocolor}{a} & \textcolor{gregoriocolor}{b} & \textcolor{gregoriocolor}{c} & \textcolor{gregoriocolor}{d} & \textcolor{gregoriocolor}{e} & \textcolor{gregoriocolor}{f} & \textcolor{gregoriocolor}{g} & \textcolor{gregoriocolor}{h} & \textcolor{gregoriocolor}{i} & \textcolor{gregoriocolor}{k} & \textcolor{gregoriocolor}{l} & \textcolor{gregoriocolor}{m} & \textcolor{gregoriocolor}{n} & \textcolor{gregoriocolor}{p} & \textcolor{gregoriocolor}{q} & \textcolor{gregoriocolor}{r} & \textcolor{gregoriocolor}{s} & \textcolor{gregoriocolor}{t} & \textcolor{gregoriocolor}{u} \\
 26 & 27 & 28 & 29 & 30 & 1 & 2 & 3 & 4 & 5 & 6 & 7 & 8 & 9 & 10 & 11 & 12 & 13 & 14 \\
\end{tabularx}
\vspace{0.5\baselineskip}
\noindent\begin{tabularx}{\linewidth}{*{12}{>{\centering\arraybackslash}X}}
 \textcolor{gregoriocolor}{A} & \textcolor{gregoriocolor}{B} & \textcolor{gregoriocolor}{C} & \textcolor{gregoriocolor}{D} & \textcolor{gregoriocolor}{E} & F & \textcolor{gregoriocolor}{F} & \textcolor{gregoriocolor}{G} & \textcolor{gregoriocolor}{H} & \textcolor{gregoriocolor}{M} & \textcolor{gregoriocolor}{N} & \textcolor{gregoriocolor}{P} \\
 15 & 16 & 17 & 18 & 19 & 20 & 20 & 21 & 22 & 23 & 24 & 25 \\
\end{tabularx}

\begin{paracol}{2}
\selectlanguage{latin}
\lettrine[lines=2]{A}{pud} Língonas, in 
 Gállia, pássio sancti Desidérii Epíscopi, qui, cum plebem suam ab exércitu 
 Wandalórum vexári cérneret, ad Regem eórum pro ea supplicatúrus accéssit; a 
 quo statim jugulári jussus, pro óvibus sibi créditis cervícem libénter 
 exténdit, et, percússus gládio, migrávit ad Christum. Passi sunt autem 
 cum ipso et álii plures de número gregis sui, qui apud eándem urbem cónditi 
 sunt.
\switchcolumn
\selectlanguage{english}
\lettrine[lines=2]{A}{t} Langres in France, the martyrdom 
 of the holy bishop Desiderius, who visited the king to offer entreaties in 
 behalf of his people who were mistreated by the Vandal army. He was 
 immediately condemned to beheading, and willingly presenting his head to 
 receive the blow of the sword, he died for the sheep committed to his charge 
 and departed for heaven. With him suffered many of his flock, who are 
 buried in the same city.
\switchcolumn*
\selectlanguage{latin}
In Hispánia sanctórum 
 Mártyrum Epitácii Epíscopi, et Basiléi.
\switchcolumn
\selectlanguage{english}
In Spain, the holy martyrs 
 Epitacius, a bishop, and Basileus.
\switchcolumn*
\selectlanguage{latin}
In território 
 Lugdunénsi sancti Desidérii, Epíscopi Viennénsis; qui, Theodoríci Regis 
 jussu lapídibus óbrutus, martyrio coronátur.
\switchcolumn
\selectlanguage{english}
In the territory of Lyons, St. 
 Desiderius, bishop of Vienne, who was crowned with martyrdom by being 
 stoned at the order of King Theodoric.
\switchcolumn*
\selectlanguage{latin}
In Africa sanctórum 
 Mártyrum Quinctiáni, Lúcii et Juliáni, qui, in persecutióne Wandálica passi, 
 ætérnas corónas meruérunt.
\switchcolumn
\selectlanguage{english}
In Africa, the holy martyrs 
 Quintian, Lucius, and Julian, who merited eternal crowns by their 
 sufferings, during the persecution of the Vandals.
\switchcolumn*
\selectlanguage{latin}
In Cappadócia 
 commemorátio sanctórum Mártyrum, qui, in persecutióne Maximiáni Galérii, 
 confráctis crúribus, necáti sunt; item eórum, qui eódem témpore, in 
 Mesopotámia, appénsi pédibus in sublíme, cápite verso deórsum, suffocáti 
 fumo et lento igne consúmpti, martyrium complevérunt.
\switchcolumn
\selectlanguage{english}
In Cappadocia, the commemoration of 
 the holy martyrs who died by having their legs crushed, in the persecution 
 of Maximian Galerius. Also in Mesopotamia, those martyrs who, at the 
 same time, were suspended in the air with their heads downward, suffocated 
 with smoke, and consumed by a slow fire, thus fulfilling their martyrdom.
\switchcolumn*
\selectlanguage{latin}
Synnadæ, in Phrygia, 
 sancti Michaélis Epíscopi.
\switchcolumn
\selectlanguage{english}
At Synnada in Phrygia, St. Michael, 
 bishop.
\switchcolumn*
\selectlanguage{latin}
Ipso die sancti 
 Mercuriális Epíscopi.
\switchcolumn
\selectlanguage{english}
The same day, St. Mercurialis, 
 bishop.
\switchcolumn*
\selectlanguage{latin}
Neápoli, in Campánia, 
 sancti Euphébii Epíscopi.
\switchcolumn
\selectlanguage{english}
At Naples in Campania, St. 
 Euphebius, bishop.
\switchcolumn*
\selectlanguage{latin}
Romæ sancti Joánnis 
 Baptístæ de Rossi, Presbyteri et Confessóris, patiéntia et caritáte in 
 evangelizándis paupéribus insígnis.
\switchcolumn
\selectlanguage{english}
At Rome, St. John Baptist de Rossi, 
 priest and confessor, a man illustrious for his patience and his zeal in 
 preaching the Gospel to the poor.
\switchcolumn*
\selectlanguage{latin}
Apud Núrsiam sanctórum 
 Eutychii et Floréntii Monachórum, quorum méminit beátus Gregórius Papa.
\switchcolumn
\selectlanguage{english}
At Norcia, Saints Eutychius and 
 Florentius, monks, mentioned by the blessed Pope Gregory.
\switchcolumn*
\selectlanguage{latin}
\end{paracol}


% ---- martyrology/mart05/mart0524.htm
\needspace{10\baselineskip}
\begin{paracol}{2}
\selectlanguage{latin}
\begin{center}{\color{gregoriocolor} Nono Kaléndas Júnii. 
 Luna\dots\ }\end{center}
\switchcolumn
\selectlanguage{english}
\begin{center}{\color{gregoriocolor} The 
 24\textsuperscript{th} Day of May. The\dots\ Day of the Moon.}\end{center}
\end{paracol}

\noindent\begin{tabularx}{\linewidth}{*{19}{>{\centering\arraybackslash}X}}
 \textcolor{gregoriocolor}{a} & \textcolor{gregoriocolor}{b} & \textcolor{gregoriocolor}{c} & \textcolor{gregoriocolor}{d} & \textcolor{gregoriocolor}{e} & \textcolor{gregoriocolor}{f} & \textcolor{gregoriocolor}{g} & \textcolor{gregoriocolor}{h} & \textcolor{gregoriocolor}{i} & \textcolor{gregoriocolor}{k} & \textcolor{gregoriocolor}{l} & \textcolor{gregoriocolor}{m} & \textcolor{gregoriocolor}{n} & \textcolor{gregoriocolor}{p} & \textcolor{gregoriocolor}{q} & \textcolor{gregoriocolor}{r} & \textcolor{gregoriocolor}{s} & \textcolor{gregoriocolor}{t} & \textcolor{gregoriocolor}{u} \\
 27 & 28 & 29 & 30 & 1 & 2 & 3 & 4 & 5 & 6 & 7 & 8 & 9 & 10 & 11 & 12 & 13 & 14 & 15 \\
\end{tabularx}
\vspace{0.5\baselineskip}
\noindent\begin{tabularx}{\linewidth}{*{12}{>{\centering\arraybackslash}X}}
 \textcolor{gregoriocolor}{A} & \textcolor{gregoriocolor}{B} & \textcolor{gregoriocolor}{C} & \textcolor{gregoriocolor}{D} & \textcolor{gregoriocolor}{E} & F & \textcolor{gregoriocolor}{F} & \textcolor{gregoriocolor}{G} & \textcolor{gregoriocolor}{H} & \textcolor{gregoriocolor}{M} & \textcolor{gregoriocolor}{N} & \textcolor{gregoriocolor}{P} \\
 16 & 17 & 18 & 19 & 20 & 21 & 21 & 22 & 23 & 24 & 25 & 26 \\
\end{tabularx}

\begin{paracol}{2}
\selectlanguage{latin}
\lettrine[lines=2]{A}{ntiochíæ} natális 
 sancti Mánahen, qui fuit Heródis Tetrárchæ collactáneus, atque, Doctor et 
 Prophéta exsístens sub grátia novi Testaménti, in eádem urbe quiévit.
\switchcolumn
\selectlanguage{english}
\lettrine[lines=2]{A}{t} Antioch, the birthday of St. 
 Manahen, foster-brother of Herod the Tetrach. He was a doctor and 
 prophet under the grace of the New Testament, and his remains now lie in the 
 city of Antioch.
\switchcolumn*
\selectlanguage{latin}
Item beátæ Joánnæ, 
 uxóris Chuzæ, procuratóris Heródis, quam Lucas Evangelísta commémorat.
\switchcolumn
\selectlanguage{english}
Also, blessed Joanna, wife of Chuza, 
 Herod's steward, mentioned by the evangelist St. Luke.
\switchcolumn*
\selectlanguage{latin}
In Portu Románo 
 natális sancti Vincéntii Mártyris.
\switchcolumn
\selectlanguage{english}
At Porto, the birthday of St. 
 Vincent, martyr.
\switchcolumn*
\selectlanguage{latin}
Nannéte, in Británnia minóre, beatórum Mártyrum Donatiáni et Rogatiáni fratrum, qui, sub 
 Diocletiáno Imperatóre, pro constántia fídei, in cárcerem missi, et in 
 equúleo suspénsi ac laniáti, deínde láncea militári confóssi, novíssime 
 cápite præcísi sunt.
\switchcolumn
\selectlanguage{english}
At Nantes in Brittany, in the time 
 of Emperor Diocletian, the blessed martyrs Donatian and Rogatian, brothers, 
 who, because of their constancy in the faith, were sent to prison, stretched 
 on the rack, and lacerated. Finally, they were pierced through with a 
 soldier's lance, and then beheaded.
\switchcolumn*
\selectlanguage{latin}
In Istria sanctórum 
 Mártyrum Zoélli, Servílii, Felícis, Silváni, et Dióclis.
\switchcolumn
\selectlanguage{english}
In Istria, the holy martyrs Zoellus, 
 Servilius, Felix, Silvanus, and Diocles.
\switchcolumn*
\selectlanguage{latin}
Eódem die sanctórum 
 Mártyrum Melétii, ducis exércitus, ac Sociórum ejus ducentórum et 
 quinquagínta duórum mílitum; qui divérso mortis génere martyrium 
 complevérunt.
\switchcolumn
\selectlanguage{english}
Also, the holy martyrs Meletius, 
 who was a military officer, and two hundred and fifty-two of his companions, 
 who achieved their martyrdom by various kinds of deaths.
\switchcolumn*
\selectlanguage{latin}
Item sanctárum 
 Mártyrum Susánnæ, Marciánæ et Palládiæ, eorúndem mílitum uxórum, quæ, una 
 cum párvulis suis, confráctæ sunt.
\switchcolumn
\selectlanguage{english}
Also, the holy martyrs Susanna, 
 Marciana, and Palladia, wives of the soldiers just mentioned, who were put 
 to death with their young children.
\switchcolumn*
\selectlanguage{latin}
Medioláni sancti 
 Robustiáni Mártyris.
\switchcolumn
\selectlanguage{english}
At Milan, St. Robustian, martyr.
\switchcolumn*
\selectlanguage{latin}
Bríxiæ sanctæ Afræ 
 Mártyris, quæ sub Hadriáno Imperatóre passa est.
\switchcolumn
\selectlanguage{english}
At Brescia, St. Afra, martyr, who 
 suffered under Emperor Hadrian.
\switchcolumn*
\selectlanguage{latin}
In monastério 
 Lirinénsi, in Gállia, sancti Vincéntii Presbyteri, doctrína et sanctitáte 
 conspícui.
\switchcolumn
\selectlanguage{english}
In the monastery of Lerins, St. 
 Vincent, a priest eminent for learning and sanctity.
\switchcolumn*
\selectlanguage{latin}
Marróchii, in Africa, 
 beáti Joánnis de Prado, Sacerdótis ex Ordine Minórum et Mártyris; qui, in 
 prædicatióne Evangélii, post víncula, cárceres, flagélla plurimáque ália 
 torménta, pro Christo fórtiter toleráta, per ignem martyrium consummávit.
\switchcolumn
\selectlanguage{english}
At Morocco in Africa, the passion 
 of blessed John of Prado, priest and martyr of the Order of Friars Minor. 
 While preaching the Gospel, he was bound, imprisoned, and scourged; and 
 after enduring with fortitude many other torments for Christ, fulfilled his 
 martyrdom by fire.
\switchcolumn*
\selectlanguage{latin}
Bonóniæ Translátio 
 sancti Domínici Confessóris, témpore Gregórii Papæ Noni.
\switchcolumn
\selectlanguage{english}
At Bologna, the translation of St. 
 Dominic, confessor, in the time of Pope Gregory IX.
\switchcolumn*
\selectlanguage{latin}
\end{paracol}


% ---- martyrology/mart05/mart0525.htm
\needspace{10\baselineskip}
\begin{paracol}{2}
\selectlanguage{latin}
\begin{center}{\color{gregoriocolor} Octávo Kaléndas Júnii. 
 Luna\dots\ }\end{center}
\switchcolumn
\selectlanguage{english}
\begin{center}{\color{gregoriocolor} The 
 25\textsuperscript{th} Day of May. The\dots\ Day of the Moon.}\end{center}
\end{paracol}

\noindent\begin{tabularx}{\linewidth}{*{19}{>{\centering\arraybackslash}X}}
 \textcolor{gregoriocolor}{a} & \textcolor{gregoriocolor}{b} & \textcolor{gregoriocolor}{c} & \textcolor{gregoriocolor}{d} & \textcolor{gregoriocolor}{e} & \textcolor{gregoriocolor}{f} & \textcolor{gregoriocolor}{g} & \textcolor{gregoriocolor}{h} & \textcolor{gregoriocolor}{i} & \textcolor{gregoriocolor}{k} & \textcolor{gregoriocolor}{l} & \textcolor{gregoriocolor}{m} & \textcolor{gregoriocolor}{n} & \textcolor{gregoriocolor}{p} & \textcolor{gregoriocolor}{q} & \textcolor{gregoriocolor}{r} & \textcolor{gregoriocolor}{s} & \textcolor{gregoriocolor}{t} & \textcolor{gregoriocolor}{u} \\
 28 & 29 & 30 & 1 & 2 & 3 & 4 & 5 & 6 & 7 & 8 & 9 & 10 & 11 & 12 & 13 & 14 & 15 & 16 \\
\end{tabularx}
\vspace{0.5\baselineskip}
\noindent\begin{tabularx}{\linewidth}{*{12}{>{\centering\arraybackslash}X}}
 \textcolor{gregoriocolor}{A} & \textcolor{gregoriocolor}{B} & \textcolor{gregoriocolor}{C} & \textcolor{gregoriocolor}{D} & \textcolor{gregoriocolor}{E} & F & \textcolor{gregoriocolor}{F} & \textcolor{gregoriocolor}{G} & \textcolor{gregoriocolor}{H} & \textcolor{gregoriocolor}{M} & \textcolor{gregoriocolor}{N} & \textcolor{gregoriocolor}{P} \\
 17 & 18 & 19 & 20 & 21 & 22 & 22 & 23 & 24 & 25 & 26 & 27 \\
\end{tabularx}

\begin{paracol}{2}
\selectlanguage{latin}
\lettrine[lines=2]{S}{alérni} deposítio 
 beáti Gregórii Séptimi, Papæ et Confessóris, ecclesiásticæ libertátis 
 propugnatóris ac defensóris acérrimi.
\switchcolumn
\selectlanguage{english}
\lettrine[lines=2]{A}{t} Salerno, the death of blessed 
 Pope Gregory VII, a most zealous protector and champion of Church liberty.
\switchcolumn*
\selectlanguage{latin}
Romæ, via Nomentána, 
 natális beáti Urbáni Primi, Papæ et Mártyris, cujus exhortatióne et doctrína 
 complúres (in quibus fuére Tibúrtius et Valeriánus) Christi suscepérunt 
 fidem, et pro ea martyrium subiérunt. Ipse quoque, in persecutióne 
 Alexándri Sevéri, multa passus pro Ecclésia Dei, tandem, cervícibus 
 abscíssis, martyrio coronátus est.
\switchcolumn
\selectlanguage{english}
At Rome, on the Via Nomentana, the 
 birthday of blessed Urban, pope and martyr, by whose exhortation and 
 teaching many persons, among whom were Tiburtius and Valerian, received the 
 faith of Christ and suffered martyrdom for it. He himself endured many 
 afflictions for the Church of God, and was crowned with martyrdom by being 
 beheaded in the persecution of Alexander Severus.
\switchcolumn*
\selectlanguage{latin}
Girvi, in Anglia, 
 tránsitus sancti Bedæ Venerábilis, Presbyteri, Confessóris et Ecclésiæ 
 Doctóris, sanctitáte et eruditióne celebérrimi. Ipsíus autem festum 
 recólitur sexto Kaléndas Júnii.
\switchcolumn
\selectlanguage{english}
At Jarrow in England, the death of 
 St. Venerable Bede, priest, confessor and doctor of the Church, well known 
 for his sanctity and scholarship. His feast, however, is celebrated on 
 the 27th day of May.
\switchcolumn*
\selectlanguage{latin}
Floréntiæ natális 
 sanctæ Maríæ-Magdalénæ de Pazzis, Vírginis, ex Ordine Carmelitárum, vita et 
 sanctitáte illústris. Ejus vero festívitas quarto Kaléndas Júnii 
 celebrátur.
\switchcolumn
\selectlanguage{english}
At Florence, the birthday of St. 
 Mary Magdalen de Pazzi, a virgin of the Order of the Carmelites, who is 
 famed for her holy life. Her feast is observed on the 29th of May.
\switchcolumn*
\selectlanguage{latin}
Doróstori, in Mysia 
 inferióre, item natális sanctórum Mártyrum Pasícratis, Valentiónis et 
 aliórum duórum, simul coronatórum.
\switchcolumn
\selectlanguage{english}
At Silistria in Bulgaria, the 
 birthday of the holy martyrs Pasicrates, Valentio, and two others crowned 
 with them.
\switchcolumn*
\selectlanguage{latin}
Medioláni sancti 
 Dionysii Epíscopi, qui, ab Ariáno Imperatóre Constántio in Cappadóciam pro 
 fide cathólica relegátus, ibídem, propióre Martyribus título, spíritum Deo 
 réddidit. Ejus sacrum corpus ab Aurélio Epíscopo transmíssum est 
 Mediolánum ad beátum Ambrósium Epíscopum; cui piæ actióni offícium quoque 
 sancti Basilíi Magni accessísse tráditur.
\switchcolumn
\selectlanguage{english}
At Milan, Bishop St. Denis, who for 
 the Catholic faith was exiled into Cappadocia by the Arian emperor 
 Constantius, where he yielded his soul to God in a manner almost like that 
 of the martyrs. His revered body was sent to blessed Bishop Ambrose at 
 Milan, by Bishop Aurelius, with the help, it is said, of St. Basil the 
 Great.
\switchcolumn*
\selectlanguage{latin}
Floréntiæ natális 
 sancti Zenóbii, ejúsdem civitátis Epíscopi, vitæ sanctitáte ac miraculórum 
 glória conspícui.
\switchcolumn
\selectlanguage{english}
At Florence, the birthday of St. 
 Zenobius, bishop of that city, renowned for the sanctity of his life and his 
 glorious miracles.
\switchcolumn*
\selectlanguage{latin}
In Británnia sancti 
 Aldélmi, Epíscopi Schireburgénsis.
\switchcolumn
\selectlanguage{english}
In England, St. Aldhelm, bishop of 
 Sherborne.
\switchcolumn*
\selectlanguage{latin}
In território 
 Tricassíno sancti Leónis Confessóris.
\switchcolumn
\selectlanguage{english}
In the territory of Troyes, St. 
 Leo, confessor.
\switchcolumn*
\selectlanguage{latin}
Lutétiæ Parisiórum 
 sanctæ Magdalénæ-Sophíæ Barat, Fundatrícis Institúti Sorórum a Sacro Corde 
 Jesu; quæ pro Christiána puellárum informatióne valde adlaborávit, et a Pio 
 Papa Undécimo in sanctárum Vírginum catálogum fuit reláta.
\switchcolumn
\selectlanguage{english}
At Paris, St. Madeleine-Sophie 
 Barat, foundress of the Congregation of the Sisters of the Sacred Heart, who 
 devoted her labours for the Christian education of girls. She was 
 added to the list of holy virgins by Pope Pius XI.
\switchcolumn*
\selectlanguage{latin}
Vérulis, in Hérnicis, 
 Translátio sanctæ Maríæ Jacóbi; cujus sacrum corpus 
 plúrimis miráculis 
 illustrátur.
\switchcolumn
\selectlanguage{english}
At Veroli in Campania, the 
 translation of St. Mary, the mother of James, whose revered body is noted 
 for many miracles.
\switchcolumn*
\selectlanguage{latin}
Assísii, in Umbria, 
 item Translátio sancti Francísci Confessóris, témpore Gregórii Papæ Noni.
\switchcolumn
\selectlanguage{english}
At Assisi in Umbria, the 
 translation of St. Francis, confessor, in the time of Pope Gregory IX.
\switchcolumn*
\selectlanguage{latin}
\end{paracol}


% ---- martyrology/mart05/mart0526.htm
\needspace{10\baselineskip}
\begin{paracol}{2}
\selectlanguage{latin}
\begin{center}{\color{gregoriocolor} Séptimo Kaléndas Júnii. 
 Luna\dots\ }\end{center}
\switchcolumn
\selectlanguage{english}
\begin{center}{\color{gregoriocolor} The 
 26\textsuperscript{th} Day of May. The\dots\ Day of the Moon.}\end{center}
\end{paracol}

\noindent\begin{tabularx}{\linewidth}{*{19}{>{\centering\arraybackslash}X}}
 \textcolor{gregoriocolor}{a} & \textcolor{gregoriocolor}{b} & \textcolor{gregoriocolor}{c} & \textcolor{gregoriocolor}{d} & \textcolor{gregoriocolor}{e} & \textcolor{gregoriocolor}{f} & \textcolor{gregoriocolor}{g} & \textcolor{gregoriocolor}{h} & \textcolor{gregoriocolor}{i} & \textcolor{gregoriocolor}{k} & \textcolor{gregoriocolor}{l} & \textcolor{gregoriocolor}{m} & \textcolor{gregoriocolor}{n} & \textcolor{gregoriocolor}{p} & \textcolor{gregoriocolor}{q} & \textcolor{gregoriocolor}{r} & \textcolor{gregoriocolor}{s} & \textcolor{gregoriocolor}{t} & \textcolor{gregoriocolor}{u} \\
 29 & 30 & 1 & 2 & 3 & 4 & 5 & 6 & 7 & 8 & 9 & 10 & 11 & 12 & 13 & 14 & 15 & 16 & 17 \\
\end{tabularx}
\vspace{0.5\baselineskip}
\noindent\begin{tabularx}{\linewidth}{*{12}{>{\centering\arraybackslash}X}}
 \textcolor{gregoriocolor}{A} & \textcolor{gregoriocolor}{B} & \textcolor{gregoriocolor}{C} & \textcolor{gregoriocolor}{D} & \textcolor{gregoriocolor}{E} & F & \textcolor{gregoriocolor}{F} & \textcolor{gregoriocolor}{G} & \textcolor{gregoriocolor}{H} & \textcolor{gregoriocolor}{M} & \textcolor{gregoriocolor}{N} & \textcolor{gregoriocolor}{P} \\
 18 & 19 & 20 & 21 & 22 & 23 & 23 & 24 & 25 & 26 & 27 & 28 \\
\end{tabularx}

\begin{paracol}{2}
\selectlanguage{latin}
\lettrine[lines=2]{R}{omæ} sancti Philíppi 
 Nérii, Presbyteri et Confessóris, qui Congregatiónis Oratórii Fundátor fuit, 
 ac virginitáte, prophetíæ dono, et miráculis éxstitit insígnis.
\switchcolumn
\selectlanguage{english}
\lettrine[lines=2]{A}{t} Rome, St. Philip Neri, priest and 
 confessor, founder of the Congregation of the Oratory, celebrated for his 
 virginal purity, the gift of prophecy, and miracles.
\switchcolumn*
\selectlanguage{latin}
Item Romæ sancti 
 Eleuthérii, Papæ et Mártyris, qui multos nóbiles Romanórum ad fidem Christi 
 perdúxit, et sanctos Damiánum et Fugátium in Británniam misit, qui Lúcium 
 Regem, cum uxóre ipsíus ac toto fere pópulo, baptizárunt.
\switchcolumn
\selectlanguage{english}
Also at Rome, St. Eleutherius, pope 
 and martyr, who converted to the Christian faith many noble Romans. He 
 sent Saints Damian and Fugatius to England, and they baptized King Lucius, 
 his wife, and almost all his people.
\switchcolumn*
\selectlanguage{latin}
Cantuáriæ, in Anglia, 
 natális sancti Augustíni, Epíscopi et Confessóris; qui, una cum áliis, a 
 beáto Gregório Papa missus, genti Anglórum sacrum Christi Evangélium 
 prædicávit., ibíque, virtútibus et miráculis gloriósus, obdormívit in 
 Dómino. Ejus tamen festívitas quinto Kaléndas Júnii recólitur.
\switchcolumn
\selectlanguage{english}
At Canterbury in England, St. 
 Augustine, bishop, who was sent there with others by blessed Pope Gregory, 
 and who preached the Gospel of Christ to the English nation. 
 Celebrated for virtues and miracles, he went peacefully to his rest in the 
 Lord. The 28th of May is observed as his feast.
\switchcolumn*
\selectlanguage{latin}
Athénis item natális 
 beáti Quadráti, Apostolórum discípuli, qui, in persecutióne Hadriáni, fide 
 et indústria sua cóngregans Ecclésiam grándi terróre dispérsam, librum pro 
 Christiánæ religiónis defensióne, valde 
 útilem et Apostólica doctrína dignum, 
 eídem Imperatóri porréxit.
\switchcolumn
\selectlanguage{english}
At Athens, during the persecution 
 of Hadrian, the birthday of blessed Quadratus, a disciple of the apostles, 
 who collected by his zealous work the faithful who had dispersed through 
 terror, and presented to the emperor a book which was an excellent apology 
 of the Christian religion, worthy of an apostle.
\switchcolumn*
\selectlanguage{latin}
Romæ sanctórum 
 Mártyrum Simítrii Presbyteri, et aliórum vigínti duórum; qui sub Antoníno 
 Pio passi sunt.
\switchcolumn
\selectlanguage{english}
At Rome, the holy martyrs Simitrius, 
 priest, and twenty-two others who suffered under Antoninus Pius.
\switchcolumn*
\selectlanguage{latin}
Viénnæ, in Gállia, 
 sancti Zacharíæ, Epíscopi et Mártyris, qui sub Trajáno passus est.
\switchcolumn
\selectlanguage{english}
At Vienne, St. Zacharias, bishop and 
 martyr, who suffered under Trajan.
\switchcolumn*
\selectlanguage{latin}
In Africa sancti 
 Quadráti Mártyris, in cujus solemnitáte sanctus Augustínus sermónem hábuit.
\switchcolumn
\selectlanguage{english}
In Africa, St. Quadratus, martyr, 
 on whose feast day St. Augustine preached a sermon.
\switchcolumn*
\selectlanguage{latin}
Tudérti, in Umbria, 
 natális sanctórum Mártyrum Felicíssimi, Heráclii et Paulíni.
\switchcolumn
\selectlanguage{english}
At Todi in Umbria, the birthday of 
 the holy martyrs Felicissimus, Heraclius, and Paulinus.
\switchcolumn*
\selectlanguage{latin}
In território Antisiodorénsi pássio sancti Prisci Mártyris, qui, cum ingénti multitúdine 
 fidélium Christi, cápite cæsus est.
\switchcolumn
\selectlanguage{english}
In the territory of Auxerre, the 
 passion of St. Priscus, martyr, along with a great multitude of other 
 Christians.
\switchcolumn*
\selectlanguage{latin}
In civitáte Quiténsi, 
 Æquatoriánæ Ditiónis, sanctæ Maríæ Annæ a Jesu de Parédes Vírginis, e tértio 
 Ordine sancti Francísci, austeritáte et in próximum caritáte præcláræ, quam 
 Pius Papa Duodécimus sanctárum Vírginum catálogo adnumerávit.
\switchcolumn
\selectlanguage{english}
In the city of Quito in Ecuador, 
 St. María Ana de Jesù de Paredes, a third order Franciscan, well known for 
 her austerity and charity towards her neighbour. Pope Pius XII 
 numbered her in the book of Virgins.
\switchcolumn*
\selectlanguage{latin}
\end{paracol}


% ---- martyrology/mart05/mart0527.htm
\needspace{10\baselineskip}
\begin{paracol}{2}
\selectlanguage{latin}
\begin{center}{\color{gregoriocolor} Sexto Kaléndas Júnii. 
 Luna\dots\ }\end{center}
\switchcolumn
\selectlanguage{english}
\begin{center}{\color{gregoriocolor} The 
 27\textsuperscript{th} Day of May. The\dots\ Day of the Moon.}\end{center}
\end{paracol}

\noindent\begin{tabularx}{\linewidth}{*{19}{>{\centering\arraybackslash}X}}
 \textcolor{gregoriocolor}{a} & \textcolor{gregoriocolor}{b} & \textcolor{gregoriocolor}{c} & \textcolor{gregoriocolor}{d} & \textcolor{gregoriocolor}{e} & \textcolor{gregoriocolor}{f} & \textcolor{gregoriocolor}{g} & \textcolor{gregoriocolor}{h} & \textcolor{gregoriocolor}{i} & \textcolor{gregoriocolor}{k} & \textcolor{gregoriocolor}{l} & \textcolor{gregoriocolor}{m} & \textcolor{gregoriocolor}{n} & \textcolor{gregoriocolor}{p} & \textcolor{gregoriocolor}{q} & \textcolor{gregoriocolor}{r} & \textcolor{gregoriocolor}{s} & \textcolor{gregoriocolor}{t} & \textcolor{gregoriocolor}{u} \\
 30 & 1 & 2 & 3 & 4 & 5 & 6 & 7 & 8 & 9 & 10 & 11 & 12 & 13 & 14 & 15 & 16 & 17 & 18 \\
\end{tabularx}
\vspace{0.5\baselineskip}
\noindent\begin{tabularx}{\linewidth}{*{12}{>{\centering\arraybackslash}X}}
 \textcolor{gregoriocolor}{A} & \textcolor{gregoriocolor}{B} & \textcolor{gregoriocolor}{C} & \textcolor{gregoriocolor}{D} & \textcolor{gregoriocolor}{E} & F & \textcolor{gregoriocolor}{F} & \textcolor{gregoriocolor}{G} & \textcolor{gregoriocolor}{H} & \textcolor{gregoriocolor}{M} & \textcolor{gregoriocolor}{N} & \textcolor{gregoriocolor}{P} \\
 19 & 20 & 21 & 22 & 23 & 24 & 24 & 25 & 26 & 27 & 28 & 29 \\
\end{tabularx}

\begin{paracol}{2}
\selectlanguage{latin}
\lettrine[lines=2]{S}{ancti} Bedæ Venerábilis, 
 Presbyteri, Confessóris et Ecclésiæ Doctóris; qui migrávit in cælum octávo 
 Kaléndas Júnii.
\switchcolumn
\selectlanguage{english}
\lettrine[lines=2]{S}{t.} Venerable Bede, priest, 
 confessor, and doctor of the Church, who went to heaven on the 25th of May.
\switchcolumn*
\selectlanguage{latin}
Sancti Joánnis Primi, 
 Papæ et Mártyris; cujus dies natális quintodécimo Kaléndas Júnii refértur, 
 sed festívitas hac die, ob Translatiónem córporis ejus, potíssimum 
 celebrátur.
\switchcolumn
\selectlanguage{english}
St. John I, pope and martyr. 
 His birthday is observed on the 18th of May, but his feast is celebrated 
 today because of the translation of his revered body.
\switchcolumn*
\selectlanguage{latin}
Doróstori, in Mysia 
 inferióre, pássio beáti Júlii, qui, témpore Alexándri Imperatóris, cum esset 
 veteránus et eméritæ milítiæ, comprehénsus est ab officiálibus, et Máximo 
 Prǽsidi oblátus; quo præsénte, cum exsecrarétur idóla, et Christi nomen 
 constantíssime confiterétur, capitáli senténtia punítus est.
\switchcolumn
\selectlanguage{english}
At Silistria in Bulgaria, during 
 the reign of Emperor Alexander, the martyrdom of blessed Julius, a veteran 
 soldier in retirement, who was arrested by the officials and presented to 
 the governor Maximus. Having denounced the idols in his presence, and 
 confessed the name of Christ with utmost constancy, he was condemned to 
 capital punishment.
\switchcolumn*
\selectlanguage{latin}
In pago Atrebaténsi 
 sancti Ranúlfi Mártyris.
\switchcolumn
\selectlanguage{english}
In the district of Arras, St. 
 Ralph, martyr.
\switchcolumn*
\selectlanguage{latin}
Apud Soram sanctæ 
 Restitútæ, Vírginis et Mártyris; quæ, sub Aureliáno Imperatóre et Agáthio 
 Procónsule, fídei certámine suscépto, dǽmonum ímpetus, paréntum blandítias 
 et tortórum sævítiam superávit, ac demum, cum áliis Christiánis, cápite 
 truncáta, martyrio decoráta est.
\switchcolumn
\selectlanguage{english}
At Sora, in the time of Emperor 
 Aurelian and the proconsul Agathius, St. Restituta, virgin and martyr, who 
 overcame in a trial for the faith the violence of the demons, the affections 
 of her family, and the cruelty of the executioners. Being finally 
 beheaded with other Christians, she obtained the honour of martyrdom.
\switchcolumn*
\selectlanguage{latin}
Aráusicæ, in Gálliis, 
 sancti Eutrópii Epíscopi, virtútibus atque miráculis illústris.
\switchcolumn
\selectlanguage{english}
At Orange in France, St. Eutropius, 
 a bishop illustrious for virtues and miracles.
\switchcolumn*
\selectlanguage{latin}
Herbípoli, in Germánia, 
 sancti Brunónis, Epíscopi et Confessóris.
\switchcolumn
\selectlanguage{english}
At Wurzburg in Germany, St. Bruno, 
 bishop and confessor.
\switchcolumn*
\selectlanguage{latin}
\end{paracol}


% ---- martyrology/mart05/mart0528.htm
\needspace{10\baselineskip}
\begin{paracol}{2}
\selectlanguage{latin}
\begin{center}{\color{gregoriocolor} Quinto Kaléndas Júnii. 
 Luna\dots\ }\end{center}
\switchcolumn
\selectlanguage{english}
\begin{center}{\color{gregoriocolor} The 
 28\textsuperscript{th} Day of May. The\dots\ Day of the Moon.}\end{center}
\end{paracol}

\noindent\begin{tabularx}{\linewidth}{*{19}{>{\centering\arraybackslash}X}}
 \textcolor{gregoriocolor}{a} & \textcolor{gregoriocolor}{b} & \textcolor{gregoriocolor}{c} & \textcolor{gregoriocolor}{d} & \textcolor{gregoriocolor}{e} & \textcolor{gregoriocolor}{f} & \textcolor{gregoriocolor}{g} & \textcolor{gregoriocolor}{h} & \textcolor{gregoriocolor}{i} & \textcolor{gregoriocolor}{k} & \textcolor{gregoriocolor}{l} & \textcolor{gregoriocolor}{m} & \textcolor{gregoriocolor}{n} & \textcolor{gregoriocolor}{p} & \textcolor{gregoriocolor}{q} & \textcolor{gregoriocolor}{r} & \textcolor{gregoriocolor}{s} & \textcolor{gregoriocolor}{t} & \textcolor{gregoriocolor}{u} \\
 1 & 2 & 3 & 4 & 5 & 6 & 7 & 8 & 9 & 10 & 11 & 12 & 13 & 14 & 15 & 16 & 17 & 18 & 19 \\
\end{tabularx}
\vspace{0.5\baselineskip}
\noindent\begin{tabularx}{\linewidth}{*{12}{>{\centering\arraybackslash}X}}
 \textcolor{gregoriocolor}{A} & \textcolor{gregoriocolor}{B} & \textcolor{gregoriocolor}{C} & \textcolor{gregoriocolor}{D} & \textcolor{gregoriocolor}{E} & F & \textcolor{gregoriocolor}{F} & \textcolor{gregoriocolor}{G} & \textcolor{gregoriocolor}{H} & \textcolor{gregoriocolor}{M} & \textcolor{gregoriocolor}{N} & \textcolor{gregoriocolor}{P} \\
 20 & 21 & 22 & 23 & 24 & 25 & 25 & 26 & 27 & 28 & 29 & 30 \\
\end{tabularx}

\begin{paracol}{2}
\selectlanguage{latin}
\lettrine[lines=2]{S}{ancti} Augustíni, 
 Epíscopi Cantuariénsis et Confessóris, cujus dies natális ágitur séptimo 
 Kaléndas Júnii.
\switchcolumn
\selectlanguage{english}
\lettrine[lines=2]{S}{t.} Augustine, bishop of Canterbury 
 and confessor, whose birthday is mentioned on the 26th of May.
\switchcolumn*
\selectlanguage{latin}
In Sardínia sanctórum 
 Mártyrum Æmílii, Felícis, Priámi et Luciáni, qui, pro Christo certántes, ab 
 eo glorióse coronáti sunt.
\switchcolumn
\selectlanguage{english}
In Sardinia, the holy martyrs 
 Aemilius, Priamus, and Lucian, who gained their crowns after being in the 
 combat for Christ.
\switchcolumn*
\selectlanguage{latin}
Carnóti, in Gállia, 
 sancti Caráuni Mártyris, qui sub Domitiáno Imperatóre, cápite amputátus, 
 martyrium sumpsit.
\switchcolumn
\selectlanguage{english}
At Chartres in France, under 
 Emperor Domitian, St. Caraunus, martyr, who was beheaded, and thus acquired 
 the glory of martyrdom.
\switchcolumn*
\selectlanguage{latin}
Item pássio sanctórum 
 Crescéntis, Dioscóridis, Pauli et Helládii.
\switchcolumn
\selectlanguage{english}
Also the martyrdom of the Saints 
 Crescens, Dioscorides, Paul, and Helladius.
\switchcolumn*
\selectlanguage{latin}
Thécuæ, in Palæstína, 
 sanctórum Monachórum Mártyrum, qui, témpore Theodósii junióris, a Saracénis 
 occísi sunt, quorum sacras relíquias collegérunt áccolæ, et summa 
 veneratióne illas habuérunt.
\switchcolumn
\selectlanguage{english}
At Thecua in Palestine, the saintly 
 monks who became martyrs by being killed by the Saracens, in the time of 
 Theodosius the Younger. Their venerable remains were gathered by the 
 inhabitants and preserved with greatest reverence.
\switchcolumn*
\selectlanguage{latin}
Corínthi sanctæ 
 Helicónidis Mártyris, témpore Gordiáni Imperatóris. Hæc primum, sub 
 Perénnio Prǽside, multis torméntis afflícta, deínde, sub ejus successóre 
 Justíno, íterum cruciáta, sed ab Angelis liberáta est; demum, disséctis 
 mammis, ferísque objécta atque igne probáta, cápitis obtruncatióne martyrium 
 complévit.
\switchcolumn
\selectlanguage{english}
At Corinth, St. Helconides, martyr, 
 who was first subjected to torments in the reign of Emperor Gordian, under 
 the governor Perennius, and then again tortured under his successor Justin, 
 but was delivered by an angel. Her breasts were cut away, she was 
 exposed to wild beasts and to fire, and finally her martyrdom was fulfilled 
 by beheading.
\switchcolumn*
\selectlanguage{latin}
Lutétiæ Parisiórum 
 sancti Germáni, Epíscopi et Confessóris; qui quantæ sanctitátis quantíque 
 fúerit mériti, quibus étiam miráculis clarúerit, litterárum moniméntis 
 Fortunátus Epíscopus consignávit.
\switchcolumn
\selectlanguage{english}
At Paris, St. Germanus, bishop and 
 confessor, whose fame for holiness, merit, and miracles has been handed down 
 to us by the writings of Bishop Fortunatus.
\switchcolumn*
\selectlanguage{latin}
Medioláni sancti 
 Senatóris Epíscopi, virtútibus et eruditióne claríssimi.
\switchcolumn
\selectlanguage{english}
At Milan, St. Senator, bishop, who 
 was very well known for his virtues and his learning.
\switchcolumn*
\selectlanguage{latin}
Urgéllæ, in Hispánia 
 Tarraconénsi, sancti Justi Epíscopi.
\switchcolumn
\selectlanguage{english}
At Urgel in Spain, Bishop St. 
 Justus.
\switchcolumn*
\selectlanguage{latin}
Floréntiæ sancti Pódii, 
 Epíscopi et Confessóris.
\switchcolumn
\selectlanguage{english}
At Florence, St. Podius, bishop and 
 confessor.
\switchcolumn*
\selectlanguage{latin}
Apud Nováriam sancti 
 Bernárdi a Menthóne, Confessóris, qui in monte Jovis, super Alpes, in 
 Valésia, celebérrimum cœnóbium et hospítium exstrúxit, atque a Pio Papa 
 Undécimo, non modo Alpium íncolis vel viatóribus, sed iis étiam qui eárum 
 juga ascendéndo se exércent, cæléstis Patrónus est attribútus.
\switchcolumn
\selectlanguage{english}
At Novara, St. Bernard of Menton, 
 confessor. On Mount Jou in the Alps of Valais in Switzerland, he 
 founded the famous monastery and hospice. Pope Pius XI appointed him 
 the heavenly patron not only of those who live in or travel across the Alps, 
 but of all mountain climbers.
\switchcolumn*
\selectlanguage{latin}
\end{paracol}


% ---- martyrology/mart05/mart0529.htm
\needspace{10\baselineskip}
\begin{paracol}{2}
\selectlanguage{latin}
\begin{center}{\color{gregoriocolor} Quarto Kaléndas Júnii. 
 Luna\dots\ }\end{center}
\switchcolumn
\selectlanguage{english}
\begin{center}{\color{gregoriocolor} The 
 29\textsuperscript{th} Day of May. The\dots\ Day of the Moon.}\end{center}
\end{paracol}

\noindent\begin{tabularx}{\linewidth}{*{19}{>{\centering\arraybackslash}X}}
 \textcolor{gregoriocolor}{a} & \textcolor{gregoriocolor}{b} & \textcolor{gregoriocolor}{c} & \textcolor{gregoriocolor}{d} & \textcolor{gregoriocolor}{e} & \textcolor{gregoriocolor}{f} & \textcolor{gregoriocolor}{g} & \textcolor{gregoriocolor}{h} & \textcolor{gregoriocolor}{i} & \textcolor{gregoriocolor}{k} & \textcolor{gregoriocolor}{l} & \textcolor{gregoriocolor}{m} & \textcolor{gregoriocolor}{n} & \textcolor{gregoriocolor}{p} & \textcolor{gregoriocolor}{q} & \textcolor{gregoriocolor}{r} & \textcolor{gregoriocolor}{s} & \textcolor{gregoriocolor}{t} & \textcolor{gregoriocolor}{u} \\
 2 & 3 & 4 & 5 & 6 & 7 & 8 & 9 & 10 & 11 & 12 & 13 & 14 & 15 & 16 & 17 & 18 & 19 & 20 \\
\end{tabularx}
\vspace{0.5\baselineskip}
\noindent\begin{tabularx}{\linewidth}{*{12}{>{\centering\arraybackslash}X}}
 \textcolor{gregoriocolor}{A} & \textcolor{gregoriocolor}{B} & \textcolor{gregoriocolor}{C} & \textcolor{gregoriocolor}{D} & \textcolor{gregoriocolor}{E} & F & \textcolor{gregoriocolor}{F} & \textcolor{gregoriocolor}{G} & \textcolor{gregoriocolor}{H} & \textcolor{gregoriocolor}{M} & \textcolor{gregoriocolor}{N} & \textcolor{gregoriocolor}{P} \\
 21 & 22 & 23 & 24 & 25 & 26 & 26 & 27 & 28 & 29 & 30 & 1 \\
\end{tabularx}

\begin{paracol}{2}
\selectlanguage{latin}
\lettrine[lines=2]{S}{anctæ} Maríæ-Magdalénæ 
 de Pazzis, ex Ordine Carmelitárum, Vírginis, cujus dies natális octávo 
 Kaléndas Júnii recensétur.
\switchcolumn
\selectlanguage{english}
\lettrine[lines=2]{S}{t.} Mary Magdalene of Pazzi of the 
 Order of Carmelites, and virgin. Her birthday was mentioned on the 
 25th of May.
\switchcolumn*
\selectlanguage{latin}
Romæ, via Aurélia, natális sancti Restitúti Mártyris.
\switchcolumn
\selectlanguage{english}
At Rome, on the Via Aurelia, the 
 birthday of St. Restitutus, martyr.
\switchcolumn*
\selectlanguage{latin}
Apud Icónium, in 
 Lycaónia, pássio sanctórum Conónis, et fílii annórum duódecim, qui, sub 
 Aureliáno Imperatóre, cratículæ, prunis suppósitis et óleo superinfúso 
 candéntis, suspensiónis in equúleo atque ignis pœnam constánter passi, ad 
 extrémum, málleo lígneo mánibus eórum contrítis, spíritum emisérunt.
\switchcolumn
\selectlanguage{english}
At Iconium in Lycaonia, in the time 
 of Emperor Aurelian, the martyrdom of the Saints Conon and his son, a child 
 twelve years of age, who were laid on a grate over burning coals sprinkled 
 with oil, placed on the rack, and exposed to the fire. Finally their 
 hands were crushed with a mallet, and they breathed their last.
\switchcolumn*
\selectlanguage{latin}
In agro Tridentíno 
 natális sanctórum Mártyrum Sisínii, Martyrii et Alexándri; qui, témpore 
 Honórii Imperatóris, in Anáuniæ pártibus (ut scribit in vita sancti Ambrósii 
 Paulínus), persequéntibus Gentílibus, martyrii corónam adépti sunt.
\switchcolumn
\selectlanguage{english}
In the district of Trent, in the 
 time of Emperor Honorius, the birthday of the holy martyrs Sisinius, 
 Martyrius, and Alexander, who were persecuted by the heathens of Anaunia, 
 and obtained the crown of martyrdom, all of which is told by Paulinus in the
 Life of Ambrose.
\switchcolumn*
\selectlanguage{latin}
Cameríni pássio 
 sanctórum mille quingentórum et vigínti quinque Mártyrum.
\switchcolumn
\selectlanguage{english}
At Camerino, the passion of fifteen 
 hundred and twenty-five holy martyrs.
\switchcolumn*
\selectlanguage{latin}
Cæsaréæ Philíppi 
 sanctárum Mártyrum Theodósiæ, quæ sancti Procópii Mártyris 
 éxstitit mater, 
 et aliárum duódecim nobílium matronárum; quæ, in persecutióne Diocletiáni, 
 cápitis obtruncatióne consummátæ sunt.
\switchcolumn
\selectlanguage{english}
At Caesarea Philippi, the holy 
 martyrs Theodosia, mother of the martyr St. Procopius, and twelve other 
 noble women, whose lives were ended by their being beheaded in the 
 persecution of Diocletian.
\switchcolumn*
\selectlanguage{latin}
Tréviris beáti 
 Maximíni, Epíscopi et Confessóris; a quo sanctus Athanásius Epíscopus, ob 
 persecutiónem Arianórum éxsulans honorífice suscéptus fuit.
\switchcolumn
\selectlanguage{english}
At Treves, blessed Maximinus, 
 bishop and confessor, who received with honour the patriarch St. Athanasius 
 when he was banished by the Arian persecutors.
\switchcolumn*
\selectlanguage{latin}
Verónæ sancti Máximi 
 Epíscopi.
\switchcolumn
\selectlanguage{english}
At Verona, St. Maximus, bishop.
\switchcolumn*
\selectlanguage{latin}
Arcáni, in Látio, 
 sancti Eleuthérii Confessóris.
\switchcolumn
\selectlanguage{english}
At Arcano in Lazio, St. Eleutherius, 
 confessor.
\switchcolumn*
\selectlanguage{latin}
\end{paracol}


% ---- martyrology/mart05/mart0530.htm
\needspace{10\baselineskip}
\begin{paracol}{2}
\selectlanguage{latin}
\begin{center}{\color{gregoriocolor} Tértio Kaléndas Júnii. 
 Luna\dots\ }\end{center}
\switchcolumn
\selectlanguage{english}
\begin{center}{\color{gregoriocolor} The 
 30\textsuperscript{th} Day of May. The\dots\ Day of the Moon.}\end{center}
\end{paracol}

\noindent\begin{tabularx}{\linewidth}{*{19}{>{\centering\arraybackslash}X}}
 \textcolor{gregoriocolor}{a} & \textcolor{gregoriocolor}{b} & \textcolor{gregoriocolor}{c} & \textcolor{gregoriocolor}{d} & \textcolor{gregoriocolor}{e} & \textcolor{gregoriocolor}{f} & \textcolor{gregoriocolor}{g} & \textcolor{gregoriocolor}{h} & \textcolor{gregoriocolor}{i} & \textcolor{gregoriocolor}{k} & \textcolor{gregoriocolor}{l} & \textcolor{gregoriocolor}{m} & \textcolor{gregoriocolor}{n} & \textcolor{gregoriocolor}{p} & \textcolor{gregoriocolor}{q} & \textcolor{gregoriocolor}{r} & \textcolor{gregoriocolor}{s} & \textcolor{gregoriocolor}{t} & \textcolor{gregoriocolor}{u} \\
 3 & 4 & 5 & 6 & 7 & 8 & 9 & 10 & 11 & 12 & 13 & 14 & 15 & 16 & 17 & 18 & 19 & 20 & 21 \\
\end{tabularx}
\vspace{0.5\baselineskip}
\noindent\begin{tabularx}{\linewidth}{*{12}{>{\centering\arraybackslash}X}}
 \textcolor{gregoriocolor}{A} & \textcolor{gregoriocolor}{B} & \textcolor{gregoriocolor}{C} & \textcolor{gregoriocolor}{D} & \textcolor{gregoriocolor}{E} & F & \textcolor{gregoriocolor}{F} & \textcolor{gregoriocolor}{G} & \textcolor{gregoriocolor}{H} & \textcolor{gregoriocolor}{M} & \textcolor{gregoriocolor}{N} & \textcolor{gregoriocolor}{P} \\
 22 & 23 & 24 & 25 & 26 & 27 & 27 & 28 & 29 & 30 & 1 & 2 \\
\end{tabularx}

\begin{paracol}{2}
\selectlanguage{latin}
\lettrine[lines=2]{S}{ancti} Felícis Primi, 
 Papæ et Mártyris, cujus dies natális tértio Kaléndas Januárii recensétur.
\switchcolumn
\selectlanguage{english}
\lettrine[lines=2]{P}{ope} St. Felix I, martyr, whose 
 birthday is commemorated on the 30th of December.
\switchcolumn*
\selectlanguage{latin}
Túrribus, in Sardínia, sanctórum Mártyrum Gabíni et Críspuli.
\switchcolumn
\selectlanguage{english}
At Torres in Sardinia, the holy 
 martyrs Gabinus and Crispulus.
\switchcolumn*
\selectlanguage{latin}
Antiochíæ sanctórum 
 Syci et Palatíni, qui, pro Christi nómine, multa torménta passi sunt.
\switchcolumn
\selectlanguage{english}
At Antioch, Saints Sycus and 
 Palatinus, who endured many torments for the name of Christ.
\switchcolumn*
\selectlanguage{latin}
Ravénnæ sancti 
 Exsuperántii, Epíscopi et Confessóris.
\switchcolumn
\selectlanguage{english}
At Ravenna, St. Exuperantius, 
 bishop and confessor.
\switchcolumn*
\selectlanguage{latin}
Papíæ sancti Anastásii 
 Epíscopi.
\switchcolumn
\selectlanguage{english}
At Pavia, St. Anastasius, bishop.
\switchcolumn*
\selectlanguage{latin}
Cæsaréæ, in Cappadócia, 
 sanctórum Basilíi et Emméliæ uxóris, qui fuérunt paréntes beatórum Basilíi 
 Magni et Gregórii Nysséni ac Petri Sebasténsis Episcopórum, atque Macrínæ 
 Vírginis. Hi vero sancti cónjuges, témpore Galérii Maximiáni, extórres 
 facti, Pónticas solitúdines incolúere; et post persecutiónem, fíliis suárum 
 relíctis virtútum herédibus, in pace quievérunt.
\switchcolumn
\selectlanguage{english}
At Caesarea in Cappadocia, the 
 Saints Basil and his wife Emmelia, parents of St. Basil the Great, St. 
 Gregory of Nyssa, St. Peter of Sebastopol, bishops, and St. Macrina, virgin. 
 They lived in exile in the deserts of Pontus during the reign of Galerius 
 Maximian, and after the persecution they died in peace, leaving their 
 children as heirs of their virtues.
\switchcolumn*
\selectlanguage{latin}
Híspali, in Hispánia, 
 sancti Ferdinándi Tértii, Castéllæ et Legiónis Regis, ob virtútum 
 præstántiam cognoménto Sancti; qui, fídei propagándæ zelo clarus, tandem, 
 devíctis Mauris, ad cæléste regnum, terréno relícto, felíciter evolávit.
\switchcolumn
\selectlanguage{english}
At Seville in Spain, St. Ferdinand 
 III, king of Castile and Leon. He was surnamed the Saint on account of 
 his eminent virtues; he was celebrated for his zeal in spreading the faith. 
 After conquering the Moors he left his kingdom on earth to pass happily to 
 that of heaven.
\switchcolumn*
\selectlanguage{latin}
Rotómagi sanctæ Joánnæ 
 Arcénsis Vírginis, Puéllæ Aurelianénsis appellátæ, quæ cum fórtiter pro 
 pátria dimicásset, tandem, in hóstium potestátem trádita, iníquo judício 
 condemnáta est et igne combústa; atque a Benedícto Décimo quinto, Pontífice 
 Máximo, Sanctórum fastis adscrípta.
\switchcolumn
\selectlanguage{english}
At Rouen, St. Joan of Arc, virgin, 
 called the Maid of Orleans. After fighting heroically for her 
 fatherland, she was at the end delivered into the hands of the enemies, 
 condemned by an unjust judge, and burned at the stake. The supreme 
 Pontiff Benedict XV placed her name on the canon of the saints.
\switchcolumn*
\selectlanguage{latin}
\end{paracol}


% ---- martyrology/mart05/mart0531.htm
\needspace{10\baselineskip}
\begin{paracol}{2}
\selectlanguage{latin}
\begin{center}{\color{gregoriocolor} Prídie Kaléndas Júnii. 
 Luna\dots\ }\end{center}
\switchcolumn
\selectlanguage{english}
\begin{center}{\color{gregoriocolor} The 
 31\textsuperscript{st} Day of May. The\dots\ Day of the Moon.}\end{center}
\end{paracol}

\noindent\begin{tabularx}{\linewidth}{*{19}{>{\centering\arraybackslash}X}}
 \textcolor{gregoriocolor}{a} & \textcolor{gregoriocolor}{b} & \textcolor{gregoriocolor}{c} & \textcolor{gregoriocolor}{d} & \textcolor{gregoriocolor}{e} & \textcolor{gregoriocolor}{f} & \textcolor{gregoriocolor}{g} & \textcolor{gregoriocolor}{h} & \textcolor{gregoriocolor}{i} & \textcolor{gregoriocolor}{k} & \textcolor{gregoriocolor}{l} & \textcolor{gregoriocolor}{m} & \textcolor{gregoriocolor}{n} & \textcolor{gregoriocolor}{p} & \textcolor{gregoriocolor}{q} & \textcolor{gregoriocolor}{r} & \textcolor{gregoriocolor}{s} & \textcolor{gregoriocolor}{t} & \textcolor{gregoriocolor}{u} \\
 4 & 5 & 6 & 7 & 8 & 9 & 10 & 11 & 12 & 13 & 14 & 15 & 16 & 17 & 18 & 19 & 20 & 21 & 22 \\
\end{tabularx}
\vspace{0.5\baselineskip}
\noindent\begin{tabularx}{\linewidth}{*{12}{>{\centering\arraybackslash}X}}
 \textcolor{gregoriocolor}{A} & \textcolor{gregoriocolor}{B} & \textcolor{gregoriocolor}{C} & \textcolor{gregoriocolor}{D} & \textcolor{gregoriocolor}{E} & F & \textcolor{gregoriocolor}{F} & \textcolor{gregoriocolor}{G} & \textcolor{gregoriocolor}{H} & \textcolor{gregoriocolor}{M} & \textcolor{gregoriocolor}{N} & \textcolor{gregoriocolor}{P} \\
 23 & 24 & 25 & 26 & 27 & 28 & 28 & 29 & 30 & 1 & 2 & 3 \\
\end{tabularx}

\begin{paracol}{2}
\selectlanguage{latin}
\lettrine[lines=1]{F}{estum} beátæ Maríæ Vírginis Regínæ.
\switchcolumn
\selectlanguage{english}
\lettrine[lines=1]{F}{east} of the Queenship of the Blessed Virgin Mary.
\switchcolumn*
\selectlanguage{latin}
Romæ sanctæ Petroníllæ 
 Vírginis, fíliæ beáti Petri Apóstoli, quæ, conjúgium nóbilis viri Flacci 
 spernens, et, accéptis triduánis ad deliberándum indúciis, ínterim jejúniis 
 et oratiónibus vacans, tértia die, mox ut Christi Sacraméntum accépit, 
 emísit spíritum.
\switchcolumn
\selectlanguage{english}
At Rome, St. Petronilla, virgin, 
 disciple of the blessed apostle Peter. She refused to marry Flaccus, a 
 nobleman, and was granted three days for deliberation. She spent these 
 days in fasting and in prayer, and on the third day, after having received 
 the Sacrament of the Body of Christ, she yielded up her soul.
\switchcolumn*
\selectlanguage{latin}
Aquiléjæ sanctórum 
 Mártyrum fratrum Cántii, Cantiáni et Cantianíllæ, qui, cum essent ex 
 illústri Aniciórum progénie, sub Diocletiáno et Maximiáno Imperatóribus, ob 
 Christiánæ fídei constántiam, una cum pædagógo suo Proto, cápite plexi sunt.
\switchcolumn
\selectlanguage{english}
At Aquileia, the holy martyrs 
 Cantius, Cantian, and Cantianilla, members of one family, which belonged to 
 the illustrious line of the Anicii. For their attachment to the 
 Christian faith, they were condemned to capital punishment with their tutor, 
 Protus, in the time of Emperors Diocletian and Maximian.
\switchcolumn*
\selectlanguage{latin}
Túrribus, in Sardínia, sancti Crescentiáni Mártyris.
\switchcolumn
\selectlanguage{english}
At Torres in Sardinia, St. 
 Crescentian, martyr.
\switchcolumn*
\selectlanguage{latin}
Apud Comános, in Ponto, 
 sancti Hérmiæ mílitis, qui, sub Antoníno Imperatóre, de innúmeris et 
 sævíssimis torméntis divína ope liberátus, carníficem convértit ad Christum, 
 et ejúsdem martyrii corónæ partícipem fecit; quam tamen ipse primus, gládio 
 obtruncátus, accépit.
\switchcolumn
\selectlanguage{english}
At Comana in Pontus during the 
 reign of Emperor Antoninus, St. Hermias, a soldier. Being miraculously 
 delivered from many horrible torments, he converted his executioner to 
 Christ, and made him partaker of the crown which he was first to receive by 
 being beheaded.
\switchcolumn*
\selectlanguage{latin}
Verónæ sancti Lupicíni 
 Epíscopi.
\switchcolumn
\selectlanguage{english}
At Verona, St. Lupicinus, bishop.
\switchcolumn*
\selectlanguage{latin}
Romæ sancti Paschásii, 
 Diáconi et Confessóris, cujus méminit beátus Gregórius Papa.
\switchcolumn
\selectlanguage{english}
At Rome, St. Paschasius, deacon and 
 confessor, who is mentioned by blessed Pope Gregory.
\switchcolumn*
\selectlanguage{latin}
\end{paracol}
